\input{preamble}

% OK, start here.
%
\begin{document}

\title{Exercices}


\maketitle

\phantomsection
\label{section-phantom}

\tableofcontents


\section{Algèbre}
\label{section-algebra}

\noindent
Cette première section ne contient que quelques questions diverses.

\begin{exercise}
\label{exercise-isomorphism-localizations}
Soient $A$ un anneau et ${\mathfrak m}$ un idéal maximal. Dans $A[X]$,
posons $\tilde {\mathfrak m}_1 = ({\mathfrak m}, X)$ et
$\tilde {\mathfrak m}_2 = ({\mathfrak m}, X-1)$. Montrer
que
$$
A[X]_{\tilde {\mathfrak m}_1} \cong A[X]_{\tilde {\mathfrak m}_2}.
$$
\end{exercise}

\begin{exercise}
\label{exercise-coherent}
Trouver un exemple d'anneau non noethérien $R$ tel que tout idéal
de type fini de $R$ soit de présentation finie comme $R$-module.
(On dit qu'un anneau est {\it cohérent} si cette dernière propriété est vérifiée.)
\end{exercise}

\begin{exercise}
\label{exercise-flat-ideals-pid}
Supposons que $(A, {\mathfrak m}, k)$ soit un anneau local noethérien. Pour tout
$A$-module de type fini $M$, notons $r(M)$ le nombre minimal
de générateurs de $M$ comme $A$-module. D'après le lemme de Nakayama, ce nombre vaut
$\dim_k M/{\mathfrak m} M = \dim_k M \otimes_A k$.
\begin{enumerate}
\item Montrer que $r(M \otimes_A N) = r(M)r(N)$.
\item Soit $I\subset A $ un idéal tel que $r(I) > 1$. Montrer que
$r(I^2) < r(I)^2$.
\item En déduire que, si tout idéal de $A$ est un module plat, alors
$A$ est un anneau principal (ou un corps).
\end{enumerate}
\end{exercise}

\begin{exercise}
\label{exercise-not-isomorphic}
Soit $k$ un corps. Montrer que les paires suivantes de
$k$-algèbres ne sont pas isomorphes :
\begin{enumerate}
\item $k[x_1, \ldots, x_n]$ et $k[x_1, \ldots, x_{n + 1}]$ pour tout
$n\geq 1$.
\item $k[a, b, c, d, e, f]/(ab + cd + ef)$ et $k[x_1, \ldots, x_n]$
pour $n = 5$.
\item $k[a, b, c, d, e, f]/(ab + cd + ef)$ et $k[x_1, \ldots, x_n]$
pour $n = 6$.
\end{enumerate}
\end{exercise}

\begin{remark}
\label{remark-simple-geometric}
L'idée de cet exercice est bien entendu de trouver, dans chaque cas,
un argument simple plutôt que d'appliquer un « gros » théorème.
Il est néanmoins utile de se laisser guider par des principes généraux.
\end{remark}

\begin{exercise}
\label{exercise-silly}
Algèbre. (Exercice un peu futile, qui devrait être facile.)
\begin{enumerate}
\item Donner un exemple d'un anneau $A$ et d'une
suite exacte courte non scindée de $A$-modules
$$
0 \to M_1 \to M_2 \to M_3 \to 0.
$$
\item Donner un exemple d'une suite non scindée de $A$-modules
comme ci-dessus et d'un homomorphisme fidèlement plat $A \to B$ tel que
$$
0 \to M_1\otimes_A B \to M_2\otimes_A B \to M_3\otimes_A B \to 0.
$$
soit scindée comme suite de $B$-modules.
\end{enumerate}
\end{exercise}

\begin{exercise}
\label{exercise-field-kummer}
Supposons que le corps $k$ contienne une racine primitive $n$-ième
de l'unité $\zeta$. Cela signifie que $\zeta^n = 1$, mais que $\zeta^m\not = 1$ pour
$0 < m < n$.
\begin{enumerate}
\item Montrer que la caractéristique de $k$ est première à $n$.
\item Supposons que $a \in k$ soit un élément de $k$ qui ne soit une
puissance $d$-ième dans $k$ pour aucun diviseur $d$ de $n$ tel que $n \geq d > 1$. Montrer que
$k[x]/(x^n-a)$ est un corps. (Indication : considérer un corps de décomposition de
$x^n-a$ et utiliser la théorie de Galois.)
\end{enumerate}
\end{exercise}

\begin{exercise}
\label{exercise-valuation}
Soit $\nu : k[x]\setminus \{0\}  \to {\mathbf Z}$ une application
ayant les propriétés suivantes : $\nu(fg) = \nu(f) + \nu(g)$ lorsque
$f$, $g$ sont non nuls, et $\nu(f + g) \geq min(\nu(f), \nu(g))$ lorsque
$f$, $g$ et $f + g$ sont non nuls, et $\nu(c) = 0$ pour tout $c\in k^*$.
\begin{enumerate}
\item Montrer que, si $f$, $g$ et $f + g$ sont non nuls et si
$\nu(f) \not = \nu(g)$, alors $\nu(f + g) = min(\nu(f), \nu(g))$.
\item Montrer que, si $f = \sum a_i x^i$, $f\not = 0$, alors
$\nu(f) \geq min(\{i\nu(x)\}_{a_i\not = 0})$. Quand a-t-on égalité ?
\item Montrer que, si $\nu$ prend une valeur négative, alors
$\nu(f) = -n \deg(f)$ pour un certain $n\in {\mathbf N}$.
\item Supposons que $\nu(x) \geq 0$. Montrer que
$\{f \mid f = 0, \ \text{ou}\ \nu(f) > 0\}$ est un idéal premier de $k[x]$.
\item Décrire toutes les applications $\nu$ possibles.
\end{enumerate}
\end{exercise}


\noindent
Soit $A$ un anneau. Un {\it idempotent} est un élément $e \in A$
tel que $e^2 = e$. Les éléments $1$ et $0$ sont toujours idempotents.
Un {\it idempotent non trivial} est un idempotent distinct de zéro et de l'unité.
Deux idempotents $e, e' \in A$ sont dits {\it orthogonaux}
si $ee' = 0$.

\begin{exercise}
\label{exercise-product}
Soit $A$ un anneau. Montrer que $A$ est un produit de deux anneaux non nuls si
et seulement si $A$ possède un idempotent non trivial.
\end{exercise}

\begin{exercise}
\label{exercise-lift-idempotents}
Soient $A$ un anneau et $I \subset A$ un idéal localement nilpotent.
Montrer que l'application $A \to A/I$ induit une bijection sur les idempotents.
(Indication : il peut être plus facile de commencer par le cas où $I$ est nilpotent.
Utiliser ensuite la « réduction noethérienne absolue » pour se ramener à ce cas.)
\end{exercise}







% Fin de la tranche source louée.
\section{Limites inductives}
\label{section-colimits}


\begin{definition}
\label{definition-directed-poset}
Un {\it ensemble préordonné filtrant} est un ensemble non vide $I$ muni d'un préordre $\leq$
tel que, pour toute paire $i, j \in I$, il existe un $k \in I$ tel que
$i \leq k$ et $j \leq k$. Un {\it système inductif d'anneaux} indexé par $I$ est la donnée d'un
anneau $A_i$ pour chaque $i \in I$ et d'un homomorphisme d'anneaux $\varphi_{ij} : A_i \to A_j$
dès que $i \leq j$, de sorte que le composé $A_i \to A_j \to A_k$ soit égal à
$A_i \to A_k$ dès que $i \leq j \leq k$.
\end{definition}

\noindent
On définit de même les systèmes inductifs de groupes, de modules sur un anneau fixé,
d'espaces vectoriels sur un corps, etc.

\begin{exercise}
\label{exercise-directed-colimit}
Soit $I$ un ensemble préordonné filtrant et soit
$(A_i, \varphi_{ij})$ un système inductif d'anneaux indexé par $I$.
Montrer qu'il existe un anneau $A$ et des homomorphismes $\varphi_i : A_i \to A$
tels que $\varphi_j \circ \varphi_{ij} = \varphi_i$ pour tous $i \leq j$
et possédant la propriété universelle suivante : étant donnés un anneau $B$
et des homomorphismes $\psi_i : A_i \to B$ tels que
$\psi_j \circ \varphi_{ij} = \psi_i$ pour tous $i \leq j$, il
existe un unique homomorphisme d'anneaux $\psi : A \to B$ tel que
$\psi_i = \psi \circ \varphi_i$.
\end{exercise}

\begin{definition}
\label{definition-colimit}
L'anneau $A$ construit dans l'exercice \ref{exercise-directed-colimit}
est appelé la {\it limite inductive} du système. Notation : $\colim A_i$.
\end{definition}

\begin{exercise}
\label{exercise-prime-in-colimit}
Soit $(I, \geq)$ un ensemble préordonné filtrant et soit
$(A_i, \varphi_{ij})$ un système inductif d'anneaux indexé par $I$, de limite inductive
$A$. Montrer qu'il existe une bijection
$$
\Spec(A) = \{(\mathfrak p_i)_{i \in I} \mid
\mathfrak p_i \subset A_i \text{ et }
\mathfrak p_i = \varphi_{ij}^{-1}(\mathfrak p_j)\ \forall i \leq j\}
\subset \prod\nolimits_{i \in I} \Spec(A_i)
$$
L'ensemble figurant au membre de droite de l'égalité est la limite projective des ensembles
$\Spec(A_i)$. Notation : $\lim \Spec(A_i)$.
\end{exercise}

\begin{exercise}
\label{exercise-colimit-surjective}
Soit $(I, \geq)$ un ensemble préordonné filtrant et soit
$(A_i, \varphi_{ij})$ un système inductif d'anneaux indexé par $I$, de limite inductive
$A$. Supposons que l'application $\Spec(A_j) \to \Spec(A_i)$ soit
surjective pour tous $i \leq j$. Montrer que
$\Spec(A) \to \Spec(A_i)$ est surjective pour tout $i$.
(Indication : on peut essayer d'utiliser le théorème de Tychonoff, mais il existe aussi une
démonstration algébrique directe essentiellement triviale fondée sur
Algèbre, lemme \ref{algebra-lemma-in-image}.)
\end{exercise}

\begin{exercise}
\label{exercise-integral-colimit-finite}
Soit une extension entière d'anneaux $A \subset B$. Montrer que
$\Spec(B) \to \Spec(A)$ est surjective.
Utiliser les exercices ci-dessus, le fait que l'assertion est vraie pour une extension finie
d'anneaux (ce qui a été démontré dans les cours), et montrer que
$B = \colim B_i$ est une limite inductive filtrante d'extensions finies $A \subset B_i$.
\end{exercise}

\begin{exercise}
\label{exercise-colimit-tensor}
Soit $(I, \geq)$ un ensemble préordonné filtrant.
Soit $A$ un anneau et soit $(N_i, \varphi_{i, i'})$ un système inductif de
$A$-modules indexé par $I$. Soit en outre $M$ un $A$-module. Montrer
que
$$
\colim_{i\in I} M \otimes_A N_i\cong
M \otimes_A \Big( \colim_{i\in I} N_i\Big).
$$
\end{exercise}

\begin{definition}
\label{definition-finite-presentation}
Un module $M$ sur $R$ est dit de {\it présentation finie} sur
$R$ s'il est isomorphe au conoyau d'un homomorphisme de modules libres de type fini
$ R^{\oplus n} \to R^{\oplus m}$.
\end{definition}

\begin{exercise}
\label{exercise-colimit-modules}
Montrer que tout module sur un anneau quelconque est
\begin{enumerate}
\item la limite inductive de ses sous-modules de type fini, et
\item, d'une certaine manière, une limite inductive de modules de présentation finie.
\end{enumerate}
\end{exercise}





\section{Catégories additives et abéliennes}
\label{section-additive}

\begin{exercise}
\label{exercise-filtered-vector-spaces}
Soit $k$ un corps. Soit $\mathcal{C}$ la catégorie des espaces vectoriels filtrés
sur $k$ ; voir
Homologie, définition \ref{homology-definition-filtered}
pour la définition d'un objet filtré d'une catégorie quelconque.
\begin{enumerate}
\item Montrer que c'est une catégorie additive (expliquer soigneusement ce qu'est la
somme directe de deux objets).
\item Soit $f : (V, F) \to (W, F)$ un morphisme de $\mathcal{C}$.
Montrer que $f$ possède un noyau et un conoyau (expliquer précisément
ce que sont le noyau et le conoyau de $f$).
\item Donner un exemple d'un morphisme de $\mathcal{C}$ tel que
le morphisme canonique $\Coim(f) \to \Im(f)$ ne soit pas un isomorphisme.
\end{enumerate}
\end{exercise}

\begin{exercise}
\label{exercise-torsion-free}
Soit $R$ un anneau intègre noethérien. Soit $\mathcal{C}$ la catégorie des
$R$-modules sans torsion de type fini.
\begin{enumerate}
\item Montrer que c'est une catégorie additive.
\item Soit $f : N \to M$ un morphisme de $\mathcal{C}$.
Montrer que $f$ possède un noyau et un conoyau (veiller à définir précisément
ce que sont le noyau et le conoyau de $f$).
\item Donner un exemple d'un anneau intègre noethérien $R$ et d'un morphisme de
$\mathcal{C}$ tel que le morphisme canonique $\Coim(f) \to \Im(f)$
ne soit pas un isomorphisme.
\end{enumerate}
\end{exercise}

\begin{exercise}
\label{exercise-other}
Donner un exemple d'une catégorie additive qui possède des noyaux
et des conoyaux, mais qui soit autre que celles des
exercices \ref{exercise-filtered-vector-spaces} et
\ref{exercise-torsion-free}.
\end{exercise}





\section{Produit tensoriel}
\label{section-tensor-product}

\noindent
Les produits tensoriels sont introduits dans Algèbre, section
\ref{algebra-section-tensor-product}.
Soit $R$ un anneau. Soit $\text{Mod}_R$ la catégorie des $R$-modules.
Nous dirons qu'un foncteur $F : \text{Mod}_R \to \text{Mod}_R$
\begin{enumerate}
\item est additif si
$F : \Hom_R(M, N) \to \Hom_R(F(M), F(N))$
est un homomorphisme de groupes abéliens
pour tous $R$-modules $M, N$ ; voir
Homologie, définition \ref{homology-definition-preadditive} ;
\item est $R$-linéaire si $F : \Hom_R(M, N) \to \Hom_R(F(M), F(N))$ est $R$-linéaire
pour tous $R$-modules $M, N$ ;
\item est exact à droite si, pour toute suite exacte courte
$0 \to M_1 \to M_2 \to M_3 \to 0$, la suite
$F(M_1) \to F(M_2) \to F(M_3) \to 0$ est exacte ;
\item est exact à gauche si, pour toute suite exacte courte
$0 \to M_1 \to M_2 \to M_3 \to 0$, la suite
$0 \to F(M_1) \to F(M_2) \to F(M_3)$ est exacte ;
\item commute aux sommes directes si, étant donnés un ensemble
$I$ et des $R$-modules $M_i$, les homomorphismes
$F(M_i) \to F(\bigoplus M_i)$ induisent un isomorphisme
$\bigoplus F(M_i) = F(\bigoplus M_i)$.
\end{enumerate}

\begin{exercise}
\label{exercise-characterize-tensor-functor}
Soit $R$ un anneau. On conserve les notations ci-dessus.
\begin{enumerate}
\item Donner un exemple d'un anneau $R$ et d'un foncteur additif
$F : \text{Mod}_R \to \text{Mod}_R$ qui ne soit pas $R$-linéaire.
\item Soit $N$ un $R$-module. Montrer que le foncteur
$F(M) = M \otimes_R N$ est $R$-linéaire,
exact à droite et commute aux sommes directes.
\item Réciproquement, montrer que tout foncteur $F : \text{Mod}_R \to \text{Mod}_R$
qui est $R$-linéaire,
exact à droite et commute aux sommes directes est de la
forme $F(M) = M \otimes_R N$ pour un certain $R$-module $N$.
\item Montrer que, si dans (3) on supprime l'hypothèse
selon laquelle $F$ commute aux sommes directes, la conclusion n'est
plus vraie.
\end{enumerate}
\end{exercise}





\section{Homomorphismes plats d'anneaux}
\label{section-flat}

\begin{exercise}
\label{exercise-localization-flat}
Soit $S$ une partie multiplicative de l'anneau $A$.
\begin{enumerate}
\item Pour un $A$-module $M$, montrer que $S^{-1}M = S^{-1}A \otimes_A M$.
\item Montrer que $S^{-1}A$ est plat sur $A$.
\end{enumerate}
\end{exercise}

\begin{exercise}
\label{exercise-examples-not-flat}
Trouver une injection $M_1 \to M_2$ de $A$-modules telle que
$M_1\otimes N \to M_2 \otimes N$ ne soit pas injective dans les
cas suivants :
\begin{enumerate}
\item $A = k[x, y]$ et $N = (x, y) \subset A$. (Ici et ci-dessous, $k$ est un corps.)
\item $A = k[x, y]$ et $N = A/(x, y)$.
\end{enumerate}
\end{exercise}

\begin{exercise}
\label{exercise-flat-not-projective}
Donner un exemple d'un anneau $A$ et d'un $A$-module de type fini $M$
qui soit plat mais non projectif comme $A$-module.
\end{exercise}

\begin{remark}
\label{remark-flat-not-projective}
Si $M$ est de présentation finie et plat sur $A$,
alors $M$ est projectif sur $A$. L'exemple devra donc faire
intervenir un anneau $A$ qui ne soit pas noethérien. Je connais un exemple
où $A$ est l'anneau des fonctions de classe ${\mathcal C}^\infty$ sur ${\mathbf R}$.
\end{remark}

\begin{exercise}
\label{exercise-flat-not-free-dvr}
Trouver un module plat mais non libre sur ${\mathbf Z}_{(2)}$.
\end{exercise}

\begin{exercise}
\label{exercise-flat-deformations}
Déformations plates.
\begin{enumerate}
\item Supposons que $k$ soit un corps et que $k[\epsilon]$ soit l'algèbre des
nombres duaux $k[\epsilon] = k[x]/(x^2)$, avec $\epsilon = \bar x$. Montrer que, pour
toute $k$-algèbre $A$, il existe une $k[\epsilon]$-algèbre plate $B$ telle que
$A$ soit isomorphe à $B/\epsilon B$.
\item Supposons que $k = {\mathbf F}_p = {\mathbf Z}/p{\mathbf Z}$ et
$$
A = k[x_1, x_2, x_3, x_4, x_5, x_6]/
(x_1^p, x_2^p, x_3^p, x_4^p, x_5^p, x_6^p).
$$
Montrer qu'il existe une ${\mathbf Z}/p^2{\mathbf Z}$-algèbre plate $B$ telle
que $B/pB$ soit isomorphe à $A$. (Ainsi, $p$ joue ici le rôle de $\epsilon$.)
\item Posons maintenant $p = 2$ et considérons la même question pour
$k = {\mathbf F}_2 = {\mathbf Z}/2{\mathbf Z}$ et
$$
A = k[x_1, x_2, x_3, x_4, x_5, x_6]/
(x_1^2, x_2^2, x_3^2, x_4^2, x_5^2, x_6^2, x_1x_2 + x_3x_4 + x_5x_6).
$$
Cependant, montrer que dans ce cas il n'existe {\it pas} de
${\mathbf Z}/4{\mathbf Z}$-algèbre plate $B$ telle que $B/2B$ soit isomorphe à
$A$. (Trouver l'astuce ! Le même exemple fonctionne en toute caractéristique
$p > 0$, mais le calcul est plus difficile.)
\end{enumerate}
\end{exercise}

\begin{exercise}
\label{exercise-flat-given-residue-field-extension}
Soit $(A, {\mathfrak m}, k)$ un anneau local et soit $k'/k$
une extension finie de corps. Montrer qu'il existe un homomorphisme local et plat
d'anneaux locaux $A \to B$ tel que ${\mathfrak m}_B = {\mathfrak m} B$ et que
$B/{\mathfrak m} B$ soit
isomorphe à $k'$ comme $k$-algèbre. (Indication : commencer par le cas où
l'extension $k \subset k'$ est engendrée par un seul élément.)
\end{exercise}

\begin{remark}
\label{remark-flat-given-residue-field-extension-general}
Le même résultat vaut pour toute extension de corps $K/k$.
\end{remark}





% Fin de la tranche source assignée.




\section{Le spectre d'un anneau}
\label{section-spectrum-ring}

\begin{exercise}
\label{exercise-spec-Z}
Déterminer $\Spec(\mathbf{Z})$ en tant qu'ensemble et décrire sa topologie.
\end{exercise}

\begin{exercise}
\label{exercise-basis-opens-standard}
Soit $A$ un anneau quelconque. Pour $f\in A$, on définit
$D(f):= \{\mathfrak p \subset A \mid f \not \in \mathfrak p\}$.
Montrer que les ouverts $D(f)$ forment une base de la topologie de
$\Spec(A)$.
\end{exercise}

\begin{exercise}
\label{exercise-radical-ideals-closed}
Montrer que l'application $I\mapsto V(I)$
définit une bijection naturelle
$$
\{I\subset A\text{ tel que }I = \sqrt{I}\}
\longrightarrow
\{T\subset \Spec(A)\text{ fermée}\}
$$
\end{exercise}

\begin{definition}
\label{definition-quasi-compact}
Un espace topologique $X$ est dit {\it quasi-compact}
si, pour tout recouvrement ouvert $X = \bigcup_{i\in I} U_i$, il existe une
partie finie $\{i_1, \ldots, i_n\}\subset I$ telle que $X = U_{i_1}\cup\ldots
U_{i_n}$.
\end{definition}

\begin{exercise}
\label{exercise-spec-quasi-compact}
Montrer que $\Spec(A)$ est quasi-compact pour tout anneau $A$.
\end{exercise}

\begin{definition}
\label{definition-Hausdorff}
On dit qu'un espace topologique $X$ vérifie l'axiome de séparation $T_0$
si, pour tout couple de points $x, y\in X$ tels que $x\not = y$, il existe un
ouvert de $X$ qui contient l'un mais pas l'autre.
On dit que $X$ est {\it séparé} si, pour tout couple $x, y\in X$ tel que $x\not = y$,
il existe des ouverts disjoints $U, V$ tels que $x\in U$
et $y\in V$.
\end{definition}

\begin{exercise}
\label{exercise-not-hausdorff}
Montrer que $\Spec(A)$ n'est {\bf pas} séparé en général.
Montrer que $\Spec(A)$ est $T_0$. Donner un exemple d'espace topologique
$X$ qui n'est pas $T_0$.
\end{exercise}

\begin{remark}
\label{remark-not-hausdorff}
En général, on réserve le mot compact aux espaces quasi-compacts et
séparés.
\end{remark}

\begin{definition}
\label{definition-irreducible}
Un espace topologique $X$ est dit {\it irréductible} si $X$ est non vide
et si $X = Z_1\cup Z_2$, avec $Z_1, Z_2\subset X$ fermés, alors soit
$Z_1 = X$, soit $Z_2 = X$. Une partie $T\subset X$ d'un espace topologique
est dite {\it irréductible} si elle est un espace topologique
irréductible pour la topologie induite par $X$.
Cette définition implique que $T$ est irréductible si et seulement
si l'adhérence $\bar T$ de $T$ dans $X$ est irréductible.
\end{definition}

\begin{exercise}
\label{exercise-irreducible-spec}
Montrer que $\Spec(A)$ est irréductible si et seulement si
$Nil(A)$ est un idéal premier et que, dans ce cas, c'est l'unique
idéal premier minimal de $A$.
\end{exercise}

\begin{exercise}
\label{exercise-irreducible-prime}
Montrer qu'une partie fermée $T\subset \Spec(A)$
est irréductible si et seulement si elle est de la forme $T = V({\mathfrak p})$ pour
un idéal premier ${\mathfrak p}\subset A$.
\end{exercise}

\begin{definition}
\label{definition-generic-point}
Un point $x$ d'un espace topologique irréductible $X$ est appelé
un {\it point générique} de $X$ si $X$ est égal à l'adhérence de
la partie $\{x\}$.
\end{definition}

\begin{exercise}
\label{exercise-irreducible-T0-at-most-one-generic}
Montrer que, dans un espace $T_0$ $X$, toute partie fermée
irréductible admet au plus un point générique.
\end{exercise}

\begin{exercise}
\label{exercise-spec-sober}
Montrer que, dans $\Spec(A)$, toute
partie fermée irréductible {\it admet bien} un point générique.
Montrer en fait que l'application
${\mathfrak p} \mapsto \overline{\{{\mathfrak p}\}}$ est
une bijection de $\Spec(A)$ sur l'ensemble des parties fermées
irréductibles de $X$.
\end{exercise}

\begin{exercise}
\label{exercise-irreducible-subset-not-generic}
Donner un exemple montrant qu'une partie irréductible de
$\Spec(\mathbf{Z})$ ne possède pas nécessairement de point générique.
\end{exercise}

\begin{definition}
\label{definition-Noetherian-space}
Un espace topologique $X$ est dit {\it noethérien} si toute
suite décroissante $Z_1\supset Z_2 \supset Z_3\supset \ldots$
de parties fermées de $X$ est stationnaire.
(Il est dit {\it artinien} si toute suite croissante de parties
fermées est stationnaire.)
\end{definition}

\begin{exercise}
\label{exercise-Noetherian-spec}
Montrer que, si l'anneau $A$ est noethérien, alors
l'espace topologique $\Spec(A)$ est noethérien. Donner un
exemple montrant que la réciproque est fausse. (Même question dans le cas
artinien, si vous le souhaitez.)
\end{exercise}

\begin{definition}
\label{definition-irreducible-component}
Une partie irréductible maximale $T\subset X$ est appelée une
{\it composante irréductible} de l'espace $X$. Une telle composante
irréductible de $X$ est automatiquement une partie fermée de $X$.
\end{definition}

\begin{exercise}
\label{exercise-irreducible-in-irreducible}
Montrer que toute partie irréductible
de $X$ est contenue dans une composante irréductible de $X$.
\end{exercise}

\begin{exercise}
\label{exercise-Noetherian-finite-nr-irreducible}
Montrer qu'un espace topologique noethérien $X$
n'a qu'un nombre fini de composantes irréductibles, disons $X_1, \ldots, X_n$,
et que $X = X_1\cup X_2\cup\ldots\cup X_n$. (Noter que
tout $X$ est toujours la réunion de ses composantes irréductibles, mais que,
si $X = {\mathbf R}$ est par exemple muni de sa topologie usuelle, les composantes irréductibles
de $X$ sont les parties réduites à un point. Ce n'est guère
intéressant.)
\end{exercise}

\begin{exercise}
\label{exercise-irreducible-components-minimal-primes}
Montrer que les composantes irréductibles de $\Spec(A)$
correspondent aux idéaux premiers minimaux de $A$.
\end{exercise}

\begin{definition}
\label{definition-closed}
Un point $x\in X$ est dit {\it fermé} si $\overline{\{x\}} = \{ x\}$.
Soient $x, y$ des points de $X$. On dit que $x$ est une {\it spécialisation}
de $y$, ou que $y$ est une {\it générisation} de $x$, si
$x\in \overline{\{y\}}$.
\end{definition}

\begin{exercise}
\label{exercise-closed-maximal}
Montrer que les points fermés de $\Spec(A)$
correspondent aux idéaux maximaux de $A$.
\end{exercise}

\begin{exercise}
\label{exercise-generalization}
Montrer que ${\mathfrak p}$ est une générisation de ${\mathfrak q}$
dans $\Spec(A)$ si et seulement si ${\mathfrak p}\subset {\mathfrak q}$.
Caractériser les points fermés,
les idéaux maximaux, les points génériques et les idéaux premiers minimaux en termes de
générisation et de spécialisation. (On utilise ici la terminologie selon laquelle un point
d'un espace topologique éventuellement réductible $X$ est appelé point générique
s'il est un point générique de l'une des composantes irréductibles de $X$.)
\end{exercise}

\begin{exercise}
\label{exercise-disjoint-closed-spec}
Soient $I$ et $J$ des idéaux de $A$.
À quelle condition $V(I)$ et $V(J)$ sont-ils disjoints ?
\end{exercise}

\begin{definition}
\label{definition-connected-component}
Un espace topologique $X$ est dit {\it connexe} s'il est non vide et n'est pas la
réunion de deux ouverts non vides et disjoints. Une {\it composante connexe}
de $X$ est une partie connexe maximale. Tout point de $X$ est contenu
dans une composante connexe de $X$, et toute composante connexe de $X$ est
fermée dans $X$. (Mais, en général, une composante connexe n'est pas nécessairement ouverte dans $X$.)
\end{definition}

\begin{exercise}
\label{exercise-disconnected-spec}
Soit $A$ un anneau non nul.
Montrer que $\Spec(A)$ est non connexe
si et seulement si $A\cong B \times C$ pour certains anneaux non nuls $B, C$.
\end{exercise}

\begin{exercise}
\label{exercise-connected-component-stable-generalization}
Soit $T$ une composante connexe
de $\Spec(A)$. Montrer que $T$ est stable par générisation.
Montrer que $T$ est une partie ouverte de $\Spec(A)$ si $A$ est noethérien.
(Remarque : cette assertion est fausse lorsque $A$ est, par exemple, un produit infini de copies de
${\mathbf F}_2$. Le spectre de cet anneau est constitué d'une infinité
de points fermés.)
\end{exercise}

\begin{exercise}
\label{exercise-primes-kx}
Déterminer $\Spec(k[x])$, c'est-à-dire décrire
les idéaux premiers de cet anneau, décrire les spécialisations possibles et
décrire la topologie. (Traiter le cas où $k$ est algébriquement clos, mais
aussi celui où $k$ ne l'est pas.)
\end{exercise}

\begin{exercise}
\label{exercise-primes-kxy}
 Déterminer $\Spec(k[x, y])$, où $k$ est algébriquement
clos.
[Indication : utiliser le morphisme
$\varphi : \Spec(k[x, y]) \to \Spec(k[x])$ ; si
$\varphi({\mathfrak p}) = (0)$, alors localiser suivant
$S = \{f\in k[x] \mid f \not = 0\}$ et utiliser le résultat du cours
sur la localisation et $\Spec$.]
(Pourquoi, à votre avis, les géomètres algébristes appellent-ils cet espace l'espace affine de dimension 2 ?)
\end{exercise}

\begin{exercise}
\label{exercise-primes-Zy}
Déterminer $\Spec(\mathbf{Z}[y])$.
[Indication : procéder comme ci-dessus.] (Espace affine de dimension 1 sur $\mathbf{Z}$.)
\end{exercise}






\section{Localisation}
\label{section-localization}


\begin{exercise}
\label{exercise-submodule-localization}
Soit $A$ un anneau. Soit $S \subset A$ une partie multiplicative.
Soit $M$ un $A$-module. Soit $N \subset S^{-1}M$ un $S^{-1}A$-sous-module.
Montrer qu'il existe un $A$-sous-module $N' \subset M$ tel que
$N = S^{-1}N'$. (Ce résultat utile s'applique en particulier aux idéaux
de $S^{-1}A$.)
\end{exercise}

\begin{exercise}
\label{exercise-localize-zero}
Soit $A$ un anneau. Soit $M$ un $A$-module. Soit $m \in M$.
\begin{enumerate}
\item Montrer que $I = \{a \in A \mid am = 0\}$ est un idéal de $A$.
\item Pour un idéal premier $\mathfrak p$ de $A$, montrer que l'image de $m$
dans $M_\mathfrak p$ est nulle si et seulement si $I \not \subset \mathfrak p$.
\item Montrer que $m$ est nul si et seulement si l'image de $m$ est nulle
dans $M_\mathfrak p$ pour tout idéal premier $\mathfrak p$ de $A$.
\item Montrer que $m$ est nul si et seulement si l'image de $m$ est nulle
dans $M_\mathfrak m$ pour tout idéal maximal $\mathfrak m$ de $A$.
\item Montrer que $M = 0$ si et seulement si $M_{\mathfrak m}$ est nul
pour tout idéal maximal $\mathfrak m$.
\end{enumerate}
\end{exercise}

\begin{exercise}
\label{exercise-localization-is-field}
Trouver un couple $(A, f)$ où $A$ est un anneau intègre possédant au moins trois idéaux
premiers deux à deux distincts et où $f \in A$ est tel que le localisé
principal $A_f = \{1, f, f^2, \ldots \}^{-1}A$ soit un corps.
\end{exercise}

\begin{exercise}
\label{exercise-localize-finite-module-zero}
Soit $A$ un anneau. Soit $M$ un $A$-module de type fini. Soit $S \subset A$
une partie multiplicative. Supposons que $S^{-1}M = 0$. Montrer qu'il
existe $f \in S$ tel que le localisé principal
$M_f = \{1, f, f^2, \ldots \}^{-1}M$ soit nul.
\end{exercise}
% Borne de la tranche : la ligne source 750 est vide.
\begin{exercise}
\label{exercise-localization-is-quotient}
Donner un exemple de triplet $(A, I, S)$ où $A$ est un anneau,
$0 \not = I \not = A$ est un idéal propre non nul et $S \subset A$
est une partie multiplicative
telle que $A/I \cong S^{-1}A$ comme $A$-algèbres.
\end{exercise}




\section{Lemme de Nakayama}
\label{section-nakayama}

\begin{exercise}
\label{exercise-nakayama}
Soit $A$ un anneau.
Soit $I$ un idéal de $A$.
Soit $M$ un $A$-module.
Soient $x_1, \ldots, x_n \in M$.
Supposons satisfaites les conditions suivantes :
\begin{enumerate}
\item $M/IM$ est engendré par $x_1, \ldots, x_n$,
\item $M$ est un $A$-module de type fini,
\item $I$ est contenu dans tout idéal maximal de $A$.
\end{enumerate}
Montrer que $x_1, \ldots, x_n$ engendrent $M$. (Solution suggérée :
se ramener, par localisation, au cas d'un idéal maximal de $A$ à l'aide de
l'exercice \ref{exercise-localize-zero} et de l'exactitude de la localisation.
Se ramener ensuite à l'énoncé du lemme de Nakayama donné dans le cours
en considérant le quotient de $M$ par le sous-module engendré par
$x_1, \ldots, x_n$.)
\end{exercise}






\section{Longueur}
\label{section-length}

\begin{definition}
\label{definition-length}
Soient $A$ un anneau et $M$ un $A$-module. La
{\it longueur} de $M$ comme $A$-module est
$$
\text{longueur}_A(M)
=
\sup
\{
n
\mid
\exists\ 0 = M_0 \subset M_1 \subset \ldots \subset M_n = M,
\text{ }M_i \not = M_{i + 1}
\}.
$$
Autrement dit, c'est la borne supérieure des longueurs des chaînes de sous-modules.
\end{definition}

\begin{exercise}
\label{exercise-length-is-one}
Montrer qu'un module $M$ sur un anneau $A$ est de longueur $1$ si et seulement s'il
est isomorphe à $A/\mathfrak m$ pour un idéal maximal $\mathfrak m$
de $A$.
\end{exercise}

\begin{exercise}
\label{exercise-length-easy}
Calculer la longueur des modules suivants sur les anneaux indiqués.
Expliquer brièvement (!) chaque réponse. (On pourra utiliser librement l'additivité de
la fonction longueur dans les suites exactes courtes ; voir
Algebra, lemme \ref{algebra-lemma-length-additive}.)
\begin{enumerate}
\item La longueur de $\mathbf{Z}/120\mathbf{Z}$ sur $\mathbf{Z}$.
\item La longueur de $\mathbf{C}[x]/(x^{100} + x + 1)$ sur $\mathbf{C}[x]$.
\item La longueur de $\mathbf{R}[x]/(x^4 + 2x^2 + 1)$ sur $\mathbf{R}[x]$.
\end{enumerate}
\end{exercise}

\begin{exercise}
\label{exercise-compute-length}
Soit $A = k[x, y]_{(x, y)}$ l'anneau local du plan affine en
l'origine. On pourra imposer au corps $k$ toute hypothèse que l'on souhaite. Supposons
que $f = x^3 + x^2y^2 + y^{100}$ et $g = y^3 - x^{999}$. Quelle est la longueur
de $A/(f, g)$ comme $A$-module ? (Une démarche possible consiste à considérer
l'idéal engendré par $f$ et $g$ dans les quotients de la forme $A/{\mathfrak m}_A^n=
k[x, y]/(x, y)^n$, où $n$ varie. Essayer de trouver $n$ tel que
$A/(f, g)+{\mathfrak m}_A^n \cong A/(f, g)+{\mathfrak m}_A^{n + 1}$
et utiliser NAK.)
\end{exercise}






\section{Idéaux premiers associés}
\label{section-ass}

\noindent
Les idéaux premiers associés sont étudiés dans
Algebra, section \ref{algebra-section-ass}.

\begin{exercise}
\label{exercise-compute-ass}
Calculer l'ensemble des idéaux premiers associés à chacun des modules suivants.
\begin{enumerate}
\item $R = k[x, y]$ et $M = R/(xy(x + y))$,
\item $R = \mathbf{Z}[x]$ et $M = R/(300x + 75)$, et
\item $R = k[x, y, z]$ et $M = R/(x^3, x^2y, xz)$.
\end{enumerate}
Comme d'habitude, $k$ désigne un corps.
\end{exercise}

\begin{exercise}
\label{exercise-prime-power-not-primary}
Donner un exemple d'anneau noethérien $R$ et d'idéal premier $\mathfrak p$
tel que $\mathfrak p$ ne soit pas le seul idéal premier associé à $R/\mathfrak p^2$.
\end{exercise}

\begin{exercise}
\label{exercise-product-primes-not-only-associated}
Soit $R$ un anneau noethérien possédant deux idéaux premiers incomparables
$\mathfrak p$, $\mathfrak q$, c'est-à-dire $\mathfrak p \not \subset \mathfrak q$
et $\mathfrak q \not \subset \mathfrak p$.
\begin{enumerate}
\item Montrer que, pour $N = R/(\mathfrak p \cap \mathfrak q)$,
on a $\text{Ass}(N) = \{\mathfrak p, \mathfrak q\}$.
\item Montrer par un exemple que le module
$M = R/\mathfrak p \mathfrak q$ peut
avoir un idéal premier associé distinct de $\mathfrak p$ et de $\mathfrak q$.
\end{enumerate}
\end{exercise}





\section{Groupes d'Ext}
\label{section-ext}

\noindent
Les groupes d'Ext sont définis dans Algebra, section \ref{algebra-section-ext}.

\begin{exercise}
\label{exercise-compute-ext-abelian-groups}
Calculer tous les groupes d'Ext $\Ext^i(M, N)$ des modules donnés
dans la catégorie des $\mathbf{Z}$-modules
(appelée aussi catégorie des groupes abéliens).
\begin{enumerate}
\item $M = \mathbf{Z}$ et $N = \mathbf{Z}$,
\item $M = \mathbf{Z}/4\mathbf{Z}$ et $N = \mathbf{Z}/8\mathbf{Z}$,
\item $M = \mathbf{Q}$ et $N = \mathbf{Z}/2\mathbf{Z}$, et
\item $M = \mathbf{Z}/2\mathbf{Z}$ et $N = \mathbf{Q}/\mathbf{Z}$.
\end{enumerate}
\end{exercise}

\begin{exercise}
\label{exercise-compute-in-regular}
Soit $R = k[x, y]$, où $k$ est un corps.
\begin{enumerate}
\item Montrer directement que le complexe de Koszul
$$
0 \to R
\xrightarrow{
\left(
\begin{matrix}
y \\
-x
\end{matrix}
\right)
}
R^{\oplus 2} \xrightarrow{(x, y)} R \xrightarrow{f \mapsto f(0, 0)} k \to 0
$$
est exact.
\item Calculer $\Ext^i_R(k, k)$, où $k = R/(x, y)$ est considéré comme un $R$-module.
\end{enumerate}
\end{exercise}

\begin{exercise}
\label{exercise-infinitely-many-nonzero-ext}
Donner un exemple d'anneau noethérien $R$ et de modules de type fini $M$, $N$
tels que $\Ext^i_R(M, N)$ soit non nul pour tout $i \geq 0$.
\end{exercise}

\begin{exercise}
\label{exercise-infinite-ext}
Donner un exemple d'anneau $R$ et d'idéal $I$ tels que
$\Ext^1_R(R/I, R/I)$ ne soit pas un $R$-module de type fini.
(On sait que cela ne peut pas se produire si $R$ est noethérien, d'après
Algebra, lemme \ref{algebra-lemma-ext-noetherian}.)
\end{exercise}







\section{Profondeur}
\label{section-depth}

\noindent
La profondeur est définie dans Algebra, section \ref{algebra-section-depth},
et étudiée plus avant dans
Dualizing Complexes, section \ref{dualizing-section-depth}.

\begin{exercise}
\label{exercise-compute-depth}
Soient $R$ un anneau, $I \subset R$ un idéal et $M$ un $R$-module.
Calculer $\text{prof}_I(M)$ dans les cas suivants.
\begin{enumerate}
\item $R = \mathbf{Z}$, $I = (30)$, $M = \mathbf{Z}$,
\item $R = \mathbf{Z}$, $I = (30)$, $M = \mathbf{Z}/(300)$,
\item $R = \mathbf{Z}$, $I = (30)$, $M = \mathbf{Z}/(7)$,
\item $R = k[x, y, z]/(x^2 + y^2 + z^2)$, $I = (x, y, z)$, $M = R$,
\item $R = k[x, y, z, w]/(xz, xw, yz, yw)$, $I = (x, y, z, w)$,
$M = R$.
\end{enumerate}
Ici, $k$ est un corps. Dans les deux derniers cas, on pourra localiser
en l'idéal maximal $I$.
\end{exercise}

\begin{exercise}
\label{exercise-depth-not-inherited-localization}
Donner un exemple d'anneau local noethérien $(R, \mathfrak m, \kappa)$
de profondeur $\geq 1$ et d'idéal premier $\mathfrak p$ tels que
\begin{enumerate}
\item $\text{prof}_\mathfrak m(R) \geq 1$,
\item $\text{prof}_\mathfrak p(R_\mathfrak p) = 0$, et
\item $\dim(R_\mathfrak p) \geq 1$.
\end{enumerate}
Si l'on ne demande pas (3), l'exercice est trop facile. Pourquoi ?
\end{exercise}

\begin{exercise}
\label{exercise-depth-torsion-free}
Soit $(R, \mathfrak m)$ un anneau local noethérien intègre.
Soit $M$ un $R$-module de type fini.
\begin{enumerate}
\item Si $M$ est sans torsion, montrer que $M$ est de profondeur au moins $1$ sur $R$.
\item Donner un exemple où la profondeur vaut $1$.
\end{enumerate}
\end{exercise}

\begin{exercise}
\label{exercise-depth-examples}
Pour tous $m \geq n \geq 0$, donner un exemple d'anneau local
noethérien $R$ tel que $\dim(R) = m$ et $\text{prof}(R) = n$.
\end{exercise}

\begin{exercise}
\label{exercise-make-depth-1}
Soit $(R, \mathfrak m)$ un anneau local noethérien.
Soit $M$ un $R$-module de type fini.
Montrer qu'il existe une suite exacte courte canonique
$$
0 \to K \to M \to Q \to 0
$$
ayant les propriétés suivantes :
\begin{enumerate}
\item $\text{prof}(Q) \geq 1$,
\item $K$ est nul ou $\text{Supp}(K) = \{\mathfrak m\}$, et
\item $\text{longueur}_R(K) < \infty$.
\end{enumerate}
Indication : en utilisant la propriété noethérienne, montrer qu'il
existe un sous-module $K$ maximal parmi ceux qui vérifient (2),
puis que $Q = M/K$ vérifie (1) et que $K$ vérifie (3).
\end{exercise}

\begin{exercise}
\label{exercise-make-depth-2}
Soit $(R, \mathfrak m)$ un anneau local noethérien.
Soit $M$ un $R$-module de type fini et de profondeur $\geq 2$.
Soit $N \subset M$ un sous-module non nul.
\begin{enumerate}
\item Montrer que $\text{prof}(N) \geq 1$.
\item Montrer que $\text{prof}(N) = 1$ si et seulement si le
module quotient $M/N$ vérifie $\text{prof}(M/N) = 0$.
\item Montrer qu'il existe un sous-module $N' \subset M$ tel que
$N \subset N'$ et que le quotient soit de longueur finie, c'est-à-dire que
$\text{longueur}_R(N'/N) < \infty$, et tel que $N'$ soit
de profondeur $\geq 2$. Indication : appliquer l'exercice \ref{exercise-make-depth-1} à $M/N$
et prendre pour $N'$ l'image inverse de $K$.
\end{enumerate}
\end{exercise}

\begin{exercise}
\label{exercise-Hartshorne-reduced}
Soit $(R, \mathfrak m)$ un anneau local noethérien.
Supposons $R$ réduit, c'est-à-dire que $R$ ne possède aucun élément
nilpotent non nul. Supposons en outre que $R$ possède deux
idéaux premiers minimaux distincts $\mathfrak p$ et $\mathfrak q$.
\begin{enumerate}
\item Montrer que la suite de $R$-modules
$$
0 \to R \to R/\mathfrak p \oplus R/\mathfrak q \to
R/\mathfrak p + \mathfrak q \to 0
$$
est exacte (le vérifier en chaque terme). Les applications sont
$x \mapsto (x \bmod \mathfrak p, x \bmod \mathfrak q)$
et $(y \bmod \mathfrak p, z \bmod \mathfrak q) \mapsto
(y - z \bmod \mathfrak p + \mathfrak q)$.
\item Montrer que si $\text{prof}(R) \geq 2$, alors
$\dim(R/\mathfrak p + \mathfrak q) \geq 1$.
\item Montrer que si $\text{prof}(R) \geq 2$, alors
$U = \Spec(R) \setminus \{\mathfrak m\}$ est un espace topologique
connexe.
\end{enumerate}
Cela démontre un cas très particulier du théorème de connexité de Hartshorne,
selon lequel le spectre épointé $U$ d'un anneau local
noethérien de $\text{prof} \geq 2$ est connexe.
\end{exercise}

\begin{exercise}
\label{exercise-depth-2}
Soit $(R, \mathfrak m)$ un anneau local noethérien.
Supposons que $x, y \in \mathfrak m$ forment une suite régulière de longueur $2$.
Pour tout $n \geq 2$, montrer qu'il n'existe pas de $a, b \in R$ tels que
$$
x^{n - 1}y^{n - 1} = a x^n + b y^n
$$
Indication : commencer par traiter le cas $n = 2$ afin de voir comment procéder.
Remarque : ce résultat admet une vaste généralisation
appelée la conjecture des monômes.
\end{exercise}







\section{Modules et anneaux de Cohen-Macaulay}
\label{section-CM}

\noindent
Les modules de Cohen-Macaulay sont étudiés dans
Algebra, section \ref{algebra-section-CM},
et les anneaux de Cohen-Macaulay sont étudiés dans
Algebra, section \ref{algebra-section-CM-ring}.

\begin{exercise}
\label{exercise-examples-CM}
Dans chacun des cas suivants, répondre par oui ou par non.
Aucune explication ni démonstration n'est nécessaire.
\begin{enumerate}
\item Soit $p$ un nombre premier.
L'anneau local $\mathbf{Z}_{(p)}$ est-il un anneau local de Cohen-Macaulay ?
\item Soit $p$ un nombre premier.
L'anneau local $\mathbf{Z}_{(p)}$ est-il un anneau local régulier ?
\item Soit $k$ un corps.
L'anneau local $k[x]_{(x)}$ est-il un anneau local de Cohen-Macaulay ?
\item Soit $k$ un corps.
L'anneau local $k[x]_{(x)}$ est-il un anneau local régulier ?
\item Soit $k$ un corps.
L'anneau local $(k[x, y]/(y^2 - x^3))_{(x, y)} =
k[x, y]_{(x, y)}/(y^2 - x^3)$ est-il un anneau local de Cohen-Macaulay ?
\item Soit $k$ un corps.
L'anneau local $(k[x, y]/(y^2, xy))_{(x, y)} =
k[x, y]_{(x, y)}/(y^2, xy)$ est-il un anneau local de Cohen-Macaulay ?
\end{enumerate}
\end{exercise}













\section{Singularités}
\label{section-singularities}

\begin{exercise}
\label{exercise-singularities}
Soit $k$ un corps quelconque. Supposons que $A = k[[x, y]]/(f)$ et
$B = k[[u, v]]/(g)$, où $f = xy$ et $g = uv + \delta$ avec
$\delta \in (u, v)^3$. Montrer que $A$ et $B$ sont des anneaux isomorphes.
\end{exercise}

\begin{remark}
\label{remark-singularities}
Une singularité d'une courbe sur un corps $k$ est appelée
point double ordinaire si l'anneau local complété de la courbe au
point est de la forme $k'[[x, y]]/(f)$, où (a) $k'$ est une extension finie
séparable de $k$, (b) le terme initial de $f$ est de degré deux, c'est-à-dire qu'il
est de la forme $q = ax^2 + bxy + cy^2$ pour des $a, b, c\in k'$ non tous nuls, et
(c) $q$ est une forme quadratique non dégénérée sur $k'$ (en caractéristique 2,
cela signifie que $b$ est non nul). En général, il y a une classe d'isomorphisme
de tels anneaux pour chaque classe d'isomorphisme de couples $(k', q)$.
\end{remark}

\begin{exercise}
\label{exercise-periodic-resolution}
Soit $R$ un anneau. Soit $n \geq 1$. Soient $A$, $B$ des matrices $n \times n$ à
coefficients dans $R$ telles que $AB = f 1_{n \times n}$
pour un non-diviseur de zéro $f$ de $R$. Posons $S = R/(f)$. Montrer que
$$
\ldots \to
S^{\oplus n} \xrightarrow{B}
S^{\oplus n} \xrightarrow{A}
S^{\oplus n} \xrightarrow{B}
S^{\oplus n} \to \ldots
$$
est exacte.
\end{exercise}







\section{Ensembles constructibles}
\label{section-constructible}

\noindent
Soit $k$ un corps algébriquement clos, par exemple le corps $\mathbf{C}$
des nombres complexes. Soit $n \geq 0$. Un polynôme $f \in k[x_1, \ldots, x_n]$
définit par évaluation une fonction $f : k^n \to k$. Un sous-ensemble $Z \subset k^n$
est appelé un {\it ensemble algébrique} s'il est l'ensemble des zéros communs d'une
famille de polynômes.

\begin{exercise}
\label{exercise-finite-nr-equations}
Démontrer qu'un ensemble algébrique peut toujours s'écrire comme l'ensemble des zéros
d'un nombre fini de polynômes.
\end{exercise}

\noindent
Avec les notations ci-dessus, un sous-ensemble $E \subset k^n$ est dit
{\it constructible} s'il est une réunion finie d'ensembles de la forme
$Z \cap \{f \not = 0\}$, où $f$ est un polynôme.

\begin{exercise}
\label{exercise-constructible-classical}
Montrer les assertions suivantes :
\begin{enumerate}
\item le complémentaire d'un ensemble constructible est constructible,
\item une réunion finie d'ensembles constructibles est constructible,
\item une intersection finie d'ensembles constructibles est constructible, et
\item tout ensemble constructible $E$ peut s'écrire comme une réunion disjointe finie
$E = \coprod E_i$, chaque $E_i$ étant de la forme $Z \cap \{f \not = 0\}$,
où $Z$ est un ensemble algébrique et $f$ un polynôme.
\end{enumerate}
\end{exercise}

\begin{exercise}
\label{exercise-division-with-remainder}
Soit $R$ un anneau. Soit
$f = a_d x^d + a_{d - 1} x^{d - 1} + \ldots + a_0 \in R[x]$.
(Comme d'habitude, cette notation signifie que $a_0, \ldots, a_d \in R$.) Soit $g \in R[x]$.
Démontrer qu'il existe $N \geq 0$ et $r, q \in R[x]$ tels que
$$
a_d^N g = q f + r
$$
avec $\deg(r) < d$, c'est-à-dire que, pour certains $c_i \in R$, on a
$r = c_0 + c_1 x + \ldots + c_{d - 1}x^{d - 1}$.
\end{exercise}







\section{Nullstellensatz de Hilbert}
\label{section-Hilbert-Nullstellensatz}


\begin{exercise}
\label{exercise-uncountable}
{\it Un argument saugrenu utilisant les nombres complexes !}
Soit ${\mathbf C}$ le corps des nombres complexes. Soit $V$ un espace vectoriel
sur ${\mathbf C}$. Le spectre d'un opérateur linéaire
$T : V \to V$ est l'ensemble des nombres complexes $\lambda \in {\mathbf C}$
tels que l'opérateur $T - \lambda \text{id}_V$ ne soit pas inversible.
\begin{enumerate}
\item Montrer que $\mathbf{C}(X)$
est de dimension non dénombrable sur ${\mathbf C}$.
\item Montrer que tout opérateur linéaire sur $V$ possède un
spectre non vide si la dimension de $V$ est finie ou
dénombrable.
\item Montrer que si une ${\mathbf C}$-algèbre de type fini
$R$ est un corps, alors l'application ${\mathbf C}\to R$ est un isomorphisme.
\item Montrer que tout idéal maximal ${\mathfrak m}$ de
${\mathbf C}[x_1, \ldots, x_n]$ est de la forme
$(x_1-\alpha_1, \ldots, x_n-\alpha_n)$ pour certains $\alpha_i \in {\mathbf C}$.
\end{enumerate}
\end{exercise}

\begin{remark}
\label{remark-HNSS}
Soit $k$ un corps. Alors, pour tout entier $n\in {\mathbf N}$ et
tout idéal maximal ${\mathfrak m} \subset k[x_1, \ldots, x_n]$,
le quotient $k[x_1, \ldots, x_n]/{\mathfrak m}$ est une extension finie
de $k$. Ce résultat sera démontré plus loin dans le cours. Bien entendu
(le vérifier), il entraîne un énoncé analogue pour les idéaux maximaux
des $k$-algèbres de type fini. L'exercice précédent le démontre
dans le cas $k = {\mathbf C}$.
\end{remark}

\begin{exercise}
\label{exercise-Hilbert-Nullstellensatz}
Soit $k$ un corps. Utiliser la remarque \ref{remark-HNSS}.
\begin{enumerate}
\item Soit $R$ une $k$-algèbre. Supposons que $\dim_k R < \infty$
et que $R$ soit intègre. Montrer que $R$ est un corps.
\item Supposons que $R$ soit une $k$-algèbre de type fini et que
$f\in R$ ne soit pas nilpotent. Montrer qu'il existe un idéal maximal
${\mathfrak m} \subset R$ tel que $f\not\in {\mathfrak m}$.
\item Montrer par un exemple que cette assertion est fausse lorsque $R$
n'est pas de type fini sur un corps.
\item Montrer que tout idéal radical $I \subset {\mathbf C}[x_1, \ldots, x_n]$
est l'intersection des idéaux maximaux qui le contiennent.
\end{enumerate}
\end{exercise}

\begin{remark}
\label{remark-Hilbert-Nullstellensatz}
C'est le Nullstellensatz de Hilbert. Il affirme en effet
que les parties fermées de $\Spec(k[x_1, \ldots, x_n])$
(qui correspondent aux idéaux radicaux d'après un exercice précédent)
sont déterminées par les points fermés qu'elles contiennent.
\end{remark}

\begin{exercise}
\label{exercise-product-matrices-ring}
Soit $A =
{\mathbf C}[x_{11}, x_{12}, x_{21}, x_{22}, y_{11}, y_{12}, y_{21}, y_{22}]$.
Soit $I$ l'idéal de $A$ engendré par les coefficients de la
matrice $XY$, où
$$
X = \left(
\begin{matrix}
x_{11} & x_{12}\\
x_{21} & x_{22}
\end{matrix}
\right)
\quad\text{et}\quad
Y = \left(
\begin{matrix}
y_{11} & y_{12}\\
y_{21} & y_{22}
\end{matrix}
\right).
$$
Déterminer les composantes irréductibles du sous-ensemble fermé $V(I)$ de
$\Spec(A)$.
(Il s'agit de les décrire et de donner des équations pour chacune. Il n'est pas
nécessaire de démontrer que les équations écrites définissent des idéaux premiers.) Indications :
\begin{enumerate}
\item On peut utiliser le Nullstellensatz de Hilbert, et il suffit de déterminer des
sous-ensembles localement fermés irréductibles qui recouvrent l'ensemble des points fermés de
$V(I)$.
\item Il y a deux composantes faciles à trouver.
\item L'image d'un ensemble irréductible par une application continue est
irréductible.
\end{enumerate}
\end{exercise}




\section{Dimension}
\label{section-dimension}

\begin{exercise}
\label{exercise-dimension-bigger-one-finite-nr-primes}
Construire un anneau $A$ de dimension $> 1$ qui n'a qu'un nombre fini d'idéaux premiers.
\end{exercise}

\begin{exercise}
\label{exercise-hypersurface-in-A2-dimension-one}
Soit $f \in \mathbf{C}[x, y]$ un polynôme non constant.
Montrer que $\mathbf{C}[x, y]/(f)$ est de dimension $1$.
\end{exercise}

\begin{exercise}
\label{exercise-dimension-polynomial-ring}
Soit $(R, \mathfrak m)$ un anneau local noethérien.
Soit $n \geq 1$. Soit $\mathfrak m' = (\mathfrak m, x_1, \ldots, x_n)$
dans l'anneau de polynômes $R[x_1, \ldots, x_n]$.
Montrer que
$$
\dim(R[x_1, \ldots, x_n]_{\mathfrak m'}) = \dim(R) + n.
$$
\end{exercise}





\section{Anneaux caténaires}
\label{section-catenary}

\begin{definition}
\label{definition-catenary}
On dit qu'un anneau noethérien $A$ est {\it caténaire}
si, pour tout triplet d'idéaux premiers
${\mathfrak p}_1 \subset {\mathfrak p}_2 \subset {\mathfrak p}_3$,
on a
$$
ht({\mathfrak p}_3 / {\mathfrak p}_1) = ht({\mathfrak p}_3/{\mathfrak p}_2) +
ht({\mathfrak p}_2/{\mathfrak p}_1).
$$
Ici, $ht(\mathfrak p/\mathfrak q)$ désigne la hauteur de
$\mathfrak p/\mathfrak q$ dans l'anneau $A/\mathfrak q$.
Autrement dit,
$$
ht(\mathfrak p/\mathfrak q) =
\dim(A_\mathfrak p/\mathfrak qA_\mathfrak p) =
\dim((A/\mathfrak q)_\mathfrak p) =
\dim((A/\mathfrak q)_{\mathfrak p/\mathfrak q})
$$
Un espace topologique $X$ est {\it caténaire} si, étant donnés $T \subset T' \subset X$
avec $T$ et $T'$ fermés et irréductibles, il existe une chaîne maximale
de parties fermées irréductibles
$$
T = T_0 \subset T_1 \subset \ldots \subset T_n = T'
$$
et toutes ces chaînes ont la même longueur (finie).
\end{definition}

\begin{exercise}
\label{exercise-catenary-the-same}
Montrer que la notion d'anneau caténaire définie dans
Algebra, définition \ref{algebra-definition-catenary}
coïncide avec celle de la définition \ref{definition-catenary}
pour les anneaux noethériens.
\end{exercise}

\begin{exercise}
\label{exercise-Noetherian-local-domain-dim-2-catenary}
Montrer qu'un anneau local noethérien intègre de dimension $2$ est caténaire.
\end{exercise}

\begin{exercise}
\label{exercise-finite-type-over-field-catenary}
Soit $k$ un corps.
Montrer qu'une $k$-algèbre de type fini est caténaire.
\end{exercise}

\begin{exercise}
\label{exercise-example-no-dim-function}
Donner un exemple d'espace topologique fini, sobre et caténaire
$X$ qui n'admette aucune fonction de dimension $\delta : X \to \mathbf{Z}$.
Ici, $\delta : X \to \mathbf{Z}$ est une fonction de dimension si,
pour $x, y \in X$, on a
\begin{enumerate}
\item $x \leadsto y$ et $x \not = y$ implique $\delta(x) > \delta(y)$,
\item $x \leadsto y$ et $\delta(x) \geq \delta(y) + 2$ implique
qu'il existe $z \in X$ tel que $x \leadsto z \leadsto y$ et
$\delta(x) > \delta(z) > \delta(y)$.
\end{enumerate}
Décrire clairement cet espace et expliquer brièvement pourquoi il
ne peut admettre de fonction de dimension.
\end{exercise}




\section{Corps des fractions}
\label{section-fraction-fields}

\begin{exercise}
\label{exercise-find-fraction-field}
Considérer l'anneau intègre
$$
{\mathbf Q}[r, s, t]/(s^2-(r-1)(r-2)(r-3), t^2-(r + 1)(r + 2)(r + 3)).
$$
Trouver un anneau intègre de la forme ${\mathbf Q}[x, y]/(f)$ dont le
corps des fractions soit isomorphe à celui de l'anneau ci-dessus.
\end{exercise}



\section{Degré de transcendance}
\label{section-transcendence}

\begin{exercise}
\label{exercise-algebraic-extension}
Soient $K'/K/k$ des extensions de corps telles que $K'$ soit
algébrique sur $K$. Montrer que $\text{deg.tr.}_k(K) = \text{deg.tr.}_k(K')$.
(Indication : montrer que si $x_1, \ldots, x_d \in K$ sont algébriquement indépendants
sur $k$ et si $d < \text{deg.tr.}_k(K')$, alors l'extension $k(x_1, \ldots, x_d) \subset K$
ne peut être algébrique.)
\end{exercise}

\begin{exercise}
\label{exercise-growth-powers-subvector-space}
Soient $k$ un corps et $K/k$ une extension de type fini
de degré de transcendance $d$.
Si $V, W \subset K$ sont des sous-espaces vectoriels sur $k$ de dimension finie,
posons
$$
VW = \{f \in K \mid f = \sum\nolimits_{i = 1, \ldots, n} v_i w_i
\text{ pour un certain }n\text{ et }v_i \in V, w_i \in W\}
$$
C'est un sous-espace vectoriel sur $k$ de dimension finie.
Posons $V^2 = VV$, $V^3 = V V^2$, etc.
\begin{enumerate}
\item Montrer qu'on peut trouver $V \subset K$ et $\epsilon > 0$ tels que
$\dim V^n \geq \epsilon n^d$ pour tout $n \geq 1$.
\item Réciproquement, montrer que, pour tout sous-espace $V \subset K$ de dimension finie,
il existe $C > 0$ tel que $\dim V^n \leq C n^d$ pour tout $n \geq 1$.
(Une démarche possible consiste à procéder comme suit :
établir d'abord le résultat pour les sous-espaces vectoriels de $k[x_1, \ldots, x_d]$,
puis pour ceux de $k(x_1, \ldots, x_d)$.
Enfin, si $K/k(x_1, \ldots, x_d)$ est une extension finie,
choisir une base de $K$ sur $k(x_1, \ldots, x_d)$ et
raisonner en développant les éléments sur cette base.)
\item En déduire que l'on peut redéfinir le degré de transcendance en fonction
de la croissance des puissances des sous-espaces vectoriels de dimension finie de $K$.
\end{enumerate}
Cela est lié à la dimension de Gelfand--Kirillov des algèbres
(non commutatives) sur $k$.
\end{exercise}








\section{Dimension des fibres}
\label{section-dimension-fibres}

\noindent
Quelques questions liées à la formule de dimension ; voir
Algebra, section \ref{algebra-section-dimension-formula}.

\begin{exercise}
\label{exercise-nr-components-fibre}
Soit $k$ un corps algébriquement clos quelconque.
Dans la suite, $k[x]$ et $k[x, y]$ désignent les anneaux de polynômes.
\begin{enumerate}
\item Pour tout entier $n \geq 0$, trouver une extension de type fini
$k[x] \subset A$ d'anneaux intègres telle que le spectre de $A/xA$
ait exactement $n$ composantes irréductibles.
\item Construire un exemple d'extension de type fini $k[x] \subset A$
d'anneaux intègres telle que le spectre de $A/(x - \alpha)A$
soit non vide et réductible pour tout $\alpha \in k$.
\item Construire un exemple d'extension de type fini $k[x, y] \subset A$
d'anneaux intègres telle que le spectre de $A/(x - \alpha, y - \beta)A$
soit irréductible\footnote{Rappelons qu'irréductible implique non vide.}
pour tout $(\alpha, \beta) \in k^2 \setminus \{(0, 0)\}$
et que le spectre de $A/(x, y)A$ soit non vide et réductible.
\end{enumerate}
\end{exercise}

\begin{exercise}
\label{exercise-codim-1}
Soit $k$ un corps algébriquement clos quelconque.
Soit $n \geq 1$.
Soit $k[x_1, \ldots, x_n]$ l'anneau de polynômes.
Posons $\mathfrak m = (x_1, \ldots, x_n)$.
Soit $k[x_1, \ldots, x_n] \subset A$ une extension de type fini
d'anneaux intègres. Posons $d = \dim(A)$.
\begin{enumerate}
\item Montrer que $d - 1 \geq \dim(A/\mathfrak m A) \geq d - n$
si $A/\mathfrak mA \not = 0$.
\item Montrer par des exemples que toutes les valeurs peuvent se produire.
\item Montrer par un exemple que $\Spec(A/\mathfrak m A)$ peut
avoir des composantes irréductibles de dimensions différentes.
\end{enumerate}
\end{exercise}




\section{Modules localement libres de type fini}
\label{section-finite-locally-free}

\begin{definition}
\label{definition-finite-locally-free}
Soit $A$ un anneau. Rappelons qu'un $A$-module $M$ est
{\it localement libre de type fini} si, pour tout ${\mathfrak p} \in \Spec(A)$,
il existe un
$f\in A$, $f \not \in {\mathfrak p}$, tel que $M_f$ soit un $A_f$-module libre
de type fini. On dit que $M$ est un {\it module inversible} si
$M$ est localement libre de type fini et de rang $1$, c'est-à-dire si, pour tout
${\mathfrak p} \in \Spec(A)$, il existe un
$f\in A$, $f \not \in \mathfrak p$, tel que $M_f \cong A_f$
comme $A_f$-module.
\end{definition}

\begin{exercise}
\label{exercise-tensor-finite-locally-free}
Montrer que le produit tensoriel de modules localement libres de type fini
est localement libre de type fini. Montrer que le produit tensoriel de deux
modules inversibles est inversible.
\end{exercise}

\begin{definition}
\label{definition-class-group}
Soit $A$ un anneau. Le {\it groupe des classes de $A$}, parfois appelé
{\it groupe de Picard de $A$}, est l'ensemble $\Pic(A)$
des classes d'isomorphie de $A$-modules inversibles, muni
de la loi de groupe définie par le produit tensoriel (voir
l'exercice \ref{exercise-tensor-finite-locally-free}).
\end{definition}

\noindent
Notons que le groupe des classes de $A$ est trivial si et seulement si tout module
inversible est isomorphe à un module libre de rang 1.

\begin{exercise}
\label{exercise-class-group-trivial}
Montrer que les groupes des classes des anneaux suivants sont triviaux :
\begin{enumerate}
\item un anneau de polynômes $A = k[x]$, où $k$ est un corps,
\item l'anneau des entiers $A = \mathbf{Z}$,
\item un anneau de polynômes $A = k[x, y]$, où $k$ est un corps, et
\item l'anneau quotient $k[x, y]/(xy)$, où $k$ est un corps.
\end{enumerate}
\end{exercise}

\begin{exercise}
\label{exercise-class-group-not-trivial}
Montrer que le groupe des classes de l'anneau
$A = k[x, y]/(y^2 - f(x))$, où $k$ est un corps de caractéristique différente de $2$
et où $f(x) = (x - t_1) \ldots (x - t_n)$ avec $t_1, \ldots, t_n \in k$
distincts et $n \geq 3$ impair, n'est pas trivial. (Indication : montrer que
l'idéal $(y, x - t_1)$ définit un élément non trivial de $\Pic(A)$.)
\end{exercise}

\begin{exercise}
\label{exercise-trace-det}
Soit $A$ un anneau.
\begin{enumerate}
\item Supposons que $M$ soit un $A$-module localement libre de type fini et
que $\varphi : M \to M$ soit un endomorphisme. Définir/construire
la {\it trace} et le {\it déterminant} de $\varphi$, et montrer que cette
construction est « fonctorielle par rapport au triplet $(A, M, \varphi)$ ».
\item Montrer que, si $M, N$ sont des $A$-modules localement libres de type fini,
et si $\varphi : M \to N$ et $\psi : N \to M$, alors
$\text{Trace}(\varphi \circ \psi) = \text{Trace}(\psi \circ \varphi)$ et
$\det(\varphi \circ \psi) = \det(\psi \circ \varphi)$.
\item Lorsque $M$ est localement libre de type fini, montrer que
$\text{Trace}$ définit une application $A$-linéaire $\text{End}_A(M) \to A$ et que
$\det$ définit une application multiplicative $\text{End}_A(M) \to A$.
\end{enumerate}
\end{exercise}

\begin{exercise}
\label{exercise-trace-det-rings}
Supposons maintenant que $B$ soit une $A$-algèbre localement libre
de type fini comme $A$-module ; autrement dit, $B$ est une $A$-algèbre
localement libre de type fini.
\begin{enumerate}
\item Définir $\text{Trace}_{B/A}$ et $\text{Norme}_{B/A}$ au moyen de
$\text{Trace}$ et de $\det$ dans l'exercice \ref{exercise-trace-det}.
\item Soit $b\in B$ et soit $\pi : \Spec(B) \to \Spec(A)$
le morphisme induit. Montrer que $\pi(V(b)) = V(\text{Norme}_{B/A}(b))$.
(Rappelons que $V(f) = \{ {\mathfrak p} \mid f \in {\mathfrak p}\}$.)
\item (Changement de base.) Supposons que $i : A \to A'$ soit un homomorphisme d'anneaux. Posons
$B' = B \otimes_A A'$. Indiquer pourquoi $i(\text{Norme}_{B/A}(b))$ est égal à
$\text{Norme}_{B'/A'}(b \otimes 1)$.
\item Calculer $\text{Norme}_{B/A}(b)$ lorsque
$B = A \times A \times A \times \ldots \times A$
et $b = (a_1, \ldots, a_n)$.
\item Calculer la norme de $y-y^3$ pour l'homomorphisme fini et plat
${\mathbf Q}[x] \to {\mathbf Q}[y]$, $x \to y^n$. (Indication : utiliser
le « changement de base »
$A = {\mathbf Q}[x] \subset A' = {\mathbf Q}(\zeta_n)(x^{1/n})$.)
\end{enumerate}
\end{exercise}



\section{Recollement}
\label{section-glueing}

\begin{exercise}
\label{exercise-cover}
Supposons que $A$ soit un anneau et que $M$ soit un $A$-module.
Soit donnée une famille d'éléments $f_i$, $i \in I$, de $A$ telle que
$$
\Spec(A) = \bigcup D(f_i).
$$
\begin{enumerate}
\item Montrer que, si $M_{f_i}$ est un $A_{f_i}$-module de type fini,
alors $M$ est un $A$-module de type fini.
\item Montrer que, si $M_{f_i}$ est un $A_{f_i}$-module plat,
alors $M$ est un $A$-module plat.
(C'est presque tautologique lorsqu'on y réfléchit.)
\end{enumerate}
\end{exercise}

\begin{remark}
\label{remark-cover}
Dans le langage de la géométrie algébrique, cela signifie que la propriété
d'« être de type fini » ou d'« être plat » est locale pour la topologie de Zariski
(en un sens approprié). On peut aussi le montrer pour la propriété
d'« être de présentation finie ».
\end{remark}

\begin{exercise}
\label{exercise-cover-ring-map}
Supposons que $A \to B$ soit un homomorphisme d'anneaux.
Soient $f_i \in A$, $i \in I$, et $g_j \in B$, $j \in J$, deux familles
d'éléments telles que
$$
\Spec(A) = \bigcup D(f_i)
\quad\text{et}\quad
\Spec(B) = \bigcup D(g_j).
$$
Montrer que, si $A_{f_i} \to B_{f_ig_j}$ est de type fini pour tous $i, j$,
alors $A \to B$ est de type fini.
\end{exercise}




\section{Montée et descente}
\label{section-going-up}


\begin{definition}
\label{definition-GU-GD}
Soit $\phi : A \to B$ un homomorphisme d'anneaux. On dit
que le {\it théorème de montée} est valable pour $\phi$ si la
condition suivante est satisfaite :
\begin{itemize}
\item[(GU)] pour tous ${\mathfrak p}, {\mathfrak p}' \in \Spec(A)$ tels que
${\mathfrak p} \subset {\mathfrak p}'$, et pour tout $P \in \Spec(B)$ situé
au-dessus de ${\mathfrak p}$, il existe $P'\in \Spec(B)$ situé
au-dessus de ${\mathfrak p}'$ tel que $P \subset P'$.
\end{itemize}
De même, on dit que le {\it théorème de descente} est valable pour $\phi$
si la condition suivante est satisfaite :
\begin{itemize}
\item[(GD)] pour tous ${\mathfrak p}, {\mathfrak p}' \in \Spec(A)$ tels que
${\mathfrak p} \subset {\mathfrak p}'$, et pour tout
$P' \in \Spec(B)$ situé
au-dessus de ${\mathfrak p}'$, il existe $P\in \Spec(B)$ situé
au-dessus de ${\mathfrak p}$ tel que $P \subset P'$.
\end{itemize}
\end{definition}

\begin{exercise}
\label{exercise-GU-GD}
Dans chacun des cas suivants, déterminer lesquelles des conditions
(GU) et (GD) sont satisfaites, et expliquer pourquoi. (Utiliser toute proposition,
tout théorème ou tout lemme disponible, mais vérifier les hypothèses dans chaque cas.)
\begin{enumerate}
\item $k$ est un corps, $A = k$, $B = k[x]$.
\item $k$ est un corps, $A = k[x]$, $B = k[x, y]$.
\item $A = {\mathbf Z}$, $B = {\mathbf Z}[1/11]$.
\item $k$ est un corps algébriquement clos, $A = k[x, y]$,
$B = k[x, y, z]/(x^2-y, z^2-x)$.
\item $A = {\mathbf Z}$, $B = {\mathbf Z}[i, 1/(2 + i)]$.
\item $A = {\mathbf Z}$, $B = {\mathbf Z}[i, 1/(14 + 7i)]$.
\item $k$ est un corps algébriquement clos, $A = k[x]$,
$B = k[x, y, 1/(xy-1)]/(y^2-y)$.
\end{enumerate}
\end{exercise}

\begin{exercise}
\label{exercise-image}
Soit $A$ un anneau. Soit $B = A[x]$ l'algèbre polynomiale à
une indéterminée sur $A$. Soit $f = a_0 + a_1 x + \ldots + a_r x^r \in B = A[x]$.
Démontrer soigneusement que l'image de $D(f)$ dans $\Spec(A)$ est égale à
$D(a_0) \cup \ldots \cup D(a_r)$.
\end{exercise}

\begin{exercise}
\label{exercise-images}
Soit $k$ un corps algébriquement clos. Calculer l'image dans
$\Spec(k[x, y])$
des morphismes suivants :
\begin{enumerate}
\item $\Spec(k[x, yx^{-1}]) \to \Spec(k[x, y])$, où
$k[x, y] \subset k[x, yx^{-1}] \subset k[x, y, x^{-1}]$.
(Indication : pour éviter toute confusion, donner un autre nom à l'élément $yx^{-1}$.)
\item $\Spec(k[x, y, a, b]/(ax-by-1))\to \Spec(k[x, y])$.
\item $\Spec(k[t, 1/(t-1)]) \to \Spec(k[x, y])$, induit par $x
\mapsto t^2$,
et $y \mapsto t^3$.
\item $k = {\mathbf C}$ (corps des nombres complexes),
$\Spec(k[s, t]/(s^3 + t^3-1)) \to \Spec(k[x, y])$, où
$x\mapsto s^2$, $y \mapsto t^2$.
\end{enumerate}
\end{exercise}

\begin{remark}
\label{remark-elimination-theory}
La détermination de l'image comme ci-dessus se fait habituellement au moyen de la théorie de l'élimination.
\end{remark}




\section{Idéaux de Fitting}
\label{section-fitting-ideals}

\begin{exercise}
\label{exercise-fitting}
Soient $R$ un anneau et $M$ un $R$-module de type fini.
Choisir une présentation
$$
\bigoplus\nolimits_{j \in J} R \longrightarrow R^{\oplus n}
\longrightarrow M \longrightarrow 0.
$$
de $M$. Noter que l'homomorphisme $R^{\oplus n} \to M$ est défini par une famille
d'éléments $x_1, \ldots, x_n$ de $M$. Les éléments $x_i$
sont des {\it générateurs} de $M$. L'homomorphisme $\bigoplus_{j \in J} R \to R^{\oplus n}$
est défini par une matrice $n \times J$ notée $A$, à coefficients dans $R$.
Autrement dit, $A = (a_{ij})_{i = 1, \ldots, n, j \in J}$.
Les colonnes $(a_{1j}, \ldots, a_{nj})$, $j \in J$, de $A$
sont appelées les {\it relations}. Tout vecteur $(r_i) \in R^{\oplus n}$
tel que $\sum r_i x_i = 0$ est une combinaison linéaire des colonnes de $A$.
Bien entendu, tout $R$-module de type fini possède de nombreuses présentations différentes.
\begin{enumerate}
\item Montrer que l'idéal engendré par les mineurs $(n - k) \times (n - k)$ de
$A$ est indépendant du choix de la présentation.
Cet idéal est l'{\it idéal de Fitting d'indice $k$ de $M$}. On le note $Fit_k(M)$.
\item Montrer que
$Fit_0(M) \subset Fit_1(M) \subset Fit_2(M) \subset \ldots$.
(Indication : utiliser le développement d'un déterminant suivant une colonne.)
\item Montrer que les conditions suivantes sont équivalentes :
\begin{enumerate}
\item $Fit_{r - 1}(M) = (0)$ et $Fit_r(M) = R$, et
\item $M$ est localement libre de rang $r$.
\end{enumerate}
\end{enumerate}
\end{exercise}




\section{Fonctions de Hilbert}
\label{section-hilbert}

\begin{definition}
\label{definition-numerical-polynomial}
Un {\it polynôme numérique} est un polynôme $f(x) \in {\mathbf Q}[x]$
tel que $f(n) \in {\mathbf Z}$ pour tout entier $n$.
\end{definition}

\begin{definition}
\label{definition-graded-module}
Un {\it module gradué} $M$ sur un anneau $A$ est un $A$-module $M$
muni d'une décomposition en somme directe de la forme
$
\bigoplus\nolimits_{n \in {\mathbf Z}} M_n
$
en sous-$A$-modules. On dit que $M$ est {\it localement fini} si tous les
$M_n$ sont des $A$-modules de type fini. Supposons que $A$ soit un anneau noethérien et
que $\varphi$ soit une {\it fonction d'Euler-Poincar\'e} sur les $A$-modules de type fini.
Cela signifie que, pour tout $A$-module de type fini $M$, on se donne un
entier $\varphi(M) \in {\mathbf Z}$ et que, pour toute suite exacte courte
$$
0
\longrightarrow
M'
\longrightarrow
M
\longrightarrow
M''
\longrightarrow
0
$$
on a $\varphi(M) = \varphi(M') + \varphi(M'')$. La {\it fonction de Hilbert}
d'un module gradué localement fini $M$ (relativement à $\varphi$) est la
fonction $\chi_\varphi(M, n) = \varphi(M_n)$. On dit que $M$ admet un
{\it polynôme de Hilbert} s'il existe un polynôme numérique
$P_\varphi$ tel que $\chi_\varphi(M, n) = P_\varphi(n)$ pour tout entier $n$
assez grand.
\end{definition}

\begin{definition}
\label{definition-graded-algebra}
Une {\it $A$-algèbre graduée} est un $A$-module gradué
$B = \bigoplus_{n \geq 0} B_n$ muni d'une application $A$-bilinéaire
$$
B \times B \longrightarrow B, \ (b, b') \longmapsto bb'
$$
qui munit $B$ d'une structure de $A$-algèbre telle que $B_n \cdot B_m \subset B_{n + m}$.
Enfin, un {\it module gradué $M$ sur une $A$-algèbre graduée $B$} est la donnée
d'un $A$-module gradué $M$ muni d'une structure (compatible) de $B$-module
telle que $B_n \cdot M_d \subset M_{n + d}$. On peut alors définir les
{\it homomorphismes de modules/anneaux gradués}, les {\it sous-modules gradués},
les {\it idéaux gradués}, les {\it suites exactes de modules gradués}, etc.
\end{definition}

\begin{exercise}
\label{exercise-Euler-Poincare-field}
Soit le corps $A = k$. Quelles sont toutes les fonctions d'Euler-Poincar\'e
possibles sur les $A$-modules de type fini dans ce cas ?
\end{exercise}

\begin{exercise}
\label{exercise-Euler-Poincare-Z}
Soit $A ={\mathbf Z}$. Quelles sont toutes les fonctions d'Euler-Poincar\'e
possibles sur les $A$-modules de type fini dans ce cas ?
\end{exercise}

\begin{exercise}
\label{exercise-Euler-Poincare-node}
Soit $A = k[x, y]/(xy)$, où $k$ est un corps algébriquement clos. Quelles sont, dans ce cas, toutes les
fonctions d'Euler-Poincar\'e possibles sur les $A$-modules de type fini ?
\end{exercise}

\begin{exercise}
\label{exercise-kernel-locally-finite}
Supposons que $A$ soit noethérien. Montrer que le noyau d'un homomorphisme
de $A$-modules gradués localement finis est localement fini.
\end{exercise}

\begin{exercise}
\label{exercise-no-hilbert}
Soit $k$ un corps, posons $A = k$ et munissons $B = k[x, y]$ de la graduation
déterminée par $\deg(x) = 2$ et $\deg(y) = 3$. Soit $\varphi(M) = \dim_k(M)$.
Calculer la fonction de Hilbert de $B$ comme $k$-module gradué. Existe-t-il
un polynôme de Hilbert dans ce cas ?
\end{exercise}

\begin{exercise}
\label{exercise-no-hilbert-or-is-there}
Soit $k$ un corps, posons $A = k$ et munissons $B = k[x, y]/(x^2, xy)$ de la graduation
déterminée par $\deg(x) = 2$ et $\deg(y) = 3$. Soit $\varphi(M) = \dim_k(M)$.
Calculer la fonction de Hilbert de $B$ comme $k$-module gradué. Existe-t-il
un polynôme de Hilbert dans ce cas ?
\end{exercise}

\begin{exercise}
\label{exercise-hilbert-to-compute}
Soit $k$ un corps et posons $A = k$. Soit $\varphi(M) = \dim_k(M)$.
Fixons $d\in {\mathbf N}$. Considérons l'algèbre graduée sur $A$
$B = k[x, y, z]/(x^d + y^d + z^d)$, où $x, y, z$ sont tous de degré $1$.
Calculer la fonction de Hilbert de $B$. Existe-t-il un polynôme de Hilbert
dans ce cas ?
\end{exercise}



\section{Proj d'un anneau}
\label{section-proj-ring}

\begin{definition}
\label{definition-homogeneous-ideal}
Soit $R$ un anneau gradué. Un idéal {\it homogène} est simplement un idéal
$I \subset R$ qui est en outre un sous-module gradué de $R$. De manière équivalente,
c'est un idéal engendré par des éléments homogènes. De manière équivalente, si
$f \in I$ et si
$$
f = f_0 + f_1 + \ldots + f_n
$$
est la décomposition de $f$ en composantes homogènes dans $R$, alors $f_i \in I$
pour tout $i$.
\end{definition}

\begin{definition}
\label{definition-Proj-R}
On définit le {\it spectre homogène $\text{Proj}(R)$}
de l'anneau gradué $R$ comme l'ensemble des idéaux premiers homogènes
${\mathfrak p}$ de $R$ tels que $R_{+} \not \subset {\mathfrak p}$.
Noter que $\text{Proj}(R)$ est une partie de $\Spec(R)$ et porte donc naturellement
la topologie induite.
\end{definition}

\begin{definition}
\label{definition-Dplus-Vplus}
Soit $R = \oplus_{d \geq 0} R_d$ un anneau gradué, soit $f\in R_d$ et
supposons que $d \geq 1$. On définit {\it $R_{(f)}$} comme le sous-anneau de
$R_f$ formé des éléments de la forme $r/f^n$, où $r$ est homogène et
$\deg(r) = nd$. On définit en outre
$$
D_{+}(f) = \{ {\mathfrak p} \in \text{Proj}(R) | f \not\in {\mathfrak p} \}.
$$
Enfin, pour un idéal homogène $I \subset R$, on définit
$V_{+}(I) = V(I) \cap \text{Proj}(R)$.
\end{definition}

\begin{exercise}
\label{exercise-topology-proj}
Sur la topologie de $\text{Proj}(R)$. Avec les définitions et les notations
ci-dessus, démontrer les assertions suivantes.
\begin{enumerate}
\item Montrer que $D_{+}(f)$ est ouvert dans $\text{Proj}(R)$.
\item Montrer que $D_{+}(ff') = D_{+}(f) \cap D_{+}(f')$.
\item Soit $g = g_0 + \ldots + g_m$ un élément
de $R$ avec $g_i \in R_i$. Exprimer $D(g) \cap \text{Proj}(R)$
en fonction des $D_{+}(g_i)$, $i \geq 1$, et de $D(g_0) \cap \text{Proj}(R)$.
Aucune démonstration n'est nécessaire.
\item Soit $g\in R_0$ un élément homogène de degré $0$.
Exprimer $D(g) \cap \text{Proj}(R)$ en fonction des $D_{+}(f_\alpha)$
pour une famille convenable d'éléments homogènes $f_\alpha \in R$
de degré strictement positif.
\item Montrer que la famille $\{D_{+}(f)\}$ d'ouverts forme une
base de la topologie de $\text{Proj}(R)$.
\item
\label{item-bijection}
Montrer qu'il existe une bijection canonique $D_{+}(f) \to \Spec(R_{(f)})$.
(Indication : imiter la démonstration pour $\Spec$, mais, à un certain stade, faire intervenir le
radical d'un idéal.)
\item Montrer que l'application définie en (\ref{item-bijection}) est un homéomorphisme.
\item Donner un exemple d'un anneau $R$ tel que $\text{Proj}(R)$ ne soit pas
quasi-compact. Aucune démonstration n'est nécessaire.
\item Montrer que toute partie fermée $T \subset \text{Proj}(R)$ est de
la forme $V_{+}(I)$ pour un idéal homogène $I \subset R$.
\end{enumerate}
\end{exercise}

\begin{remark}
\label{remark-continuous-proj-spec}
Il existe une application continue $ \text{Proj}(R) \longrightarrow \Spec(R_0) $.
\end{remark}

\begin{exercise}
\label{exercise-iso-polynomial-ring-one-variable}
Si $R = A[X]$ avec $\deg(X) = 1$, montrer que l'application naturelle
$\text{Proj}(R) \to \Spec(A)$ est une bijection et même
un homéomorphisme.
\end{exercise}

\begin{exercise}
\label{exercise-blowing-up-I}
Éclatement : première partie.
Dans cet exercice, $R = Bl_I(A) = A \oplus I \oplus I^2 \oplus \ldots$.
Considérer l'application naturelle $b : \text{Proj}(R) \to \Spec(A)$.
Poser $U = \Spec(A) - V(I)$. Montrer que
$$
b : b^{-1}(U) \longrightarrow U
$$
est un homéomorphisme.
On peut donc identifier $U$ à une partie ouverte de $\text{Proj}(R)$.
Soit $Z \subset \Spec(A)$ un sous-schéma fermé irréductible
de point générique $\xi \in Z$. Supposer que $\xi \not\in V(I)$,
autrement dit $Z \not\subset V(I)$, autrement dit
$\xi \in U$, autrement dit $Z\cap U \not = \emptyset$. On définit
le {\it transformé pur} $Z'$ de $Z$ comme l'adhérence de l'unique
point $\xi'$ situé au-dessus de $\xi$. Autrement dit,
$Z'$ est l'adhérence dans $\text{Proj}(R)$ de la partie localement fermée
$Z\cap U \subset U \subset \text{Proj}(R)$.
\end{exercise}

\begin{exercise}
\label{exercise-blowing-up-II}
Éclatement : deuxième partie.
Soit $A = k[x, y]$, où $k$ est un corps, et soit $I = (x, y)$.
Soit $R$ l'algèbre d'éclatement associée à $A$ et $I$.
\begin{enumerate}
\item Montrer que les transformés purs de $Z_1 = V(\{x\})$ et de
$Z_2 = V(\{y\})$ sont disjoints.
\item Montrer que les transformés purs de $Z_1 = V(\{x\})$ et de
$Z_2 = V(\{x-y^2\})$ ne sont pas disjoints.
\item Trouver un idéal $J \subset A$ tel que $V(J) = V(I)$
et tel que les transformés purs de $Z_1 = V(\{x\})$ et de
$Z_2 = V(\{x-y^2\})$ dans l'éclatement le long de $J$ soient disjoints.
\end{enumerate}
\end{exercise}

\begin{exercise}
\label{exercise-proj-when-empty}
Soit $R$ un anneau gradué.
\begin{enumerate}
\item Montrer que $\text{Proj}(R)$ est vide si $R_n = (0)$ pour tout $n >> 0$.
\item Montrer que $\text{Proj}(R)$ est un espace topologique irréductible
si $R$ est intègre et si $R_{+}$ n'est pas nul. (Rappelons que l'espace
topologique vide n'est pas irréductible.)
\end{enumerate}
\end{exercise}

\begin{exercise}
\label{exercise-blowing-up-III}
Éclatement : troisième partie.
On reprend $A$, $I$, $U$ et $Z$ comme dans la définition du transformé pur.
Soit $Z = V({\mathfrak p})$ pour un idéal premier ${\mathfrak p}$. Poser $\bar A =
A/{\mathfrak p}$ et noter
$\bar I$ l'image de $I$ dans $\bar A$.
\begin{enumerate}
\item Montrer qu'il existe un homomorphisme d'anneaux surjectif
$R: = Bl_I(A) \to \bar R: = Bl_{\bar I}(\bar A)$.
\item Montrer que l'homomorphisme d'anneaux ci-dessus induit une application bijective
de $\text{Proj}(\bar R)$ sur le transformé pur $Z'$ de $Z$. (Cette question
n'est pas si facile. Indication : utiliser 5(b) ci-dessus.)
\item En déduire que le transformé pur $Z' = V_{+}(P)$, où
$P \subset R$ est l'idéal homogène défini par
$P_d = I^d \cap {\mathfrak p}$.
\item Supposer que $Z_1 = V({\mathfrak p})$ et
$Z_2 = V({\mathfrak q})$ sont des parties fermées irréductibles
définies par des idéaux premiers, que $Z_1 \not \subset Z_2$
et que $Z_2 \not \subset Z_1$. Montrer que l'éclatement de l'idéal
$I = {\mathfrak p} + {\mathfrak q}$ sépare les
transformés purs de $Z_1$ et de $Z_2$,
c'est-à-dire $Z_1' \cap Z_2' = \emptyset$. (Indication : considérer les idéaux homogènes
$P$ et $Q$ de la partie (c), puis $V(P + Q)$.)
\end{enumerate}
\end{exercise}


\section{Anneaux de Cohen-Macaulay de dimension 1}
\label{section-CM-dim-1}

\begin{definition}
\label{definition-CM}
On dit qu'un anneau local noethérien $A$ est un {\it anneau de Cohen-Macaulay}
de dimension $d$ s'il est de dimension $d$ et s'il existe un système
de paramètres $x_1, \ldots, x_d$ pour $A$ tel que $x_i$ soit un non-diviseur de zéro
dans $A/(x_1, \ldots, x_{i-1})$ pour $i = 1, \ldots, d$.
\end{definition}

\begin{exercise}
\label{exercise-CM-dim-1-I}
Anneaux de Cohen-Macaulay de dimension 1. Première partie : théorie.
\begin{enumerate}
\item Soit $(A, {\mathfrak m})$ un anneau local noethérien tel que $\dim A = 1$.
Montrer que, si $x\in {\mathfrak m}$ n'est pas un diviseur de zéro, alors
\begin{enumerate}
\item $\dim A/xA = 0$, autrement dit $A/xA$ est artinien,
autrement dit $\{x\}$ est un système de paramètres pour $A$.
\item $A$ n'a aucun idéal premier associé immergé.
\end{enumerate}
\item Réciproquement, soit $(A, {\mathfrak m})$ un anneau local noethérien de
dimension $1$. Montrer que, si $A$ n'a aucun idéal premier associé immergé, il existe
un non-diviseur de zéro dans ${\mathfrak m}$.
\end{enumerate}
\end{exercise}

\begin{exercise}
\label{exercise-CM-dim-1-II}
Anneaux de Cohen-Macaulay de dimension 1. Deuxième partie : exemples.
\begin{enumerate}
\item Soit $A$ l'anneau local en $(x, y)$ de $k[x, y]/(x^2, xy)$.
\begin{enumerate}
\item Montrer que $A$ est de dimension 1.
\item Démontrer que tout élément de ${\mathfrak m}\subset A$ est un
diviseur de zéro.
\item Trouver $z\in {\mathfrak m}$ tel que $\dim A/zA = 0$
(aucune démonstration n'est nécessaire).
\end{enumerate}
\item Soit $A$ l'anneau local en $(x, y)$ de $k[x, y]/(x^2)$.
Trouver un non-diviseur de zéro dans ${\mathfrak m}$ (aucune démonstration n'est nécessaire).
\end{enumerate}
\end{exercise}

\begin{exercise}
\label{exercise-embedding-dim-1}
Anneaux locaux de dimension de plongement $1$.
Supposer que $(A, {\mathfrak m}, k)$ soit un anneau local noethérien
de dimension de plongement $1$, c'est-à-dire
$$
\dim_k {\mathfrak m}/{\mathfrak m}^2 = 1.
$$
Montrer que la fonction $f(n) = \dim_k {\mathfrak m}^n/{\mathfrak m}^{n + 1}$
est soit constante de valeur $1$, soit telle que ses valeurs soient
$$
1, 1, \ldots, 1, 0, 0, 0, 0, 0, \ldots
$$
\end{exercise}

\begin{exercise}
\label{exercise-regular-local-dim-1}
Anneaux locaux réguliers de dimension $1$.
Supposer que $(A, {\mathfrak m}, k)$ soit un anneau local noethérien régulier de
dimension $1$. Rappelons que cela signifie que $A$ est de dimension $1$
et de dimension de plongement $1$, c'est-à-dire
$$
\dim_k {\mathfrak m}/{\mathfrak m}^2 = 1.
$$
Soit $x\in{\mathfrak m}$ un élément dont la classe dans ${\mathfrak m}/{\mathfrak
m}^2$ n'est pas nulle.
\begin{enumerate}
\item Montrer que, pour tout élément $y$
de ${\mathfrak m}$, il existe un entier $n$ tel que $y$ puisse s'écrire
$y = ux^n$, où $u\in A^\ast$ est une unité.
\item Montrer que $x$ est un non-diviseur de zéro dans $A$.
\item En déduire que $A$ est intègre.
\end{enumerate}
\end{exercise}

\begin{exercise}
\label{exercise-nonzerodivisor-graded}
Soit $(A, {\mathfrak m}, k)$ un anneau local noethérien, et soit
$Gr_{\mathfrak m}(A)$ son gradué associé.
\begin{enumerate}
\item Supposer que $x\in {\mathfrak m}^d$ a pour image un non-diviseur de zéro
$\bar x \in {\mathfrak m}^d/{\mathfrak m}^{d + 1}$ dans la composante de degré $d$ de
$Gr_{\mathfrak m}(A)$.
Montrer que $x$ est un non-diviseur de zéro.
\item Supposer que la profondeur de $A$ est au moins $1$.
Autrement dit, supposer qu'il existe un non-diviseur de zéro $y \in {\mathfrak m}$.
Dans ce cas, on peut améliorer le résultat : supposer seulement que $x\in {\mathfrak m}^d$ a pour image
l'élément $\bar x \in {\mathfrak m}^d/{\mathfrak m}^{d + 1}$ de degré $d$
de $Gr_{\mathfrak m}(A)$, lequel est un non-diviseur de zéro en degrés assez
élevés : $\exists N$ tel que, pour tout $n \geq N$, l'application
de multiplication par $\bar x$
$$
{\mathfrak m}^n/{\mathfrak m}^{n + 1} \longrightarrow
{\mathfrak m}^{n + d}/{\mathfrak m}^{n + d + 1}
$$
soit injective. Montrer alors que $x$ est un non-diviseur de zéro.
\end{enumerate}
\end{exercise}

\begin{exercise}
\label{exercise-embedding-2-dim-1}
Supposer que $(A, {\mathfrak m}, k)$ soit un anneau local noethérien de
dimension $1$. Supposer également que la dimension de plongement de $A$ soit
$2$, c'est-à-dire supposer que
$$
\dim_k {\mathfrak m}/{\mathfrak m}^2 = 2.
$$
Notation : $f(n) = \dim_k {\mathfrak m}^n/{\mathfrak m}^{n + 1}$.
Choisir des générateurs $x, y \in {\mathfrak m}$
et écrire $Gr_{\mathfrak m}(A) = k[\bar x, \bar y]/I$ pour un certain
idéal homogène $I$.
\begin{enumerate}
\item Montrer qu'il existe un élément homogène
$F\in k[\bar x, \bar y]$ tel que $I \subset (F)$, avec égalité
en tout degré assez élevé.
\item Montrer que $f(n) \leq n + 1$.
\item Montrer que, si $f(n) < n + 1$, alors $n \geq \deg(F)$.
\item Montrer que, si $f(n) < n + 1$, alors $f(n + 1) \leq f(n)$.
\item Montrer que $f(n) = \deg(F)$ pour tout $n >> 0$.
\end{enumerate}
\end{exercise}

\begin{exercise}
\label{exercise-CM-dim-1-embedding-dim-2}
Anneaux de Cohen-Macaulay de dimension 1 et de dimension de plongement 2.
Supposer que $(A, {\mathfrak m}, k)$ soit un anneau local noethérien de
Cohen-Macaulay
et de dimension $1$. Supposer également que la dimension de plongement de $A$ soit
$2$, c'est-à-dire supposer que
$$
\dim_k {\mathfrak m}/{\mathfrak m}^2 = 2.
$$
Notations : $f$, $F$, $x, y\in {\mathfrak m}$ et $I$ comme dans
l'exercice \ref{exercise-embedding-2-dim-1}. On pourra utiliser librement
tous les résultats des problèmes ci-dessus.
\begin{enumerate}
\item Supposer que $z\in {\mathfrak m}$ est un élément dont la classe
dans ${\mathfrak m}/{\mathfrak m}^2$ est la forme linéaire
$\alpha \bar x + \beta \bar y \in k[\bar x, \bar y]$
qui est première à $F$.
\begin{enumerate}
\item Montrer que $z$ est un non-diviseur de zéro dans $A$.
\item Poser $d = \deg(F)$.
Montrer que ${\mathfrak m}^n = z^{n + 1-d}{\mathfrak m}^{d-1}$
pour tout $n$ assez grand. (Indication : montrer d'abord que
$z^{n + 1-d}{\mathfrak m}^{d-1} \to {\mathfrak m}^n/{\mathfrak m}^{n + 1}$
est surjective en utilisant ce que l'on sait de $Gr_{\mathfrak m}(A)$. Puis utiliser le lemme de Nakayama.)
\end{enumerate}
\item Quelle condition sur $k$ garantit l'existence
d'un tel $z$ ? (Aucune démonstration n'est nécessaire; c'est trop facile.)

\noindent
Supposons désormais qu'il existe un $z$ comme ci-dessus. Cette hypothèse
est en fait sans restriction essentielle (au sens où l'on peut se ramener à
la situation où elle est satisfaite afin d'obtenir les résultats des
parties (d) et (e) ci-dessous).
\item Montrer maintenant que ${\mathfrak m}^\ell = z^{\ell - d + 1} {\mathfrak m}^{d-1}$
pour tout $\ell \geq d$.
\item En déduire que $I = (F)$.
\item En déduire que la fonction $f$ prend les valeurs
$$
2, 3, 4, \ldots, d-1, d, d, d, d, d, d, d, \ldots
$$
\end{enumerate}
\end{exercise}

\begin{remark}
\label{remark-CM-dim-1-embedding-dim-2}
Cela donne à penser qu'un anneau local noethérien de Cohen-Macaulay de dimension 1
et de dimension de plongement 2 est de la forme $B/FB$, où $B$ est un anneau local
régulier de dimension 2. C'est à peu près exact (sous des hypothèses
convenables de ``bonne tenue'' de l'anneau).
\end{remark}






\section{Une infinité de nombres premiers}
\label{section-many-primes}

\noindent
Section regroupant des questions singulières sur les anneaux dans lesquels
les nombres premiers non inversibles sont en nombre infini.

\begin{exercise}
\label{exercise-not-in-Q}
Donner un exemple d'une ${\mathbf Z}$-algèbre $R$ de type fini
possédant les deux propriétés suivantes :
\begin{enumerate}
\item Il n'existe aucun homomorphisme d'anneaux $R \to {\mathbf Q}$.
\item Pour tout nombre premier $p$, il existe un idéal maximal
${\mathfrak m} \subset R$ tel que $R/{\mathfrak m} \cong {\mathbf F}_p$.
\end{enumerate}
\end{exercise}

\begin{exercise}
\label{exercise-strange-fp-1}
Pour $f \in {\mathbf Z}[x, u]$, nous définissons $f_p(x)
= f(x, x^p) \bmod p \in {\mathbf F}_p[x]$. Donner un exemple
d'un $f \in {\mathbf Z}[x, u]$ tel que les deux propriétés
suivantes tiennent :
\begin{enumerate}
\item Il existe infiniment beaucoup de $p$ de sorte que
$f_p$ n'a pas de zéro dans ${\mathbf F}_p$.
\item Pour tous $p >> 0$, le facteur polynôme
$f_p$ est linéaire ou quadratique.
\end{enumerate}
\end{exercise}

\begin{exercise}
\label{exercise-strange-fp-2}
Pour $f \in {\mathbf Z}[x, y, u, v]$, nous définissons
$f_p(x, y) = f(x, y, x^p, y^p) \bmod p \in {\mathbf F}_p[x, y]$. Donner un
exemple de ``interesting'' d'un $f$ tel que $f_p$ est réductible pour
tous $p >> 0$. Par exemple, $f = xv-yu$ avec $f_p = xy^p-x^py =
xy(x^{p-1}-y^{p-1})$ est ``uninteresting''; n'importe quel $f$ dépend seulement de $x, u$ est ``uninteresting'',
etc.
\end{exercise}

\begin{remark}
\label{remark-strange-fp}
Soit $h \in {\mathbf Z}[y]$ être un polynôme monique de degré $d$.
Ensuite:
\begin{enumerate}
\item La carte $A = {\mathbf Z}[x] \to B ={\mathbf Z}[y]$,
$x \mapsto h$ est finie localement sans rang $d$.
\item Pour toutes les premières $p$
la carte $A_p = {\mathbf F}_p[x]\to B_p = {\mathbf F}_p[y]$,
$y \mapsto h(y) \bmod p$ est finie localement sans rang $d$.
\end{enumerate}
\end{remark}

\begin{exercise}
\label{exercise-strange-fp-3}
Laisse-moi. $h, A, B, A_p, B_p$ être comme dans la remarque. $f \in {\mathbf Z}[x, u]$
nous définissons $f_p(x) = f(x, x^p) \bmod p \in {\mathbf F}_p[x]$. Pour
$g \in {\mathbf Z}[y, v]$ nous définissons
$g_p(y) = g(y, y^p) \bmod p \in {\mathbf F}_p[y]$C'est vrai.
\begin{enumerate}
\item Donner un exemple de $h$ et $g$ de
sorte qu'il n'existe pas de $f$ avec la propriété
$$
f_p  =  Norm_{B_p/A_p}(g_p).
$$
\item Afficher que pour tout choix de $h$ et $g$ comme ci-dessus,
il existe un $f$ non nul tel que pour tous $p$ nous avons
$$
Norm_{B_p/A_p}(g_p)\quad\text{divides}\quad f_p .
$$
Si vous voulez vous limiter au cas $h = y^n$, même avec $n = 2$,
mais c'est vrai en général.
\item Discutez de la pertinence de cela pour les exercices 6 et 7 de l'ensemble
précédent.
\end{enumerate}
\end{exercise}

\begin{exercise}
\label{exercise-strange-fp-unsolved}
Les problèmes non résolus peuvent être très difficiles
ou faciles.
\begin{enumerate}
\item Y a-t-il un $f \in {\mathbf Z}[x, u]$ tel que $f_p$ est irréductible
pour un nombre infini de $p$? (Aperçu: Oui, cela arrive pour $f(x, u) = u -
x - 1$ et aussi pour $f(x, u) = u^2 - x^2 + 1$.)
\item Laisser $f \in {\mathbf Z}[x, u]$ non zéro, et supposer $\deg_x(f_p)
= dp + d'$ pour tous les grands $p$. (En d'autres termes $\deg_u(f) = d$ et
le coefficient $c$ de $u^d$ dans $f$ a $\deg_x(c) = d'$.) Supposons que nous puissions
écrire $d = d_1 + d_2$ et $d' = d'_1 + d'_2$ avec $d_1,
d_2 > 0$ et $d'_1, d'_2 \geq 0$ de sorte que pour tous les
grands $p$ il existe une factorisation
$$
f_p = f_{1, p} f_{2, p}
$$
avec $\deg_x(f_{1, p}) = d_1p + d'_1$. Est-il vrai que $f$ vient par
une construction standard comme dans l'exercice 4? (Plus précisément, y a-t-il un $h$
et $g$ tel que $Norm_{B_p/A_p}(g_p)$ divise $f_p$ pour tous $p >> 0$.)
\item Question analogue à celle de (b) mais maintenant avec
$f \in {\mathbf Z}[x_1, x_2, u_1, u_2]$ irréductible et en supposant simplement que $f_p(x_1,
x_2) = f(x_1, x_2, x_1^p, x_2^p) \bmod p$ facteurs pour tous $p
>> 0$.
\end{enumerate}
\end{exercise}


















\section{Catégorie dérivée finie}
\label{section-filtered-derived}

\noindent
Pour faire les exercices dans cette section, veuillez lire le
matériel dans Homology, Section \ref{homology-section-filtrations}. Nous dirons $A$ est un
objet filtré de $\mathcal{A}$, pour signifier que $A$ vient doté d'une
filtration $F$ que nous omettons de la notation.

\begin{exercise}
\label{exercise-split-injective}
Soit $\mathcal{A}$ être une catégorie abelienne.
Soit $I$ être un objet filtré de
$\mathcal{A}$. Supposons que la filtration sur
$I$ est finie et que chaque $\text{gr}^p(I)$ est un objet
injecteur de $\mathcal{A}$. Montrer qu'il
existe un isomorphisme $I \cong \bigoplus \text{gr}^p(I)$ avec
la filtration $F^p(I)$ correspondant à $\bigoplus_{p' \geq p} \text{gr}^p(I)$.
\end{exercise}

\begin{exercise}
\label{exercise-filtered-injective}
Soit $\mathcal{A}$ être une catégorie abelienne. Soit
$I$ être un objet filtré de $\mathcal{A}$. Supposez
que la filtration sur $I$ est finie.
Montrer ce qui suit sont équivalents:
\begin{enumerate}
\item Pour tout diagramme solide
$$
\xymatrix{
A \ar[r]_\alpha \ar[d] & B \ar@{-->}[ld] \\
I &
}
$$
des objets filtrés
avec (\romannumeral1) les filtrations sur $A$ et $B$ sont
finies, et (\romannumeral2) $\text{gr}(\alpha)$ injectable la flèche
pointillée existe faisant le déplacement du diagramme.
\item Chaque $\text{gr}^p I$ est injectable.
\end{enumerate}
\end{exercise}

\noindent
Notez que, compte tenu d'un morphisme $\alpha : A \to
B$ d'objets filtrés avec des filtrations finies pour dire que
$\text{gr}(\alpha)$ injectable est la même chose que dire que $\alpha$ est
un {\it monomorphisme strict} dans la catégorie
$\text{Fil}(\mathcal{A})$. À savoir, être un monomorphisme signifie $\Ker(\alpha) = 0$
et signifie strict que cela implique également
$\Ker(\text{gr}(\alpha)) = 0$. Voir Homologie, Lemma \ref{homology-lemma-characterize-strict}. (Nous utilisons seulement
le terme «injectif» pour un morphisme dans une catégorie abelienne, bien
qu'il ait un sens dans toute catégorie additive ayant des
grains.) Les exercices ci-dessus justifient la définition suivante.

\begin{definition}
\label{definition-injective-filtered}
Soit $\mathcal{A}$ être une catégorie abelienne. Soit
$I$ être un objet filtré de $\mathcal{A}$. Supposons
que la filtration sur $I$ est finie.
Nous disons que $I$ est {\it par injection filtrée} si chaque $\text{gr}^p(I)$
est un objet injectable de $\mathcal{A}$.
\end{definition}

\noindent
Nous faisons la définition suivante pour éviter d'avoir à
dire "avec une filtration finie" partout.

\begin{definition}
\label{definition-finite-filtration-category}
Que $\mathcal{A}$ soit une catégorie
abelienne. Nous indiquons {\it $\text{Fil}^f(\mathcal{A})$} la sous-catégorie complète
de $\text{Fil}(\mathcal{A})$ dont les objets sont constitués
de $A \in \Ob(\text{Fil}(\mathcal{A}))$ dont la
filtration est finie.
\end{definition}

\begin{exercise}
\label{exercise-inject-into-injective}
Laisser $\mathcal{A}$ être une catégorie abelienne.
Assumer $\mathcal{A}$ a assez d'injectifs. Laisser
$A$ être un objet de $\text{Fil}^f(\mathcal{A})$. Montrer
qu'il existe un monomorphisme strict $\alpha : A \to
I$ de $A$ dans un objet injectable filtré $I$ de $\text{Fil}^f(\mathcal{A})$.
\end{exercise}

\begin{definition}
\label{definition-filtered-quasi-isomorphism}
Laisse-moi. $\mathcal{A}$ être une catégorie
abélienne. $\alpha : K^\bullet \to L^\bullet$ être un
morphisme de complexes de $\text{Fil}(\mathcal{A})$. Nous
disons que $\alpha$ est un {\it quasi-isomorphisme
filtré} si pour chaque $p
\in \mathbf{Z}$ le morphisme $\text{gr}^p(K^\bullet) \to \text{gr}^p(L^\bullet)$
est un quasi-isomorphisme.
\end{definition}

\begin{definition}
\label{definition-filtered-acyclic}
Soit $\mathcal{A}$ être une catégorie abelienne.
Soit $K^\bullet$ être un complexe de $\text{Fil}^f(\mathcal{A})$. Nous
disons que $K^\bullet$ est {\it acyclique filtré}
si pour chaque $p \in \mathbf{Z}$ le complexe $\text{gr}^p(K^\bullet)$ est
acyclique.
\end{definition}

\begin{exercise}
\label{exercise-filtered-quasi-isomorphism}
Soit $\mathcal{A}$ être une catégorie abelienne.
Soit $\alpha : K^\bullet \to L^\bullet$ être un morphisme de complexes
délimités en dessous de $\text{Fil}^f(\mathcal{A})$. (Notez le supercrite $f$.) Montrer que
les éléments suivants sont équivalents :
\begin{enumerate}
\item $\alpha$ est un quasi-isomorphisme filtré,
\item pour chaque $p \in \mathbf{Z}$
la carte $\alpha : F^pK^\bullet \to F^pL^\bullet$ est un quasi-isomorphisme,
\item pour chaque $p \in
\mathbf{Z}$ la carte $\alpha : K^\bullet/F^pK^\bullet \to L^\bullet/F^pL^\bullet$
est un quasi-isomorphisme, et
\item le cône de
$\alpha$ (voir Catégories dérivées, Définition \ref{derived-definition-cone}) est
un complexe acyclique filtré.
\end{enumerate}
De plus, montrez que si $\alpha$ est un quasi-isomorphisme filtré,
alors $\alpha$ est aussi un quasi-isomorphisme habituel.
\end{exercise}

\begin{exercise}
\label{exercise-injective-resolution}
Laisser $\mathcal{A}$ être une catégorie abelienne.
Assumer $\mathcal{A}$ a assez d'injectifs.
Laisser $A$ être un objet de $\text{Fil}^f(\mathcal{A})$.
Montrer qu'il existe un
complexe $I^\bullet$ de $\text{Fil}^f(\mathcal{A})$, et un
morphisme $A[0] \to I^\bullet$ tel que
\begin{enumerate}
\item chaque $I^p$ est filtré par injection,
\item $I^p = 0$ pour $p < 0$, et
\item $A[0] \to I^\bullet$ est un quasi-isomorphe filtré.
\end{enumerate}
\end{exercise}

\begin{exercise}
\label{exercise-injective-resolution-complex}
Let $\mathcal{A}$ be abelian category. Assumer
$\mathcal{A}$ a assez d'injectifs. Let $K^\bullet$
be a lounded under complex of objects of
$\text{Fil}^f(\mathcal{A})$. Montrer qu'il existe un quasi-isomorphisme filtré
$\alpha : K^\bullet \to I^\bullet$ avec $I^\bullet$ un complexe
de $\text{Fil}^f(\mathcal{A})$ ayant filtré des termes injectables $I^n$,
et limité en dessous. En fait, nous pouvons
choisir $\alpha$ de telle sorte que chaque $\alpha^n$ est
un monomorphisme strict.
\end{exercise}

\begin{exercise}
\label{exercise-morphisms-lift}
Soit $\mathcal{A}$ être une catégorie abelienne.
Considérez un diagramme solide
$$
\xymatrix{
K^\bullet \ar[r]_\alpha \ar[d]_\gamma & L^\bullet \ar@{-->}[dl]^\beta \\
I^\bullet
}
$$
d'un complexe de $\text{Fil}^f(\mathcal{A})$. Assumer $K^\bullet$, $L^\bullet$
et $I^\bullet$ sont délimités ci-dessous et supposer que
chaque $I^n$ est un objet injectable filtré.
Supposons également que $\alpha$ est un quasi-isomorphisme filtré.
\begin{enumerate}
\item Il existe une carte des complexes $\beta$, ce qui fait que le
diagramme se déplace vers l'homotopie.
\item Si $\alpha$ est un monomorphisme strict à chaque degré, alors
nous pouvons trouver un $\beta$ qui fait le déplacement du diagramme.
\end{enumerate}
\end{exercise}

\begin{exercise}
\label{exercise-acyclic-is-zero}
Soit $\mathcal{A}$ être une catégorie
abelienne. Soit $K^\bullet$, $K^\bullet$ être des complexes de
$\text{Fil}^f(\mathcal{A})$.
\begin{enumerate}
\item $K^\bullet$ délimitée ci-dessous et filtrée acyclique, et
\item $I^\bullet$ délimité ci-dessous et composé d'objets injectables filtrés.
\end{enumerate}
Alors tout morphisme $K^\bullet \to I^\bullet$ est homotopique à zéro.
\end{exercise}

\begin{exercise}
\label{exercise-morphisms-equal-up-to-homotopy}
Soit $\mathcal{A}$ être une catégorie abelienne.
Considérez un diagramme solide
$$
\xymatrix{
K^\bullet \ar[r]_\alpha \ar[d]_\gamma & L^\bullet \ar@{-->}[dl]^{\beta_i} \\
I^\bullet
}
$$
de complexes de $\text{Fil}^f(\mathcal{A})$. Supposons $K^\bullet$,
$L^\bullet$ et $I^\bullet$ délimités ci-dessous et chacun d'eux $I^n$
un objet injectable filtré.
$\alpha$ un quasi-isomorphisme filtré. Deux morphismes $\beta_1, \beta_2$
faire en sorte que le diagramme se déplace
vers l'homotopie sont homotopiques.
\end{exercise}







\section{Fonctions régulières}
\label{section-regular-functions}

\begin{exercise}
\label{exercise-extra-function}
Considérez la courbe d'affine $X$ donnée par l'équation $t^2 =
s^5 + 8$ dans $\mathbf{C}^2$ avec les coordonnées $s, t$. Soit
$x \in X$ être le point avec les coordonnées $(1,
3)$. Soit $U = X \setminus \{x\}$. Prouver qu'il y
a une fonction régulière sur $U$ qui n'est pas la
restriction d'une fonction régulière sur $\mathbf{C}^2$, c'est-à-dire qu'elle n'est pas
la restriction d'un polynôme dans $s$ et $t$ à $U$.
\end{exercise}

\begin{exercise}
\label{exercise-no-extra-function}
Soit $n \geq 2$. Soit $E \subset \mathbf{C}^n$ être un sous-ensemble
fini. Montrer que toute fonction régulière sur $\mathbf{C}^n \setminus E$
est un polynôme.
\end{exercise}

\begin{exercise}
\label{exercise-cone}
Soit $X \subset \mathbf{C}^n$ être une variété d'affines. Disons que $X$
est un cône {\it} si $x = (a_1, \ldots, a_n) \in X$
et $\lambda \in \mathbf{C}$ implique $(\lambda a_1, \ldots, \lambda
a_n) \in X$. Bien sûr, si $\mathfrak p \subset \mathbf{C}[x_1,
\ldots, x_n]$ est un idéal premier généré par des polynômes homogènes dans $x_1,
\ldots, x_n$, alors la variété d'affine $X = V(\mathfrak p) \subset \mathbf{C}^n$
est un cône. Montrer que inversement
l'idéal premier $\mathfrak p \subset \mathbf{C}[x_1, \ldots, x_n]$
correspondant à un cône peut être généré par des polynômes homogènes
dans $x_1, \ldots, x_n$.
\end{exercise}

\begin{exercise}
\label{exercise-extra-function-cone}
Donner un exemple d'une variété d'affines $X \subset \mathbf{C}^n$ qui
est un cône (voir Exercice \ref{exercise-cone}) et
une fonction régulière $f$ sur $U = X \setminus \{(0, \ldots,
0)\}$ qui n'est pas la restriction d'une
fonction polynomiale sur $\mathbf{C}^n$.
\end{exercise}

\begin{exercise}
\label{exercise-regular-functions}
Dans cet exercice, nous essayons de voir ce qui se passe avec les
fonctions régulières sur les champs non-algébriques fermés. Soit $k$ être un
corps. Soit $Z \subset k^n$ être un sous-ensemble Zariski fermé localement, c'est-à-dire qu'il
existe des idéaux $I \subset J \subset k[x_1, \ldots, x_n]$ tels que
$$
Z = \{a \in k^n \mid
f(a) = 0\ \forall\ f \in I,\ \exists\ g \in J,\ g(a) \not = 0\}.
$$
Une fonction $\varphi : Z \to k$ est dite {\it régulière} si pour chaque $z
\in Z$ il existe un quartier ouvert de Zariski $z \in U \subset Z$ et
des polynômes $f, g \in k[x_1, \ldots, x_n]$ tels que $g(u)
\not = 0$ pour tous $u \in U$ et tel
que $\varphi(u) = f(u)/g(u)$ pour tous $u \in U$.
\begin{enumerate}
\item Si $k = \bar k$ et $Z = k^n$ montrent que les fonctions
régulières sont données par des polynômes. (Ce n'est que si vous n'avez pas déjà vu
cet argument.)
\item Si $k$ est fini montre que (a) chaque fonction $\varphi$ est régulière, (b)
l'anneau des fonctions régulières est fini dimensionnel sur $k$. (Si vous aimez, vous
pouvez prendre $Z = k^n$ et même $n = 1$.)
\item Si $k = \mathbf{R}$ donner un exemple d'une fonction régulière sur
$Z = \mathbf{R}$ qui n'est pas donnée par un polynôme.
\item Si $k = \mathbf{Q}_p$ donner un exemple d'une fonction régulière sur
$Z = \mathbf{Q}_p$ qui n'est pas donnée par un polynôme.
\end{enumerate}
\end{exercise}






\section{Feuilles}
\label{section-sheaves}

\noindent
Un morphisme $f : X \to Y$ d'une catégorie $\mathcal{C}$ est un monomorphisme {\it} si pour chaque
paire de morphismes $a, b : W \to X$ nous avons
$f \circ a = f \circ b \Rightarrow a = b$. Un monomorphisme dans la catégorie
des ensembles est une carte injectable des ensembles.

\begin{exercise}
\label{exercise-mono-sheaves-sets}
Prouver soigneusement qu'une carte des gerbes d'ensembles est un monomorphisme (dans
la catégorie des gerbes d'ensembles) si et seulement si les cartes induites sur
toutes les tiges sont injectables.
\end{exercise}

\noindent
Un morphisme $f : X \to Y$ d'une catégorie $\mathcal{C}$ est un isomorphisme {\it} s'il
existe un morphisme $g : Y \to X$ tel que $f \circ g = \text{id}_Y$
et $g \circ f = \text{id}_X$. Un isomorphisme dans la catégorie des ensembles est
une carte bijective des ensembles.

\begin{exercise}
\label{exercise-isomorphism-sheaves-sets}
Prouver soigneusement qu'une carte des gerbes d'ensembles est un isomorphisme (dans
la catégorie des gerbes d'ensembles) si et seulement si les cartes induites sur
toutes les tiges sont bijectives.
\end{exercise}

\noindent
Un morphisme $f : X \to Y$ d'une catégorie $\mathcal{C}$ est un épimorphisme {\it} si pour
chaque paire de morphismes $a, b : Y \to Z$ nous avons
$a \circ f = b \circ f \Rightarrow a = b$. Un épimorphisme dans
la catégorie des ensembles est une carte surjective des ensembles.

\begin{exercise}
\label{exercise-epi-sheaves-sets}
Prouver soigneusement qu'une carte des gerbes d'ensembles est un épimorphisme (dans
la catégorie des gerbes d'ensembles) si et seulement si les cartes induites sur
toutes les tiges sont surjectives.
\end{exercise}

\begin{exercise}
\label{exercise-adjoint-push-pull}
Soit $f : X \to Y$ être une
carte des espaces topologiques. Prouver $f_\ast$ puis $f^{-1}$ pour les gerbes de {\bf
sets} forment une paire de functeurs.
\end{exercise}

\begin{exercise}
\label{exercise-j-shriek}
Soit $j : U \to X$ être une immersion ouverte.
\begin{enumerate}
\item Pullback $j^{-1} : \Sh(X) \to \Sh(U)$ a une adjointe gauche
$j_{!} : \Sh(U) \to \Sh(X)$ appelée {\it extension par le jeu vide}.
\item Caractériser les tiges de $j_{!}({\mathcal G})$ pour
$\mathcal{G} \in \Sh(U)$.
\item Pullback $j^{-1} : \textit{Ab}(X) \to \textit{Ab}(U)$ a un joint gauche
$j_{!} : \textit{Ab}(U) \to \textit{Ab}(X)$ appelé {\it extension par zéro}.
\item Caractériser les tiges de $j_{!}({\mathcal G})$
pour $\mathcal{G} \in \textit{Ab}(U)$.
\end{enumerate}
Observer cette extension de zéro diffère de l'extension par le
jeu vide!
\end{exercise}

\begin{exercise}
\label{exercise-not-locally-generated-by-sections}
Soit $X = \mathbf{R}$ avec la topologie
habituelle. Soit $\mathcal{O}_X = \underline{\mathbf{Z}/2\mathbf{Z}}_X$. Soit $i :
Z = \{0\} \to X$ être l'inclusion
et laissez $\mathcal{O}_Z = \underline{\mathbf{Z}/2\mathbf{Z}}_Z$. Prouver ce qui
suit (les trois premiers suivent les définitions, mais si vous n'êtes
pas clair sur les définitions, vous devriez les élucider):
\begin{enumerate}
\item $i_*\mathcal{O}_Z$ est une gerbe de gratte-ciel.
\item Il y a une
carte de surjectif canonique
de $\underline{\mathbf{Z}/2\mathbf{Z}}_X \to i_*\underline{\mathbf{Z}/2\mathbf{Z}}_Z$.
Notez le noyau $\mathcal{I} \subset \mathcal{O}_X$.
\item $\mathcal{I}$ est une gaine idéale de $\mathcal{O}_X$.
\item La gerbe $\mathcal{I}$ sur $X$ ne peut pas
être générée
localement par des sections (comme dans Modules, Définition \ref{modules-definition-locally-generated}).
\end{enumerate}
\end{exercise}

\begin{exercise}
\label{exercise-quotient-j-shriek-Z}
Soit $X$ être un espace topologique.
Soit ${\mathcal F}$ être une gerbe abelienne sur $X$. Montrer
que ${\mathcal F}$ est le quotient d'une somme directe (éventuellement très grande) des
gerbes dont tous les termes sont de la forme
$$
j_{!}(\underline{{\mathbf Z}}_U)
$$
où $U \subset X$ est ouvert et $\underline{{\mathbf Z}}_U$ indique la masse
constante avec la valeur ${\mathbf Z}$ sur $U$.
\end{exercise}

\begin{remark}
\label{remark-direct-sum-stalk-abelian}
Que $X$ soit un espace
topologique. Dans la catégorie des gerbes abeliennes, la somme
directe d'une famille de gerbes $\{{\mathcal F}_i\}_{i\in I}$ est la gerbe associée au pré
Sheaf $U \mapsto \oplus {\mathcal F}_i(U)$. Par conséquent, la tige de la somme
directe à un point $x$ est la somme directe des tiges du
${\mathcal F}_i$ à $x$.
\end{remark}

\begin{exercise}
\label{exercise-product-over-points}
Soit $X$ être un espace topologique. Supposons qu'on nous donne une collection de
groupes abeliens $A_x$ indexés par $x \in X$. Montrer que la
règle $U \mapsto \prod_{x \in U} A_x$ avec des mappages de restriction
évidents définit une gerbe $\mathcal{G}$ de groupes abeliens. Montrer, par un exemple, que
ce n'est généralement pas le cas pour $\mathcal{G}_x = A_x$ pour $x \in X$.
\end{exercise}

\begin{exercise}
\label{exercise-modified-product-over-points}
Soit $X$, $A_x$, $\mathcal{G}$ être comme
dans l'exercice \ref{exercise-product-over-points}. Soit $\mathcal{B}$ être une
base pour la topologie de $X$, voir Topologie, Définition
\ref{topology-definition-base}. Pour $U \in \mathcal{B}$ laissez $A_U$ être
un sous-groupe $A_U \subset \mathcal{G}(U) = \prod_{x
\in U} A_x$. Supposons que pour $U \subset V$ avec
$U, V \in \mathcal{B}$ les cartes de restriction $A_V$
dans $A_U$. Pour $U \subset X$ set ouvert
$$
\mathcal{F}(U) =
\left\{
(s_x)_{x \in U}
\middle|
\begin{matrix}
\text{ for every }x\text{ in }U\text{ there exists } V \in \mathcal{B} \\
x \in V \subset U\text{ such that } (s_y)_{y \in V} \in A_V
\end{matrix}
\right\}
$$
Afficher que $\mathcal{F}$ définit une gerbe de groupes abeliens sur $X$. Montrer,
par un exemple, que ce n'est généralement pas le cas
pour $\mathcal{F}(U) = A_U$ pour $U \in \mathcal{B}$.
\end{exercise}

\begin{exercise}
\label{exercise-exact-but-not-a-stalk-functor}
Donner un exemple d'espace topologique $X$ et d'un functeur
$$
F : \Sh(X) \longrightarrow \textit{Sets}
$$
ce qui est exact (commutés avec des produits finis et des égaliseurs
et commutations avec des coproduits finis et des coégaliseurs, voir
Catégories, Section \ref{categories-section-exact-functor}), mais il n'y a pas de point $x \in
X$ tel que $F$ est isomorphe au functeur
de la tige $\mathcal{F} \mapsto \mathcal{F}_x$.
\end{exercise}




\section{Schemes}
\label{section-schemes}

\noindent
Laisse-moi. $LRS$ être la catégorie des espaces sonnés locaux. Un
schéma d'affine est un objet dans $LRS$ isomorphique dans $LRS$ à
une paire du formulaire $(\Spec(A), {\mathcal O}_{\Spec(A)})$. Un schéma est un
objet $(X,
{\mathcal O}_X)$ des $LRS$ de telle sorte que chaque point $x\in X$ a
un quartier ouvert $U \subset X$ de telle sorte que la
paire $(U, {\mathcal O}_X|_U)$ C'est un régime d'affinité.

\begin{exercise}
\label{exercise-one-point}
Trouver un espace local $1$-point qui n'est pas un schéma.
\end{exercise}

\begin{exercise}
\label{exercise-two-points}
Supposez que $X$ est un schéma dont l'espace
topologique sous-jacent a 2 points. Montrer que $X$ est un schéma affine.
\end{exercise}

\begin{exercise}
\label{exercise-discrete-finite-set-points}
Supposons que $X$ soit un schéma dont l'espace topologique sous-jacent est un
ensemble discret fini. Montrer que $X$ est un schéma affine.
\end{exercise}

\begin{exercise}
\label{exercise-three-points}
Montrer qu'il existe un régime non-affine ayant trois points.
\end{exercise}

\begin{exercise}
\label{exercise-quasi-compact-closed-point}
Supposons que $X$ soit un schéma quasi-compact non vide. Montrer
que $X$ a un point fermé.
\end{exercise}

\begin{remark}
\label{remark-open-immersion}
Lorsque $(X, {\mathcal O}_X)$ est un espace annelé et $U \subset
X$ est un sous-ensemble ouvert, alors $(U, {\mathcal O}_X|_U)$ est un espace annelé.
Notation : ${\mathcal O}_U = {\mathcal O}_X|_U$. Il existe un morphisme canonique
des espaces annelés
$$
j : (U, {\mathcal O}_U) \longrightarrow (X, {\mathcal O}_X).
$$
Si $(X, {\mathcal O}_X)$ est un espace sonné localement, de
même que $(U, {\mathcal O}_U)$
et $j$ est un morphisme d'espaces annelés locaux. $(X, {\mathcal O}_X)$ est
un régime,
il en est de même pour $(U, {\mathcal O}_U)$ et $j$ C'est un
morphisme
des schémas. $(U, {\mathcal O}_U)$ est un {\it sous-régime ouvert} des $(X, {\mathcal
O}_X)$
et que $j$ est un {\it immersion ouverte}. Plus généralement, tout
morphisme $j' : (V, {\mathcal O}_V) \to (X, {\mathcal O}_X)$ c'est-à-dire {\it isomorphe} à
un
morphisme $j : (U, {\mathcal O}_U) \to (X, {\mathcal O}_X)$ comme ci-dessus est
appelé
une immersion ouverte.
\end{remark}

\begin{exercise}
\label{exercise-open-affine-not-affine}
Donner un exemple d'un schéma d'affine $(X, {\mathcal O}_X)$ et d'un
schéma ouvert $U \subset X$ tel que $(U, {\mathcal O}_X|U)$ n'est pas un schéma
d'affine.
\end{exercise}

\begin{exercise}
\label{exercise-morphism-does-not-extend}
À titre d'exemple d'une paire de schémas
d'affines $(X, {\mathcal O}_X)$, $(Y, {\mathcal O}_Y)$, d'un
sous-régime ouvert $(U, {\mathcal O}_X|_U)$ de $X$
et d'un morphisme de schémas
$(U, {\mathcal O}_X|_U) \to (Y, {\mathcal O}_Y)$ qui
ne s'étend pas à un morphisme de
schémas $(X, {\mathcal O}_X) \to (Y, {\mathcal O}_Y)$.
\end{exercise}

\begin{exercise}
\label{exercise-closed-subscheme-does-not-extend}
(C'est assez dur.) Vu
un exemple d'un schéma $X$, et d'un sous-régime ouvert $U \subset X$ et
d'un sous-régime fermé $Z \subset U$ tel que $Z$ ne s'étend pas
à un sous-régime fermé de $X$.
\end{exercise}

\begin{exercise}
\label{exercise-not-morphism-schemes}
Donner un exemple d'un schéma $X$, d'un corps $K$ et
d'un morphisme des espaces annelés $\Spec(K) \to X$ qui
n'est PAS un morphisme des schémas.
\end{exercise}

\begin{exercise}
\label{exercise-just-kidding}
Faites tous les exercices dans \cite[Chapter II]{H}, Sections
1 et 2...\ \ Juste plaisanter!
\end{exercise}

\begin{definition}
\label{definition-integral}
Un schéma $X$ est appelé {\it intégrale} si $X$ n'est pas
vide et pour chaque affine non vide ouverte $U
\subset X$ l'anneau $\Gamma(U, \mathcal{O}_X) = \mathcal{O}_X(U)$ est un domaine.
\end{definition}

\begin{exercise}
\label{exercise-morphism-integral-schemes-surjective-stalks-not-closed}
Donner un exemple de morphisme de {\it intégrale} schémas
$f : X \to Y$ de sorte que les cartes induites ${\mathcal O}_{Y, f(x)}
\to {\mathcal O}_{X, x}$ sont surjectives pour tous $x\in X$, mais $f$
n'est pas une immersion fermée.
\end{exercise}

\begin{exercise}
\label{exercise-fibre-product-affines-not-affine}
Donner un exemple d'un produit en fibre $X \times_S Y$ tel que $X$ et $Y$ sont
des affine mais $X \times_S Y$ ne l'est pas.
\end{exercise}

\begin{remark}
\label{remark-separated-base-fibre-product-affines-affine}
Il s'avère que cela ne peut pas se produire avec $S$ séparé. Savez-vous pourquoi?
\end{remark}

\begin{exercise}
\label{exercise-not-geometrically-integral}
Donner un exemple d'un schéma
$V$ qui est un schéma 1 dimensionnel intégral de
type fini sur ${\mathbf Q}$
de sorte que $\Spec({\mathbf C}) \times_{\Spec({\mathbf Q})} V$
n'est pas intégré.
\end{exercise}

\begin{exercise}
\label{exercise-not-geometrically-reduced}
Donner un exemple d'un schéma
$V$ qui est un schéma 1 dimensionnel intégral de type fini
sur un corps $k$ de sorte que $\Spec(k') \times_{\Spec(k)} V$ ne
soit pas réduit pour une extension de corps finie $k'/k$.
\end{exercise}

\begin{remark}
\label{remark-affine-dimension}
Si votre schéma est affine alors la dimension est
la même que la dimension Krull de l'anneau sous-jacent. Ainsi, vous pouvez
utiliser les résultats des derniers semestres pour calculer la dimension.
\end{remark}






\section{Morphismes}
\label{section-morphisms}

\noindent
Une question importante est, étant donné un morphisme $\pi : X
\to S$, si le morphisme a une section ou une
section rationnelle. Voici quelques exemples d'exercices.

\begin{exercise}
\label{exercise-no-section}
Considérons le morphisme des schémas
$$
\pi :
X = \Spec(\mathbf{C}[x, t, 1/xt])
\longrightarrow
S = \Spec(\mathbf{C}[t]).
$$
\begin{enumerate}
\item Montrer qu'il n'existe pas de morphisme $\sigma : S \to
X$ tel que $\pi \circ \sigma = \text{id}_S$.
\item Montrer qu'il existe un non vide ouvert $U \subset S$ et
un morphisme $\sigma : U \to X$ tel que $\pi \circ \sigma = \text{id}_U$.
\end{enumerate}
\end{exercise}

\begin{exercise}
\label{exercise-no-rational-section}
Considérons le morphisme des schémas
$$
\pi :
X = \Spec(\mathbf{C}[x, t]/(x^2 + t))
\longrightarrow
S = \Spec(\mathbf{C}[t]).
$$
Afficher qu'il n'existe pas de non vide ouvert $U \subset S$ et
de morphisme $\sigma : U \to X$ tel que $\pi \circ \sigma = \text{id}_U$.
\end{exercise}

\begin{exercise}
\label{exercise-has-rational-section}
Soit $A, B, C \in \mathbf{C}[t]$ être des polynômes non
nuls. Considérez le morphisme des schémas
$$
\pi :
X = \Spec(\mathbf{C}[x, y, t]/(A + Bx^2 + Cy^2))
\longrightarrow
S = \Spec(\mathbf{C}[t]).
$$
Montrer qu'il existe un non vide ouvert $U \subset S$ et
un morphisme $\sigma : U \to X$ tel que $\pi \circ \sigma = \text{id}_U$.
(Aperçu: Symboliquement, écrivez $x = X/Z$, $y = Y/Z$ pour certains
$X, Y, Z \in \mathbf{C}[t]$ de degré $\leq d$ pour certains $d$,
et calculez la condition que cela résout l'équation. Ensuite, montrez, en
utilisant la théorie des dimensions, que si $d >> 0$ vous pouvez
trouver non zéro $X, Y, Z$ résolvant l'équation.)
\end{exercise}

\begin{remark}
\label{remark-tsen}
L'exercice \ref{exercise-has-rational-section} est un cas particulier de
« théorème de Tsen ». L'exercice
\ref{exercise-no-section-curve} montre que la méthode est limitée aux équations
à faible degré (coniques lorsque la base et la fibre
ont la dimension 1).
\end{remark}

\begin{exercise}
\label{exercise-no-section-curve}
Considérons le morphisme des schémas
$$
\pi :
X = \Spec(\mathbf{C}[x, y, t]
/(1 + t x^3 + t^2 y^3))
\longrightarrow
S = \Spec(\mathbf{C}[t])
$$
Afficher qu'il n'existe pas de non vide ouvert $U \subset S$ et
de morphisme $\sigma : U \to X$ tel que $\pi \circ \sigma = \text{id}_U$.
\end{exercise}

\begin{exercise}
\label{exercise-no-section-surface}
Considérons les schémas
$$
X = \Spec(\mathbf{C}[\{x_i\}_{i = 1}^{8}, s, t]
/(1 + s x_1^3 + s^2 x_2^3 +
t x_3^3 + st x_4^3 + s^2t x_5^3 +
t^2 x_6^3 + st^2 x_7^3 + s^2t^2 x_8^3))
$$
et
$$
S = \Spec(\mathbf{C}[s, t])
$$
et le morphisme des schémas
$$
\pi : X \longrightarrow S
$$
Afficher qu'il n'existe pas de non vide ouvert $U \subset S$ et
de morphisme $\sigma : U \to X$ tel que $\pi \circ \sigma = \text{id}_U$.
\end{exercise}

\begin{exercise}
\label{exercise-for-number-theorists}
(Pour les théoriciens du nombre.)
$$
Z \subset \Spec\left({\mathbf Z}[x, \frac{1 }{ x(x-1)(2x-1)}]\right)
$$
si bien que le morphisme $Z \to \Spec({\mathbf Z})$ est
fini et surjectif.
\end{exercise}

\begin{exercise}
\label{exercise-quasi-section}
Si vous n'aimez pas la théorie des nombres, vous pouvez essayer
la variante où vous regardez
$$
\Spec\left({\mathbf F}_p[t, x, \frac{1 }{ x(x-t)(tx-1)}]\right)
\longrightarrow
\Spec({\mathbf F}_p[t])
$$
et vous essayez de trouver un sous-régime fermé du schéma
supérieur qui cartographie par surjectivement le schéma inférieur. (Il y a
une raison théorique d'avoir un corps de terrain fini ici; bien
qu'il ne soit pas nécessaire dans ce cas particulier.)
\end{exercise}

\begin{remark}
\label{remark-interpretation-skolem-noether}
Dans les exercices \ref{exercise-for-number-theorists} et \ref{exercise-quasi-section} on nous donne
un morphisme $f : X \to S$ où $S$ est le
spectre d'un anneau Dedekind $A$ avec $f$ plat et surjectif avec
fibre générique géométriquement irréductible. Dans les deux cas, le résultat est qu'il
existe un morphisme surjectif fini $S' \to S$ de telle sorte que le
changement de base $X_{S'} \to S'$ Il s'avère que cela
tient si $A$ est excellent, ses champs de résidus à des
idéaux maximaux sont des extensions algébriques de champs finis, et
$\Pic(A')$ est une torsion lorsque $A'$ est la fermeture intégrale
de $A$ dans une extension finie de son corps de fraction, voir
\cite{MB1}. Par exemple, si $A = \mathbf{Z}$ ou $A = \mathbf{F}_p[x]$ ou
une extension finie de ceux-ci. Cependant,
il s'avère qu'il existe un anneau Dedekind $A$ avec des champs de
résidus finis à des primeurs maximaux dont le groupe Picard a un élément
de nontorsion, voir \cite{Goldman}, et le résultat est faux pour
le spectre d'un tel anneau. \ref{exercise-no-quasi-section} donne un exemple géométrique où
il ne fonctionne pas.
\end{remark}

\begin{exercise}
\label{exercise-no-quasi-section}
Prouver qu'il existe un $f \in \mathbf{C}[x, t]$ qui n'est pas divisible
par $t - \alpha$ pour n'importe quel $\alpha \in \mathbf{C}$
tel qu'il n'existe aucun $Z \subset \Spec(\mathbf{C}[x, t, 1/f])$ qui cartographie de
façon subjective le $\Spec(\mathbf{C}[t])$ fini. (Je pense que $f(x, t)
= (xt - 2)(x - t + 3)$ fonctionne. montrer que n'importe quel
candidat a la propriété requise n'est pas si facile.)
\end{exercise}

\begin{exercise}
\label{exercise-finite}
Laisse-moi. $A \to B$ être une homomorphisme d'anneaux de
type fini. Supposons que $\Spec(B) \to \Spec(A)$ facteurs à
travers une immersion fermée $\Spec(B) \to \mathbf{P}^n_A$ pour certains
$n$- Prouver-le. $A \to B$ est une homomorphisme d'anneaux
finie, c'est-à-dire que $B$ est fini
comme un $A$-Module. $A$ est noethérien (s'il vous plaît
supposez ceci) vous pouvez argumenter en utilisant que $H^i(Z, \mathcal{O}_Z)$ pour $i \in
\mathbf{Z}$ est un fini $A$-module pour chaque sous-régime
fermé $Z \subset \mathbf{P}^n_A$C'est vrai.
\end{exercise}

\begin{exercise}
\label{exercise-projective-finite}
Laisser $k$ être un corps algébrique fermé.
Laisser $f : X \to Y$ être un morphisme
de variétés projectives telles que $f^{-1}(\{y\})$ est fini
pour chaque point fermé $y \in
Y$. Prouver que $f$ est fini comme un morphisme
de schémas. Astuces : a) être fini est
une propriété locale, b) essayer de réduire à l'exercice \ref{exercise-finite},
et c) utiliser une immersion fermée $X \to \mathbf{P}^n_k$ pour
obtenir une immersion fermée $X \to \mathbf{P}^n_Y$ sur $Y$.
\end{exercise}






\section{Espaces mandataires}
\label{section-tangent-space}

\begin{definition}
\label{definition-dual-numbers}
Pour n'importe quel anneau $R$, nous indiquons $R[\epsilon]$ l'anneau
de {\it numéros doubles}. En tant que module $R$, il est
libre avec base $1$, $\epsilon$. La structure de l'anneau provient du réglage
$\epsilon^2 = 0$.
\end{definition}

\begin{exercise}
\label{exercise-tangent-space-Zariski}
Soit $f : X \to S$ être un morphisme des
schémas. Soit $x \in X$ être un point, laissez
$s = f(x)$. Considérez le diagramme commutif solide
$$
\xymatrix{
\Spec(\kappa(x)) \ar[r] \ar[dr] \ar@/^1pc/[rr] &
\Spec(\kappa(x)[\epsilon]) \ar@{.>}[r] \ar[d]&
X \ar[d] \\
&
\Spec(\kappa(s)) \ar[r] &
S
}
$$
avec la flèche incurvée étant le morphisme canonique de
$\Spec(\kappa(x))$ dans $X$. Si $\kappa(x)
= \kappa(s)$ montre que l'ensemble de flèches pointillées qui
font la conversion du diagramme sont dans une à une correspondance
avec l'ensemble de cartes linéaires
$$
\Hom_{\kappa(x)}(
\frac{\mathfrak m_x}{\mathfrak m_x^2 + \mathfrak m_s\mathcal{O}_{X, x}},
\kappa(x))
$$
En d'autres termes : décrire une
telle bijection. (Cela fonctionne plus généralement si $\kappa(x) \supset \kappa(s)$
est une extension algébrique séparable.)
\end{exercise}

\begin{definition}
\label{definition-tangent-space}
Laisse-moi. $f : X \to S$ être un
morphisme des schémas. $x \in X$. Nous doublons l'ensemble
des flèches pointillées de l'exercice \ref{exercise-tangent-space-Zariski} des {\it
espace tangent de $X$ sur $S$} et nous
l'indiquons $T_{X/S, x}$. Un élément de cet espace est
appelé un {\it vecteur tangent} des $X/S$ à $x$C'est vrai.
\end{definition}

\begin{exercise}
\label{exercise-simple-push-out}
Pour tout corps $K$ prouver que le diagramme
$$
\xymatrix{
\Spec(K) \ar[r] \ar[d] & \Spec(K[\epsilon_1]) \ar[d] \\
\Spec(K[\epsilon_2]) \ar[r] &
\Spec(K[\epsilon_1, \epsilon_2]/(\epsilon_1\epsilon_2))
}
$$
est un diagramme pushout dans la catégorie des schémas.
(Ici $\epsilon_i^2 = 0$ comme précédemment.)
\end{exercise}

\begin{exercise}
\label{exercise-tangent-space-vectors-space}
Laisser $f : X \to S$ être
un morphisme des schémas. Laisser $x \in X$. Définir
l'ajout de vecteurs tangents, en utilisant l'exercice \ref{exercise-simple-push-out}
et un morphisme approprié
$$
\Spec(K[\epsilon])
\longrightarrow
\Spec(K[\epsilon_1, \epsilon_2]/(\epsilon_1\epsilon_2)).
$$
De même, définissez la multiplication scalaire des vecteurs tangents (cela est plus
facile). Montrer que $T_{X/S, x}$ devient un espace $\kappa(x)$-vecteur avec
vos constructions.
\end{exercise}

\begin{exercise}
\label{exercise-compute-TS}
Soit $k$ être un corps. Considérez
le morphisme de la structure $f : X = \mathbf{A}^1_k \to \Spec(k) = S$.
\begin{enumerate}
\item Laisser $x \in X$ être un point fermé. Quelle est la dimension
de $T_{X/S, x}$?
\item Laisser $\eta \in X$ être le point générique. Quelle est la
dimension de $T_{X/S, \eta}$?
\item Considérez maintenant $X$ comme un schéma sur
$\Spec(\mathbf{Z})$. Quelles sont les dimensions de $T_{X/\mathbf{Z}, x}$
et $T_{X/\mathbf{Z}, \eta}$?
\end{enumerate}
\end{exercise}

\begin{remark}
\label{remark-tangent-space-relative}
L'exercice \ref{exercise-compute-TS} explique pourquoi il est nécessaire de considérer l'espace
tangent de $X$ sur $S$ pour obtenir une bonne notion.
\end{remark}

\begin{exercise}
\label{exercise-compute-TS-field}
Considérons le morphisme des schémas
$$
f : X = \Spec(\mathbf{F}_p(t))
\longrightarrow
\Spec(\mathbf{F}_p(t^p)) = S
$$
Calculer l'espace tangent de $X/S$ au point unique de
$X$. N'est-ce pas bizarre? Que pensez-vous qu'il se passe si vous
prenez le morphisme des schémas correspondant à $\mathbf{F}_p[t^p] \to \mathbf{F}_p[t]$?
\end{exercise}

\begin{exercise}
\label{exercise-compute-TS-cusp}
Soit $k$ être un
corps. Calculer l'espace tangent de $X/k$ au point $x = (0,
0)$ où $X = \Spec(k[x, y]/(x^2 - y^3))$.
\end{exercise}

\begin{exercise}
\label{exercise-map-tangent-spaces}
Soit $f : X \to Y$ être un morphisme des schémas
sur $S$. Soit $x \in X$ être un
point. Définissez $y = f(x)$. Supposez que la carte naturelle $\kappa(y) \to
\kappa(x)$ est bijective. Montrer, en utilisant la
définition, que $f$ induit une carte linéaire naturelle
$$
\text{d}f : T_{X/S, x} \longrightarrow T_{Y/S, y}.
$$
Correspondez-le à ce qui se passe sur
les anneaux locaux par l'exercice \ref{exercise-tangent-space-Zariski} dans le cas $\kappa(x) = \kappa(s)$.
\end{exercise}

\begin{exercise}
\label{exercise-Jacobian}
Laisser $k$ être un corps algébrique fermé.

\begin{eqnarray*}
f : \mathbf{A}_k^n & \longrightarrow & \mathbf{A}^m_k \\
(x_1, \ldots, x_n) & \longmapsto & (f_1(x_i), \ldots, f_m(x_i))
\end{eqnarray*}
être un morphisme des schémas sur $k$. Ceci
est donné par $m$ polynômes $f_1, \ldots, f_m$ dans
les variables $n$.
$$
A = \left( \frac{\partial f_j}{\partial x_i} \right)
$$
Soit $x \in \mathbf{A}^n_k$ être un point fermé.
Définissez $y = f(x)$. Montrer que la carte sur
les espaces tangents $T_{\mathbf{A}^n_k/k, x} \to T_{\mathbf{A}^m_k/k, y}$ est donnée
par la valeur de la matrice $A$ au point $x$.
\end{exercise}














\section{Feuilles à cosse cohésives }
\label{section-quasi-coherent}

\begin{definition}
\label{definition-quasi-coherent}
Laisser $X$ être
un schéma. Une feuille $\mathcal{F}$ de $\mathcal{O}_X$-modules est {\it quasi-cohérente} si pour
chaque affine ouverte $\Spec(R) = U \subset X$ la restriction
$\mathcal{F}|_U$ est de la forme $\widetilde M$ pour certains $R$-module
$M$.
\end{definition}

\noindent
Il suffit de vérifier ces conditions sur les
membres d'une couverture d'affine
ouverte de $X$. Voir Schemes, Section \ref{schemes-section-quasi-coherent} pour plus
de résultats.

\begin{definition}
\label{definition-specialization}
Soit $X$ être un espace topologique. Soit $x, x' \in
X$. Nous disons $x$ est une spécialisation {\it} de
$x'$ si et seulement si $x \in \overline{\{x'\}}$.
\end{definition}

\begin{exercise}
\label{exercise-quasi-coherent-specialization-points}
Soit $X$ être un schéma. Soit $x, x' \in X$.
Soit $\mathcal{F}$ être une masse quasi-cohérente de $\mathcal{O}_X$-modules. Supposons
que (a) $x$ soit une spécialisation de $x'$ et
(b) $\mathcal{F}_{x'} \not = 0$. Montrer que $\mathcal{F}_x \not = 0$.
\end{exercise}

\begin{exercise}
\label{exercise-O-module-specialization-points}
Trouver un exemple d'un schéma $X$, des points $x, x'
\in X$, une gaine de $\mathcal{O}_X$-modules
$\mathcal{F}$ de sorte que (a) $x$ est une spécialisation de $x'$ et
(b) $\mathcal{F}_{x'} \not = 0$ et $\mathcal{F}_x = 0$.
\end{exercise}

\begin{definition}
\label{definition-Noetherian-scheme}
Un schéma $X$ est appelé {\it localement noethérien} si et seulement
si pour chaque point $x \in X$ il existe un
affine ouvert $\Spec(R) = U \subset X$ tel que
$R$ est noethérien. Un schéma est {\it noethérien} s'il est localement noethérien et quasi-compact.
\end{definition}

\noindent
Si $X$ est localement noethérien alors n'importe quel
affine ouvert de $X$ est le
spectre d'un anneau noethérien, voir Propriétés, Lemma \ref{properties-lemma-locally-noethérien}.

\begin{definition}
\label{definition-coherent}
Let $X$ be a local noethérien
scheme. Let $\mathcal{F}$ be a quasi-coherent fleaf of
$\mathcal{O}_X$-modules. Nous disons que $\mathcal{F}$ est {\it cohérent} si, pour chaque
point $x \in X$, il existe un affine ouvert
$\Spec(R) = U \subset X$ tel que $\mathcal{F}|_U$
est isomorphe à $\widetilde M$ pour certains modules finis $R$-module $M$.
\end{definition}

\begin{exercise}
\label{exercise-extend-quasi-coherent}
Soit $X = \Spec(R)$ être un schéma d'affine.
\begin{enumerate}
\item Laisser $f \in R$.
Laisser $\mathcal{G}$ être une gaine quasi-cohérente de
$\mathcal{O}_{D(f)}$-modules sur le sous-régime
ouvert $D(f)$. Montrer que $\mathcal{G} =
\mathcal{F}|_{D(f)}$ pour certaines gaines quasi-cohérentes de
$\mathcal{O}_X$-modules $\mathcal{F}$.
\item Soit $I \subset R$ être
un idéal. Soit $i : Z \to X$ être
le sous-régime fermé de $X$
correspondant à $I$. Soit $\mathcal{G}$ être une
gaine quasi-cohérente de $\mathcal{O}_Z$-modules sur
le sous-régime fermé de $Z$. Montrer que
$\mathcal{G} = i^*\mathcal{F}$ pour une gaine quasi-cohérente de $\mathcal{O}_X$-modules $\mathcal{F}$.
(Pourquoi est-ce stupide?)
\item Supposons que $R$ est
noethérien. Soit $f \in R$.
Soit $\mathcal{G}$ être une gaine cohérente
de $\mathcal{O}_{D(f)}$-modules sur le
sous-régime ouvert $D(f)$. Montrer que $\mathcal{G} =
\mathcal{F}|_{D(f)}$ pour une gaine cohérente de $\mathcal{O}_X$-modules $\mathcal{F}$.
\end{enumerate}
\end{exercise}

\begin{remark}
\label{remark-extend-off-open}
Si $U \to X$ est une immersion quasi-compacte, alors
n'importe quelle gaine quasi-cohérente sur $U$ est la restriction d'une
gaine quasi-cohérente sur $X$. Si
$X$ est un schéma noethérien, et $U \subset X$ est ouverte,
alors toute gaine cohérente sur $U$ est la restriction d'une
gaine cohérente sur $X$.
Bien sûr, l'exercice ci-dessus est plus facile et ne devrait pas utiliser ces faits généraux.
\end{remark}
















\section{Proj et plans projectifs}
\label{section-proj}

\begin{exercise}
\label{exercise-graded-ring-specified-result}
Donner des exemples d'anneaux gradués $S$ tels que
\begin{enumerate}
\item $\text{Proj}(S)$ est un affine et non vide, et
\item $\text{Proj}(S)$ est intégrale, non vide mais non isomorphe à
${\mathbf P}^n_A$ pour tout $n\geq 0$, tout anneau $A$.
\end{enumerate}
\end{exercise}

\begin{exercise}
\label{exercise-nonconstant-morphism-proj}
Donner un exemple de morphisme non
constant des schémas ${\mathbf P}^1_{\mathbf C} \to {\mathbf P}^5_{\mathbf C}$
sur $\Spec({\mathbf C})$.
\end{exercise}

\begin{exercise}
\label{exercise-isomorphism-P1-conic}
Donner un exemple d'isomorphisme des schémas
$$
{\mathbf P}^1_{\mathbf C} \to
\text{Proj}({\mathbf C}[X_0, X_1, X_2]/(X_0^2 + X_1^2 + X_2^2))
$$
\end{exercise}

\begin{exercise}
\label{exercise-family-special-fibre-different}
Donner un exemple de morphisme des schémas
$f : X \to {\mathbf A}^1_{\mathbf C} = \Spec({\mathbf C}[T])$ de
telle sorte que la fibre (scheme théorétique) $X_t$ des $f$ sur
$t \in {\mathbf A}^1_{\mathbf C}$ est
(a) isomorphe à ${\mathbf P}^1_{\mathbf C}$ lorsque $t$ est un
point fermé qui n'est pas égal à
$0$, et b) ne sont pas isomorphes à ${\mathbf P}^1_{\mathbf C}$
lorsque $t = 0$. Nous appellerons $X_0$ des {\it Fibres spéciales}
Cela peut être fait de plusieurs façons. Essayez de donner des
exemples qui satisfont (chaque de) les contraintes supplémentaires suivantes (sauf si ce
n'est pas possible):
\begin{enumerate}
\item Pouvez-vous le faire avec une fibre projective spéciale?
\item Pouvez-vous le faire avec des fibres spéciales irréductibles et projectives?
\item Pouvez-vous le faire avec une fibre spéciale intégrale et projective?
\item Pouvez-vous le faire avec une fibre spéciale lisse et projective?
\item Pouvez-vous le faire avec $f$ un morphisme plat?
Cela signifie simplement que pour chaque affine ouverte $\Spec(A) \subset X$ la carte de
l'anneau induit $\mathbf{C}[t] \to A$ est plate, ce qui signifie que dans ce cas
tout polynôme non nul dans $t$ est un diviseur non nul sur $A$.
\item Pouvez-vous le faire avec $f$ un morphisme plat et projectif?
\item Pouvez-vous le faire avec $f$ fibre plate, projective et spéciale réduite?
\item Pouvez-vous le faire avec $f$ fibre plate, projective et spéciale irréductible?
\item Pouvez-vous le faire avec $f$ fibre intégrale plate, projective et spéciale?
\end{enumerate}
Que pensez-vous se passe-t-il lorsque vous remplacez
${\mathbf P}^1_{\mathbf C}$ par une autre variété sur ${\mathbf C}$?
(Cela peut devenir très difficile selon les variantes ci-dessus que vous
demandez.)
\end{exercise}

\begin{exercise}
\label{exercise-affine-onto-projective-space}
Soit $n \geq 1$ être n'importe quel
entier positif. Donner un exemple de morphisme
surjectif $X \to {\mathbf P}^n_{\mathbf C}$ avec $X$ affine.
\end{exercise}

\begin{exercise}
\label{exercise-morphism-proj}
Cartes de $\text{Proj}$. Soit $R$ et $S$ être notées anneaux. Supposons que
nous ayons une homomorphisme d'anneaux
$$
\psi : R \to S
$$
et un entier $e \geq 1$ tel que $\psi(R_d) \subset S_{de}$ pour tous
$d \geq 0$. (Par nos conventions, ce n'est pas un homomorphisme des anneaux
gradués, à moins que $e = 1$.)
\begin{enumerate}
\item Pour quels éléments $\mathfrak p \in \text{Proj}(S)$ y
a-t-il un point correspondant bien défini dans $\text{Proj}(R)$? Autrement dit, trouvez
un $U \subset \text{Proj}(S)$ ouvert approprié tel que $\psi$ définit une
carte continue $r_\psi : U \to \text{Proj}(R)$.
\item Donner un exemple où $U \not = \text{Proj}(S)$.
\item Donner un exemple où $U = \text{Proj}(S)$.
\item (Ne l'inscrivez pas.) Convainquez-vous que la carte continue
$U \to \text{Proj}(R)$ vient canoniquement avec une carte sur les
gerbes de sorte que $r_\psi$ est un morphisme des schémas:
$$
\text{Proj}(S) \supset U \longrightarrow \text{Proj}(R).
$$
\item Que pouvez-vous dire de cette carte
si $R = \bigoplus_{d \geq 0} S_{de}$ (en tant qu'anneau classé
avec $S_e$, $S_{2e}$, etc. en degré $1$, $2$, etc.) et
$\psi$ est la cartographie d'inclusion?
\end{enumerate}
\end{exercise}

\noindent
{\bf Notation.} Laisser $R$ être un anneau classé
comme ci-dessus et laisser $n \geq 0$ être un entier. $X =
\text{Proj}(R)$. Ensuite, il y a un module quasi-cohérent unique ${\mathcal O}_X$ ${\mathcal
O}_X(n)$ sur $X$ tel que, pour chaque élément homogène $f
\in R$ de degré positif, nous avons ${\mathcal O}_X |_{D_{+}(f)}$ est
la gaine quasi-cohérente associée au module $R_{(f)} = (R_f)_0$ $(R_f)_n$ ($ =
$éléments homogènes de degré $n$ dans $R_f = R[1/f]$). Voir
\cite[page 116+]{H}.
$$
{\mathcal O}_X(n_1) \otimes_{{\mathcal O}_X} {\mathcal O}_X(n_2)
\longrightarrow
{\mathcal O}_X(n_1 + n_2)
$$

\begin{exercise}
\label{exercise-pathologies-proj}
Pathologies dans $\text{Proj}$. Donner
des exemples de $R$ comme ci-dessus
\begin{enumerate}
\item ${\mathcal O}_X(1)$ n'est pas un module invertible ${\mathcal O}_X$.
\item ${\mathcal O}_X(1)$ est invertible, mais
la carte naturelle ${\mathcal O}_X(1) \otimes_{{\mathcal O}_X} {\mathcal O}_X(1) \to {\mathcal
O}_X(2)$ n'est PAS un isomorphisme.
\end{enumerate}
\end{exercise}

\begin{exercise}
\label{exercise-finitely-many-points-in-affine}
Soit $S$ être une bague
graduée. Soit $X = \text{Proj}(S)$.
Montrer que tout ensemble fini de points de $X$ est contenu dans un affine
standard ouvert.
\end{exercise}

\begin{exercise}
\label{exercise-prepare-glueing}
Laisse-moi. $S$ Soyez une
bague graduée. $X = \text{Proj}(S)$-
Laisse-moi faire. $Z, Z' \subset X$ être deux
sous-régimes fermés. $\varphi : Z \to Z'$
être un isomorphisme. $Z \cap Z'
= \emptyset$. Montrer ça pour n'importe quel $z
\in Z$ il existe un affine ouvert $U \subset X$ de
telle sorte que $z \in
U$, $\varphi(z) \in U$ et $\varphi(Z \cap U) = Z' \cap U$.
(Aperçu : Utiliser
l'exercice \ref{exercise-finitely-many-points-in-affine} et quelque chose qui s'apparente à Schemes, Lemma \ref{schemes-lemma-standard-open-two-affines}.)
\end{exercise}







\section{Morphismes de la ligne projective}
\label{section-from-P1}

\noindent
Dans cette section, nous étudions les morphismes de $\mathbf{P}^1$
aux schémas projectifs.

\begin{exercise}
\label{exercise-from-generic-point}
Soit $k$ être un corps. Soit $k[t] \subset k(t)$ être
l'inclusion de l'anneau polynôme dans son corps de fraction. Soit $X$
être un schéma de type fini sur $k$.
Montrer que pour tout morphisme
$$
\varphi : \Spec(k(t)) \longrightarrow X
$$
sur $k$, il existe un $f \in k[t]$ non nul et un
morphisme $\psi : \Spec(k[t, 1/f]) \to X$ sur
$k$ tel que $\varphi$ est la composition
$$
\Spec(k(t)) \longrightarrow \Spec(k[t, 1/f]) \longrightarrow X
$$
\end{exercise}

\begin{exercise}
\label{exercise-from-generic-point-Pn}
Soit $k$ être un corps. Soit $k[t] \subset k(t)$
être l'inclusion de l'anneau polynôme dans son corps de fraction.
Montrer que pour tout morphisme
$$
\varphi : \Spec(k(t)) \longrightarrow \mathbf{P}^n_k
$$
sur $k$, il existe un morphisme $\psi : \Spec(k[t]) \to \mathbf{P}^n_k$ sur
$k$ tel que $\varphi$ est la composition
$$
\Spec(k(t)) \longrightarrow \Spec(k[t]) \longrightarrow \mathbf{P}^n_k
$$
Astuce : l'image de $\varphi$ est dans un standard ouvert $D_+(T_i)$ pour
certains $i$ ; alors montrez que vous pouvez ``clairer les dénominateurs''.
\end{exercise}

\begin{exercise}
\label{exercise-from-generic-point-projective}
Soit $k$ être un corps. Soit $k[t] \subset k(t)$
être l'inclusion de l'anneau polynôme dans son corps de fraction. Soit
$X$ être un schéma projectif sur $k$.
Montrer que pour tout morphisme
$$
\varphi : \Spec(k(t)) \longrightarrow X
$$
sur $k$, il existe un morphisme $\psi : \Spec(k[t]) \to X$
sur $k$ tel que $\varphi$ est la composition
$$
\Spec(k(t)) \longrightarrow \Spec(k[t]) \longrightarrow X
$$
Conseil: utiliser Exercice \ref{exercise-from-generic-point-Pn}.
\end{exercise}

\begin{exercise}
\label{exercise-from-generic-point-P1-to-projective}
Soit $k$ être un corps. Soit $X$ être un schéma
projectif sur $k$. Soit $K$ être le corps de fonction de
$\mathbf{P}^1_k$ (voir l'indice ci-dessous).
$$
\varphi : \Spec(K) \longrightarrow X
$$
sur $k$, il existe un morphisme $\psi : \mathbf{P}^1_k \to X$
sur $k$ tel que $\varphi$ est la composition
$$
\Spec(k(t)) \longrightarrow \mathbf{P}^1_k \longrightarrow X
$$
Conseil: utiliser Exercice \ref{exercise-from-generic-point-projective} pour chacun des deux morceaux
de l'affine en couverture ouverte $\mathbf{P}^1_k = D_+(T_0)
\cup D_+(T_1)$, utiliser que $D_+(T_0)$ est le spectre
d'un anneau polynôme et que $K$ est le
corps de fraction de ce anneau polynôme.
\end{exercise}







\section{Morphismes des surfaces aux courbes}
\label{section-from-surfaces-to-curves}

\begin{exercise}
\label{exercise-points-projective-space}
Soit $R$ être un
anneau. Soit $R \to k$ être une carte de $R$ à
un corps. Soit $n
\geq 0$.
$$
\Mor_{\Spec(R)}(\Spec(k), \mathbf{P}^n_R)
=
(k^{n + 1} \setminus \{0\})/k^*
$$
où $k^*$ agit par multiplication scalaire sur $k^{n + 1}$.
À partir de maintenant, nous indiquons $(x_0 :
\ldots : x_n)$ le morphisme $\Spec(k) \to \mathbf{P}^n_k$ correspondant
à la classe d'équivalence de l'élément
$(x_0, \ldots, x_n) \in k^{n + 1} \setminus \{0\}$.
\end{exercise}

\begin{exercise}
\label{exercise-curve-projective-plane}
Soit $k$ être un corps. Soit $Z \subset \mathbf{P}^2_k$
être un sous-régime fermé irréductible et
réduit. Montrer que (a) $Z$ est un point fermé, ou
(b) il existe un $F \in k[X_0, X_1, X_2]$ homogène
irréductible de degré $> 0$ tel que $Z =
V_{+}(F)$, ou (c) $Z = \mathbf{P}^2_k$.
\end{exercise}

\begin{exercise}
\label{exercise-bezout}
Soit $k$ être un corps. Soit $Z_1, Z_2 \subset \mathbf{P}^2_k$ être
des sous-régimes fermés irréductibles de la forme $V_{+}(F)$ pour
certains $F_i \in k[X_0, X_1, X_2]$ homogènes et irréductibles de degré $> 0$.
Montrer que $Z_1 \cap Z_2$ n'est pas
vide. (Enfin, utilisez la théorie des dimensions pour estimer la
dimension de l'anneau local $k[X_0, X_1, X_2]/(F_1, F_2)$ à $0$.)
\end{exercise}

\begin{exercise}
\label{exercise-no-nonconstant-morphism-proj}
Montrer qu'il n'y a pas de morphisme non constant
des schémas $\mathbf{P}^2_{\mathbf{C}} \to \mathbf{P}^1_{\mathbf{C}}$ sur $\Spec(\mathbf{C})$. Voici un
{\it morphisme constant} est celui dont l'image est un
seul point. (Enfin: Si le
morphisme n'est pas constant considérer les fibres sur $0$ et
$\infty$ et argumentent qu'ils doivent se rencontrer pour obtenir une contradiction.)
\end{exercise}

\begin{exercise}
\label{exercise-nonconstant-morphism}
Laisse-moi. $k$ Etre un
corps. Supposons que ce soit un corps. $X \subset \mathbf{P}^3_k$ est un
sous-régime fermé donné par une seule équation homogène $F \in k[X_0, X_1, X_2, X_3]$.
En d'autres termes,
$$
X = \text{Proj}(k[X_0, X_1, X_2, X_3]/(F)) \subset \mathbf{P}^3_k
$$
Comme expliqué dans le cours. Supposons que
$$
F = X_0 G + X_1 H
$$
pour certains polynômes homogènes $G, H \in k[X_0, X_1, X_2, X_3]$ de
degré positif. Montrer que si $X_0, X_1, G, H$ n'ont pas de zéros
communs, alors il existe un morphisme non constant
$$
X \longrightarrow \mathbf{P}^1_k
$$
des schémas sur $\Spec(k)$ qui
sur les points de corps (voir Exercice \ref{exercise-points-projective-space}) ressemble à $(x_0 : x_1
: x_2 : x_3) \mapsto (x_0 : x_1)$ chaque fois que $x_0$
ou $x_1$ n'est pas zéro.
\end{exercise}





\section{Garnes inversées}
\label{section-invertible-sheaves}

\begin{definition}
\label{definition-invertible-sheaf}
Laisse-moi. $X$ être un espace sonné
local. {\it invertible ${\mathcal O}_X$-Module} le $X$ est
une masse de ${\mathcal O}_X$-modules ${\mathcal L}$ de telle sorte que chaque point
ait un quartier ouvert $U \subset X$ de telle sorte que
${\mathcal L}|_U$ est isomorphe à ${\mathcal O}_U$ comme ${\mathcal O}_U$-
On dit ça. ${\mathcal L}$ est trivial s'il est isomorphe
à ${\mathcal O}_X$ en tant que ${\mathcal O}_X$-Module.
\end{definition}

\begin{exercise}
\label{exercise-general-facts-invertible}
Les faits généraux.
\begin{enumerate}
\item Afficher qu'un module inverse ${\mathcal O}_X$ sur
un schéma $X$ est quasi-cohérent.
\item Supposons que $X\to Y$ est un morphisme d'espaces annelés
locaux, et ${\mathcal L}$ un module inversé ${\mathcal O}_Y$.
Montrer que $f^\ast {\mathcal L}$ est un module inversé ${\mathcal O}_X$.
\end{enumerate}
\end{exercise}

\begin{exercise}
\label{exercise-invertible-algebra}
Algèbre.
\begin{enumerate}
\item Montrer qu'un module invertible ${\mathcal O}_X$ sur un schéma
d'affine $\Spec(A)$ correspond à un module $A$ $M$ qui est (i)
fini, (ii) projectif, (iii) localement exempt de rang 1, et
donc (iv) plat, et (v) finiment présenté. (Sentez libre de citer
les choses du cours des derniers semestres; ou des livres d'algèbre.)
\item Supposez que $A$ est un domaine et que $M$
est un module comme dans (a). Montrer que $M$ est isomorphe comme
un $A$-module à un $I \subset A$ idéal tel que $IA_{\mathfrak p}$ est le
principal de chaque ${\mathfrak p}$ primaire.
\end{enumerate}
\end{exercise}

\begin{definition}
\label{definition-invertible-module}
Soit $R$ être un anneau. Un module invertible {\it $M$} est un module $R$
$M$ tel que $\widetilde M$ est une gaine invertible sur le spectre
de $R$. Nous disons $M$ est {\it trivial} si $M \cong R$ comme
un module $R$.
\end{definition}

\noindent
En d'autres termes, $M$ est invertible si et seulement
s'il satisfait à toutes les conditions suivantes :
il est plat, de présentation finie, projectif et localement
exempt de rang 1. (Bien sûr, il suffit qu'il soit
localement exempt de rang 1).

\begin{exercise}
\label{exercise-simple-examples-invertible}
Des exemples simples.
\begin{enumerate}
\item
\label{item-affine-line}
Soit $k$ être un corps. Soit $A
= k[x]$. Montrer que $X = \Spec(A)$ n'a que des
modules invertibles triviaux ${\mathcal O}_X$. En d'autres termes, montrez que chaque
module invertible $A$ est exempt de rang 1.
\item
\label{item-affine-line-with-0-and-1-identified}
Soit $A$ être l'anneau
$$
A = \{ f\in k[x] \mid f(0) = f(1) \}.
$$
Montrer qu'il existe un module non trivial inversé $A$, à
moins que $k = {\mathbf F}_2$. (Considérer $\Spec(A)$ comme identifiant
$0$ et $1$ dans ${\mathbf A}^1_k = \Spec(k[x])$.)
\item
\label{item-affine-line-with-cusp}
Même question que dans (\ref{item-affine-line-with-0-and-1-identified}) pour l'anneau $A =
k[x^2, x^3] \subset k[x]$ (sauf maintenant
$k = {\mathbf F}_2$ fonctionne également).
\end{enumerate}
\end{exercise}

\begin{exercise}
\label{exercise-higher-dimension-invertible}
Des dimensions plus élevées.
\begin{enumerate}
\item Prouver que chaque gerbe invertible sur deux dimensions
de l'espace d'affine est trivial. Plus précisément,
laissez ${\mathbf A}^2_k = \Spec(k[x, y])$ où $k$
est un corps. Montrer que chaque gerbe invertible sur ${\mathbf
A}^2_k$ est trivial. (Enfin, une façon de le faire
est de considérer le module correspondant $M$, de regarder
$M
\otimes_{k[x, y]} k(x)[y]$, puis d'utiliser Exercise \ref{exercise-simple-examples-invertible} (\ref{item-affine-line}) pour trouver un générateur
pour cela; puis vous devez encore penser. Une autre
façon est d'utiliser
Exercice \ref{exercise-invertible-algebra} et d'utiliser ce que nous
savons sur les idéaux de l'anneau
polynôme: les premiers de hauteur on est généré par un polynôme
irréductible; puis vous devez toujours penser.)
\item Prouver que chaque gerbe invertible sur n'importe quel
sous-régime ouvert de deux dimensions de l'espace affine est trivial. Plus précisément, laissez $U
\subset {\mathbf A}^2_k$ être un sous-régime ouvert où $k$ est un corps.
Montrer que chaque gerbe invertible sur $U$ est trivial. Conseil: Montrer que chaque
gerbe invertible sur $U$ s'étend à un sur ${\mathbf A}^2_k$. Pas facile; mais
vous pouvez le trouver dans \cite{H}.
\item Trouver un exemple de gerbe invertible
non triviale sur un cône perforé sur un corps. Plus
précisément, laissez $k$ être un corps et laissez $C = \Spec(k[x, y, z]/(xy-z^2))$. Soit
$U = C \setminus \{ (x, y, z) \}$. Trouver
une gerbe invertible non triviale sur $U$. Conseil: Il peut être plus
facile de calculer le groupe de classes d'isomorphisme des gerbes invertibles sur
$U$ que de simplement en trouver une. Notez que $U$
est couvert par les ouvertures $\Spec(k[x, y,
z, 1/x]/(xy-z^2))$ et $\Spec(k[x, y, z,
1/y]/(xy-z^2))$ qui sont `'faciles'' à traiter.
\end{enumerate}
\end{exercise}

\begin{definition}
\label{definition-picard-group}
Soit $X$ être un espace
sonné localement. Le groupe {\it Picard de $X$}
est l'ensemble $\Pic(X)$ des classes d'isomorphisme des modules inverses
$\mathcal{O}_X$ avec addition donnée par le
produit tensor. Voir Modules, Définition \ref{modules-definition-pic}. Pour un
anneau $R$ nous définissons $\Pic(R) = \Pic(\Spec(R))$.
\end{definition}

\begin{exercise}
\label{exercise-traverso}
Soit $R$ être une bague.
\begin{enumerate}
\item Montrer que si $R$ est un domaine normal de
noethérien, alors $\Pic(R) = \Pic(R[t])$. [Conseil : Il
y a une carte $R[t] \to R$, $t \mapsto 0$ qui
est un inverse gauche à la carte $R \to R[t]$.
Il suffit donc de montrer que tout module invertible $R[t]$ $M$
tel que $M/tM \cong R$ est exempt
de rang $1$. Soit $K$ être le corps fractionnaire de $R$. Choisissez une
banalisation $K[t] \to M \otimes_{R[t]} K[t]$ qui est possible par l'exercice \ref{exercise-simple-examples-invertible} (\ref{item-affine-line}).
Ajustez-le donc avec la banalisation de $M/tM$ ci-dessus.
Montrer qu'il s'agit en fait d'une banalisation de $M$ au-dessus de
$R[t]$ (c'est là que la normalité entre en jeu).]
\item Laisser $k$ être un corps.
Affichez que $\Pic(k[x^2, x^3, t]) \not = \Pic(k[x^2, x^3])$.
\end{enumerate}
\end{exercise}













\section{{\v C}ech Cohomologie}
\label{section-cech-cohomology}

\begin{exercise}
\label{exercise-cech-cohomology}
{\v C}ech cohomology. Ici $k$ est un corps.
\begin{enumerate}
\item Laisser $X$ être un schéma avec une couverture
ouverte ${\mathcal U} : X = U_1 \cup U_2$, avec $U_1
= \Spec(k[x])$, $U_2 =
\Spec(k[y])$ avec $U_1 \cap U_2 = \Spec(k[z, 1/z])$ et avec des immersions
ouvertes $U_1 \cap U_2 \to U_1$ resp.\ $U_1 \cap U_2 \to U_2$
déterminée par $x \mapsto z$ resp.\ $y \mapsto z$ (et je veux
vraiment dire cela). (Nous avons vu dans les conférences qu'un tel $X$ existe;
c'est la ligne d'affine avec zéro doublé.) Calculer ${\check H}^1({\mathcal U}, {\mathcal O})$;
p.
ex.\ donner une base pour cela comme un $k$-vectorspace.
\item Pour chaque élément
de ${\check H}^1({\mathcal U}, {\mathcal O})$
construire une séquence exacte de gerbes de ${\mathcal O}_X$-modules
$$
0 \to {\mathcal O}_X \to E \to {\mathcal O}_X \to 0
$$
telle que la limite $\delta(1) \in {\check H}^1({\mathcal U}, {\mathcal O})$
égale l'élément donné.
(Une partie du problème est de donner un sens à cela. Voir aussi
ci-dessous. Il
est également OK de montrer abstraitement une telle chose doit exister.)
\end{enumerate}
\end{exercise}

\begin{definition}
\label{definition-delta}
(Définition du delta.) Supposons que
$$
0 \to {\mathcal F}_1 \to {\mathcal F}_2 \to {\mathcal F}_3 \to 0
$$
est une courte séquence exacte de gerbes abeliennes sur n'importe quel espace
topologique $X$. La
carte des limites $\delta : H^0(X, {\mathcal F}_3) \to {\check H}^1(X,
{\mathcal F}_1)$ est définie comme suit. Prenez un élément $\tau \in H^0(X,
{\mathcal F}_3)$. Choisissez une couverture ouverte ${\mathcal U} : X =
\bigcup_{i\in I} U_i$ telle que pour chaque $i$ il existe une section
$\tilde \tau_i \in {\mathcal F}_2$ soulevant la restriction de $\tau$ à $U_i$.
$$
(i_0, i_1) \longmapsto
\tilde \tau_{i_0}|_{U_{i_0i_1}} - \tilde \tau_{i_1}|_{U_{i_0i_1}}.
$$
Il s'agit clairement d'un 1-coboundary dans le complexe {\v
C}ech ${\check C}^\ast({\mathcal U}, {\mathcal F}_2)$. Mais nous observons que
(en pensant à ${\mathcal F}_1$ comme un sous-seuil de ${\mathcal F}_2$) le
RHS est toujours une section de ${\mathcal F}_1$ sur $U_{i_0i_1}$.
Par conséquent, nous voyons que l'assignation définit une 1-cochain
dans le complexe ${\check C}^\ast({\mathcal U}, {\mathcal F}_2)$. La classe
de cohomologie de cette 1-cochain est par définition {\it $\delta(\tau)$}.
\end{definition}



\section{Cohomologie}
\label{section-cohomology}

\begin{exercise}
\label{exercise-cohomology-not-zero}
Laisse-moi. $X = \mathbf{R}$ En utilisant seulement les
propriétés formelles de la cohomologie (functorialité et longue
séquence de cohomologie exacte) montrent qu'il existe
une gerbe $\mathcal{F}$ le $X$ avec non-zéro $H^1(X, \mathcal{F})$C'est vrai.
\end{exercise}

\begin{exercise}
\label{exercise-mayer-vietoris}
Soit $X = U \cup V$ être un espace topologique écrit
comme l'union de deux s'ouvre. Ensuite, nous avons une
longue séquence exacte Mayer-Vietoris
$$
0 \to
H^0(X, \mathcal{F}) \to
H^0(U, \mathcal{F}) \oplus H^0(V, \mathcal{F}) \to
H^0(U \cap V, \mathcal{F}) \to
H^1(X, \mathcal{F}) \to \ldots
$$
Quelle propriété des gerbes injectables est essentielle pour la construction de
la longue séquence exacte de Mayer-Vietoris ? Pourquoi tient-elle ?
\end{exercise}

\begin{exercise}
\label{exercise-cohomology-two-point-space}
Soit $X$ être un espace topologique.
\begin{enumerate}
\item Afficher que $H^i(X, \mathcal{F})$ est zéro pour $i > 0$ si
$X$ a $2$ ou moins de points.
\item Et si $X$ avait des points $3$?
\end{enumerate}
\end{exercise}

\begin{exercise}
\label{exercise-cohomology-spec-local-ring}
Soit $X$ être le spectre d'un anneau local. Montrer
que $H^i(X, \mathcal{F})$ est zéro pour $i > 0$
et toute gerbe de groupes abeliens $\mathcal{F}$.
\end{exercise}

\begin{exercise}
\label{exercise-affine-morphism}
Soit $f : X \to Y$ être un
morphisme affine des schémas. Prouver que $H^i(X, \mathcal{F}) =
H^i(Y, f_*\mathcal{F})$ pour tout module quasi-cohérent $\mathcal{O}_X$ $\mathcal{F}$. N'hésitez pas
à imposer d'autres conditions sur $X$ et $Y$ et
à utiliser l'accord de {\v C}ech cohomology avec la cohomologie
pour les gerbes quasi-cohérentes et les couvertures d'affine ouvertes des schémas séparés.
\end{exercise}

\begin{exercise}
\label{exercise-affine-morphism-to-Pn}
Soit $A$ être un anneau. Soit $\mathbf{P}^n_A = \text{Proj}(A[T_0, \ldots, T_n])$ être
un espace projectif sur $A$.
Soit $\mathbf{A}^{n + 1}_A = \Spec(A[T_0, \ldots, T_n])$ et laissez
$$
U = \bigcup\nolimits_{i = 0, \ldots, n} D(T_i) \subset \mathbf{A}^{n + 1}_A
$$
être le complément de l'image de
l'immersion fermée $0 : \Spec(A)
\to \mathbf{A}^{n + 1}_A$.
$$
f : U \longrightarrow \mathbf{P}^n_A
$$
et prouver que $f_*\mathcal{O}_U
= \bigoplus_{d \in \mathbf{Z}} \mathcal{O}_{\mathbf{P}^n_A}(d)$. Plus généralement, montrer que
pour un module $A[T_0, \ldots, T_n]$ classé $M$ on
a
$$
f_*(\widetilde{M}|_U) =
\bigoplus\nolimits_{d \in \mathbf{Z}} \widetilde{M(d)}
$$
où sur le côté gauche nous avons la gaine
quasi-cohérente $\widetilde{M}$ associée à $M$ sur $\mathbf{A}^{n + 1}_A$
et sur le membre de droite nous avons
la gaine quasi-cohérente $\widetilde{M(d)}$ associée au module classé $M(d)$.
\end{exercise}

\begin{exercise}
\label{exercise-compute-Pn}
Soit $A$ être un anneau et laissez $\mathbf{P}^n_A = \text{Proj}(A[T_0, \ldots, T_n])$
être un espace projectif
sur $A$. Calculer soigneusement la cohomologie des torsions
Serre $\mathcal{O}_{\mathbf{P}^n_A}(d)$ de la gaine de structure
sur $\mathbf{P}^n_A$. N'hésitez pas à utiliser {\v C}ech cohomology et
l'accord de {\v C}ech cohomology avec la cohomologie pour
les gerbes quasi-cohérentes et les couvertures d'affines ouvertes de schémas séparés.
\end{exercise}

\begin{exercise}
\label{exercise-cohomology-hypersurface}
Soit $A$ être un anneau et laissez $\mathbf{P}^n_A = \text{Proj}(A[T_0, \ldots, T_n])$
être un espace projectif sur
$A$. Soit $F \in A[T_0, \ldots, T_n]$ être homogène
de degré $d$. Soit $X \subset \mathbf{P}^n_A$ être le
sous-régime fermé correspondant à l'idéal classé $(F)$ de $A[T_0, \ldots,
T_n]$. Que pouvez-vous dire de $H^i(X, \mathcal{O}_X)$?
\end{exercise}

\begin{exercise}
\label{exercise-characterize-finite-product-fields}
Soit $R$ être un anneau tel que pour n'importe quel functeur exact
gauche $F : \text{Mod}_R \to \textit{Ab}$ nous avons $R^iF
= 0$ pour $i > 0$. Montrer que $R$ est un produit fini
de champs.
\end{exercise}






\section{Plus de cohomologie}
\label{section-more-cohomology}

\begin{exercise}
\label{exercise-cohomology-coordinate-axes}
Soit $k$ être un corps. Soit $X \subset \mathbf{P}^n_k$
être la ``croix coordonnée''. Autrement dit, laissez $X$
être défini par les équations homogènes
$$
T_i T_j = 0\text{ for }i > j > 0
$$
où, comme d'habitude, nous écrivons $\mathbf{P}^n_k = \text{Proj}(k[T_0, \ldots, T_n])$. En
d'autres termes, $X$ est le sous-régime fermé correspondant au quotient $k[T_0,
\ldots, T_n]/(T_iT_j; i > j > 0)$ de l'anneau
polynôme. Calculer $H^i(X, \mathcal{O}_X)$ pour tous $i$Conseil: utilisation {\v C}La cohomologie.
\end{exercise}

\begin{exercise}
\label{exercise-compute-cohomology-punctured}
Soit $A$ être un anneau. Soit $I = (f_1,
\ldots, f_t)$ être un idéal fini
de $A$. Soit $U \subset \Spec(A)$ être le complément de
$V(I)$. Pour tout module $A$ $M$ écrire un
complexe de modules $A$ (en termes de $A$, $f_1, \ldots,
f_t$, $M$) dont les groupes de cohomologie donnent $H^n(U, \widetilde{M})$.
\end{exercise}

\begin{exercise}
\label{exercise-compute-cohomology-affine-space-punctured}
Soit $k$ être un corps. Soit $U \subset \mathbf{A}^d_k$
être le complément du point fermé $0$ de
$\mathbf{A}^d_k$. Calculer $H^n(U, \mathcal{O}_U)$ pour tous $n$.
\end{exercise}

\begin{exercise}
\label{exercise-find-curve-genus-one-hundred}
Soit $k$ être un corps. Trouver explicitement un schéma $X$ projectif sur
$k$ de dimension $1$ avec $H^0(X, \mathcal{O}_X) = k$
et $\dim_k H^1(X, \mathcal{O}_X) = 100$.
\end{exercise}

\begin{exercise}
\label{exercise-degree-2-cover}
Soit $f : X \to Y$ être un morphisme fini localement libre
de degré $2$. Supposons que $X$ et $Y$ sont des schémas intégraux et
que $2$ est invertible dans la gaine de structure de $Y$, c'est-à-dire $2
\in \Gamma(Y, \mathcal{O}_Y)$ est invertible. Montrer que la carte $\mathcal{O}_Y$-module
$$
f^\sharp : \mathcal{O}_Y \longrightarrow f_*\mathcal{O}_X
$$
a un inverse gauche, c'est-à-dire qu'il y a une
carte $\mathcal{O}_Y$-module $\tau : f_*\mathcal{O}_X \to
\mathcal{O}_Y$ avec $\tau \circ f^\sharp =
\text{id}$. Conclure que $H^n(Y, \mathcal{O}_Y) \to H^n(X, \mathcal{O}_X)$ est
injectable\footnote{Il existe un morphisme fini localement libre $X \to
Y$ entre des schémas intégraux de degré $2$ où
la carte $H^1(Y, \mathcal{O}_Y) \to H^1(X, \mathcal{O}_X)$ n'est pas injectable.}.
\end{exercise}

\begin{exercise}
\label{exercise-pic}
Soit $X$ être un schéma (ou un espace
sonné localement). La règle $U \mapsto \mathcal{O}_X(U)^*$ définit
une masse de groupes
désignés $\mathcal{O}_X^*$. Expliquer brièvement pourquoi le
groupe Picard de $X$ (Définition \ref{definition-picard-group})
est égal à $H^1(X, \mathcal{O}_X^*)$.
\end{exercise}

\begin{exercise}
\label{exercise-pic-nontrivial}
Donner un exemple d'un schéma d'affine $X$ avec $\Pic(X)$ non trivial.
Conclure en utilisant l'exercice \ref{exercise-pic} que $H^1(X, -)$ n'est pas le functeur
zéro pour un tel $X$.
\end{exercise}

\begin{exercise}
\label{exercise-kill-cohomology-complement}
Laisse-moi. $A$ Sois une bague. $I = (f_1, \ldots, f_t)$ être un idéal
finiment généré de $A$- Laisse-moi faire. $U \subset \Spec(A)$ soit le
complément de $V(I)$. Vu un quasi-cohérent $\mathcal{O}_{\Spec(A)}$-Module $\mathcal{F}$ et $\xi \in H^p(U,
\mathcal{F})$ avec $p > 0$, montrer qu'il existe $n > 0$ de
telle sorte que $f_i^n \xi = 0$ pour $i =
1, \ldots, t$. Conseil: Une façon possible de procéder est
d'utiliser le complexe que vous avez trouvé dans l'exercice \ref{exercise-compute-cohomology-punctured}C'est vrai.
\end{exercise}

\begin{exercise}
\label{exercise-h0-complement}
Soit $A$ être un anneau. Soit $I = (f_1, \ldots,
f_t)$ être un idéal fini de $A$. Soit $U \subset
\Spec(A)$ être le complément de $V(I)$. Soit $M$ être un
module $A$ dont $I$-torsion est zéro,
c'est-à-dire $0 = \Ker((f_1, \ldots, f_t) : M \to M^{\oplus
t})$. Montrer qu'il y a un isomorphisme canonique
$$
H^0(U, \widetilde{M}) = \colim \Hom_A(I^n, M).
$$
Attention : ce n'est pas trivial.
\end{exercise}

\begin{exercise}
\label{exercise-Noetherian}
Soit $A$ être un anneau noethérien. Soit $I$ être un idéal de
$A$. Soit $M$ être un module $A$. Soit $M[I^\infty]$ être l'ensemble
d'éléments de torsion de puissance $I$ défini par
$$
M[I^\infty] = \{x \in M \mid
\text{ there exists an }n \geq 1\text{ such that }I^nx = 0\}
$$
Définissez $M' = M/M[I^\infty]$. Ensuite, la torsion de $I$-puissance de $M'$ est
zéro.
$$
\colim \Hom_A(I^n, M) = \colim \Hom_A(I^n, M').
$$
Avertissement : ce n'est pas trivial. Astuces : (1) essayer de réduire à $M$
fini, (2) montrer tout élément de $\Ext^1_A(I^n, N)$
cartes à zéro dans $\Ext^1_A(I^{n + m}, N)$ pour certains $m >
0$ si $N = M[I^\infty]$ et $M$ fini, (3) montrer la même
chose que dans (2) pour $\Hom_A(I^n, N)$, (3) considérer la longue
séquence exacte
$$
0 \to \Hom_A(I^n, M[I^\infty]) \to \Hom_A(I^n, M) \to \Hom_A(I^n, M')
\to \Ext^1_A(I^n, M[I^\infty])
$$
pour $M$ fini et comparer avec la séquence pour $I^{n + m}$ pour conclure.
\end{exercise}








\section{Cohomologie revue}
\label{section-more-more-cohomology}


\begin{exercise}
\label{exercise-nonkillable}
Exemple d'un corps $k$, une courbe $X$ sur
$k$, une invertible $\mathcal{O}_X$-Module $\mathcal{L}$ et une classe de
cohomologie $\xi \in H^1(X, \mathcal{L})$ avec la propriété
suivante: pour chaque morphisme fini surjectif $\pi
: Y \to X$ de régimes l'élément $\xi$ tire
vers un élément non zéro de $H^1(Y, \pi^*\mathcal{L})$. Conseil: construction
$X$, $k$, $\mathcal{L}$ si bien qu'il y a
une courte séquence exacte
$0 \to \mathcal{L} \to \mathcal{O}_X \to i_*\mathcal{O}_Z \to 0$ où
$Z \subset X$ est un sous-régime fermé composé de
plus de $1$ Regardez ce qui se passe quand vous
reculez.
\end{exercise}

\begin{exercise}
\label{exercise-cohomology-cycle-curves}
Laisse-moi. $k$ être un corps algébrique fermé.
$X$ être un projetif $1$Supposons que ce soit le
cas. $X$ contient un cycle de courbes, c'est-à-dire supposons qu'il
existe un $n \geq 2$ et distincte par paire
$1$-sous-régimes dimensionnels intégrés fermés $C_1, \ldots, C_n$
et des points fermés distincts dans le sens de la
paire $x_1, \ldots, x_n \in X$ de telle sorte que
$x_n \in C_n \cap C_1$ et
$x_i \in C_i \cap C_{i
+ 1}$ pour $i = 1, \ldots, n
- 1$- Prouver-le. $H^1(X, \mathcal{O}_X)$ n'est pas zéro. $\mathcal{F}$
être l'image de la carte $\mathcal{O}_X \to \bigoplus
\mathcal{O}_{C_i}$, et montrer $H^1(X, \mathcal{F})$ n'est pas zéro
en utilisant $\kappa(x_i) = k$ et $H^0(C_i, \mathcal{O}_{C_i}) = k$. Utilisez aussi
que $H^2(X, -) = 0$ par le théorème de Grothendieck.
\end{exercise}

\begin{exercise}
\label{exercise-surface}
Laisse-moi. $X$ être une surface projective au-dessus d'un corps algébrique fermé
$k$. Prouver qu'il existe un sous-régime fermé approprié $Z \subset X$
de telle sorte que $H^1(Z, \mathcal{O}_Z)$ n'est
pas zéro. Conseil: Utilisez l'exercice \ref{exercise-cohomology-cycle-curves}C'est vrai.
\end{exercise}

\begin{exercise}
\label{exercise-surface-better}
Soit $X$ être une surface projective sur un corps fermé algébriquement
$k$. Montrer que pour chaque
$n \geq 0$ il existe un sous-schème fermé approprié
$Z \subset X$ tel que $\dim_k
H^1(Z, \mathcal{O}_Z) > n$. Expliquer seulement comment le
faire en modifiant les arguments dans
Exercice \ref{exercise-surface} et \ref{exercise-cohomology-cycle-curves}; ne donnez pas tous les détails.
\end{exercise}

\begin{exercise}
\label{exercise-surface-h2-nonzero}
Laisse-moi. $X$ être une surface projective au-dessus d'un corps algébrique
fermé $k$. Prouver qu'il existe une cohérence $\mathcal{O}_X$-Module $\mathcal{F}$ de
telle sorte que $H^2(X, \mathcal{F})$ n'est pas zéro. Conseil: Utilisez
le résultat de l'exercice \ref{exercise-surface-better} et une séquence exacte
habilement choisie.
\end{exercise}

\begin{exercise}
\label{exercise-kunneth}
Laisse-moi. $X$ et $Y$ être des régimes sur un domaine $k$ (se
sentir libre de supposer $X$ et $Y$ sont agréables, par exemple
qcqs ou projective sur $k$). Prêt $X \times Y = X
\times_{\Spec(k)} Y$ avec des projections $p : X \times Y
\to X$ et $q : X \times Y \to Y$.
Pour un quasi-cohérent $\mathcal{O}_X$-Module $\mathcal{F}$ et un quasi-cohérent $\mathcal{O}_Y$-Module $\mathcal{G}$ prouver que
$$
H^n(X \times Y,
p^*\mathcal{F} \otimes_{\mathcal{O}_{X \times Y}} q^*\mathcal{G}) =
\bigoplus\nolimits_{a + b = n}
H^a(X, \mathcal{F}) \otimes_k H^b(Y, \mathcal{G})
$$
ou simplement montrer que cela tient quand on prend des dimensions.
Points supplémentaires pour des solutions « propres ».
\end{exercise}

\begin{exercise}
\label{exercise-Oab}
Soit $k$ être un corps. Soit $X = \mathbf{P}|^1 \times \mathbf{P}^1$ être
le produit de la ligne projective au-dessus de $k$ avec
elle-même avec des projections $p : X \to \mathbf{P}^1_k$ et $q : X \to
\mathbf{P}^1_k$.
$$
\mathcal{O}(a, b) = p^*\mathcal{O}_{\mathbf{P}^1_k}(a)
\otimes_{\mathcal{O}_X} q^*\mathcal{O}_{\mathbf{P}^1_k}(b)
$$
Calculer les dimensions de $H^i(X, \mathcal{O}(a, b))$
pour tous $i, a, b$. Conseil: Utilisez Exercice \ref{exercise-kunneth}.
\end{exercise}










\section{Cohomologie et polynômes Hilbert}
\label{section-cohomology-hilbert-polynomials}

\begin{situation}
\label{situation-hilbert-polynomial}
Soit $k$ être un corps. Soit $X
= \mathbf{P}^n_k$ être $n$-dimensionnel espace projectif. Soit $\mathcal{F}$
être un module cohérent $\mathcal{O}_X$. Rappelez-vous que
$$
\chi(X, \mathcal{F}) =
\sum\nolimits_{i = 0}^n (-1)^i \dim_k H^i(X, \mathcal{F})
$$
Rappelons que la fonction {\it Hilbert polynomial} de $\mathcal{F}$ est la fonction
$$
t \longmapsto \chi(X, \mathcal{F}(t))
$$
Nous rappelons
également que $\mathcal{F}(t) = \mathcal{F} \otimes_{\mathcal{O}_X} \mathcal{O}_X(t)$ où $\mathcal{O}_X(t)$
est la $t$e torsion de la gaine de
structure comme dans Constructions, Définition \ref{constructions-definition-twist}. Dans Variétés,
sous-section \ref{varieties-subsection-hilbert} nous avons prouvé que le polynôme
Hilbert est un polynôme dans $t$.
\end{situation}

\begin{exercise}
\label{exercise-hilbert-pol-easy}
En situation \ref{situation-hilbert-polynomial}.
\begin{enumerate}
\item Si $P(t)$ est le polynôme de Hilbert de
$\mathcal{F}$, quel est le polynôme de Hilbert de $\mathcal{F}(-13)$.
\item Si $P_i$ est le polynôme de Hilbert de
$\mathcal{F}_i$, quel est le polynôme de Hilbert de $\mathcal{F}_1 \oplus \mathcal{F}_2$.
\item Si $P_i$ est le polynôme de Hilbert de $\mathcal{F}_i$
et $\mathcal{F}$ est le noyau d'une carte
de surjectif $\mathcal{F}_1 \to \mathcal{F}_2$, qu'est-ce que le polynôme de Hilbert
de $\mathcal{F}$?
\end{enumerate}
\end{exercise}

\begin{exercise}
\label{exercise-find-given-hilbert-pol-dim-1}
Dans la situation \ref{situation-hilbert-polynomial} supposez $n \geq 1$. Trouver une
gaine cohérente dont Hilbert polynôme est $t - 101$.
\end{exercise}

\begin{exercise}
\label{exercise-find-given-hilbert-pol-dim-2}
Dans Situation \ref{situation-hilbert-polynomial} supposez $n \geq 2$. Trouver une gerbe cohérente
dont Hilbert polynôme est $t^2/2 + t/2 -
1$. (C'est un peu délicat; il suffit de
montrer qu'il y a une gerbe aussi cohérente.)
\end{exercise}

\begin{exercise}
\label{exercise-bound-genus-in-degree}
Dans la situation \ref{situation-hilbert-polynomial} supposer $n \geq 2$ et $k$ algébriquement
fermé. Laisser $C \subset X$
être un sous-régime fermé intégral de dimension $1$. En d'autres termes,
$C$ est une courbe projective. Laisser $d
t + e$ être le polynôme Hilbert de $\mathcal{O}_C$ considéré
comme une gaine cohérente sur $X$.
\begin{enumerate}
\item Donner une limite supérieure sur $e$. (Aperçu
: Utilisez que $\mathcal{O}_C(t)$ n'a de cohomologie qu'en degrés $0$
et $1$ et étudiez $H^0(C, \mathcal{O}_C)$.)
\item Choisir une section globale $s$ des $\mathcal{O}_X(1)$ qui se
croisent $C$ transversalement, c'est-à-dire de manière à ce qu'il y ait
des points fermés distincts par paire $c_1, \ldots, c_r \in C$ et
une courte séquence exacte
$$
0 \to \mathcal{O}_C \xrightarrow{s} \mathcal{O}_C(1) \to
\bigoplus\nolimits_{i = 1, \ldots, r} k_{c_i} \to 0
$$
où $k_{c_i}$ est la gaine de gratte-ciel avec la valeur $k$ dans $c_i$. (Such
an $s$ existe; s'il vous plaît utilisez simplement ceci.)
Montrer que $r = d$. (Astuce: tordez la séquence et voyez ce que vous obtenez.)
\item La séquence courte donne une carte $k$ linéaire $\varphi_t :
\Gamma(C, \mathcal{O}_C(t)) \to \bigoplus_{i = 1, \ldots, d} k$ pour tout $t$. Affichez
que si cette carte est surjective pour $t \geq d - 1$.
\item Donner une limite inférieure sur $e$ en termes de $d$. (Aperçu : montrer
que $H^1(C, \mathcal{O}_C(d - 2)) = 0$ en utilisant le résultat de (3) et utiliser
la disparition.)
\end{enumerate}
\end{exercise}

\begin{exercise}
\label{exercise-three-quadrics-in-plane}
Dans la situation \ref{situation-hilbert-polynomial} supposer $n =
2$. Laisser $s_1, s_2, s_3 \in \Gamma(X,
\mathcal{O}_X(2))$ être trois équations quadriques.
$$
\mathcal{F} = \Coker\left(\mathcal{O}_X(-2)^{\oplus 3}
\xrightarrow{s_1, s_2, s_3} \mathcal{O}_X\right)
$$
Énumérez les polynômes possibles de Hilbert de ces $\mathcal{F}$.
(Essayez de visualiser les intersections des quadrics dans le plan projectif.)
\end{exercise}





\section{Courbes}
\label{section-curves}

\begin{exercise}
\label{exercise-closed-subset-vanshing}
Laisser $k$ être un corps algébrique fermé. Laisser $X$ être
une courbe de projection sur $k$. Laisser $\mathcal{L}$ être
un module inverse $\mathcal{O}_X$. Laisser $s_0,
\ldots, s_n \in H^0(X, \mathcal{L})$
être des sections globales de $\mathcal{L}$.
$$
Z \subset
\mathbf{P}^n \times X
$$
de telle sorte que le
point fermé $((\lambda_0 : \ldots : \lambda_n),
x)$ soit dans $Z$ si et seulement
si la section $\lambda_0 s_0 + \ldots +
\lambda_n s_n$ disparaît à $x$. (Enfin, décrire $Z$ affine localement.)
\end{exercise}

\begin{exercise}
\label{exercise-diagonal-cartier}
Soit $k$ être un corps algébrique fermé. Soit $X$ être une courbe lisse sur $k$.
Soit $r \geq 1$. Montrer que le sous-ensemble fermé
$$
D \subset X \times X^r = X^{r + 1}
$$
dont les points fermés sont les tuples $(x, x_1, \ldots, x_r)$ avec $x
= x_i$ pour certains $i$, a une gaine idéale invertible. En d'autres termes, montrer
que $D$
est un diviseur Cartier efficace. Astuces: réduire à $r = 1$ et utiliser
que $X$ est une courbe lisse pour dire quelque chose sur la diagonale (regardez
dans les livres pour cela).
\end{exercise}

\begin{exercise}
\label{exercise-one-divisor-in-another}
Laisse-moi. $k$ être un corps algébrique fermé. $X$ être une courbe de projection lisse sur $k$-
Laisse-moi faire. $T$ être un schéma de type fini sur $k$ et laisse
$$
D_1 \subset X \times T
\quad\text{and}\quad
D_2 \subset X \times T
$$
soit deux diviseurs Cartier efficaces tels que pour $t \in T$ les
fibres $D_{i, t} \subset X_t$ ne sont pas denses (c.-à-d. ne contiennent pas
le point générique de la courbe $X_t$).
Prouver qu'il existe un sous-scheme fermé canonique $Z \subset T$
tel qu'un point fermé $t \in T$ est dans $Z$ si
et seulement si pour le schéma des fibres théorétiques $D_{1, t}$, $D_{2, t}$
de $D_1$, $D_2$ nous avons
$$
D_{1, t} \subset D_{2, t}
$$
en tant que sous-régimes fermés de
$X_t$. Conseils: Montrer que, peut-être après avoir rétréci $T$, vous
pouvez supposer que $T = \Spec(A)$ est affine et qu'il y a un affine
ouvert $U \subset X$ tel que $D_i \subset U \times
T$. Ensuite montrez que $M_1 = \Gamma(D_1, \mathcal{O}_{D_1})$ est un module fini localement libre
$A$ (ici vous aurez besoin d'une algèbre non triviale --- demandez à vos
amis). Après rétrécissement $T$ vous pouvez supposer que $M_1$ est un module
libre $A$.
$$
\Gamma(U \times T, \mathcal{I}_{D_2}) \to M_1 = A^{\oplus N}
$$
et vous définissez $Z$ comme le sous-régime fermé coupé par l'idéal généré
par les coefficients de vecteurs dans l'image de cette carte. Expliquer pourquoi
cela fonctionne (cela nécessitera peut-être un peu plus d'algèbre commutative).
\end{exercise}

\begin{exercise}
\label{exercise-closed-subset-vanshing-r}
Laisser $k$ être un corps algébrique fermé. Laisser $X$ être une
courbe de projection lisse sur $k$. Laisser $\mathcal{L}$ être un
module inverse $\mathcal{O}_X$. Laisser $s_0, \ldots, s_n
\in H^0(X, \mathcal{L})$ être des sections globales
de $\mathcal{L}$. Laisser $r \geq 1$.
$$
Z \subset
\mathbf{P}^n \times X \times \ldots \times X = \mathbf{P}^n \times X^r
$$
de telle sorte que
le point clos $((\lambda_0 : \ldots : \lambda_n), x_1,
\ldots, x_r)$ est en $Z$ si
et seulement si la section $s_\lambda = \lambda_0
s_0 + \ldots + \lambda_n s_n$ disparaît sur le
diviseur $D = x_1 + \ldots + x_r$,
c'est-à-dire la section $s_\lambda$ est en $\mathcal{L}(-D)$. Conseil: expliquer comment
ceci suit en combinant ensuite les résultats
des exercices \ref{exercise-diagonal-cartier} et \ref{exercise-one-divisor-in-another}C'est vrai.
\end{exercise}

\begin{exercise}
\label{exercise-effective}
Soit $k$ être un corps algébrique fermé. Soit $X$ être une courbe de
projection lisse sur $k$. Soit $\mathcal{L}$ être un module invertible $\mathcal{O}_X$. Montrer
qu'il y a un sous-ensemble fermé naturel
$$
Z \subset X^r
$$
de telle sorte qu'un point clos $(x_1, \ldots, x_r)$ des $X^r$
est en $Z$ si et seulement si $\mathcal{L}(-x_1 - \ldots
-x_r)$ a une section globale non zéro. Conseil: utiliser Exercice \ref{exercise-closed-subset-vanshing-r}C'est vrai.
\end{exercise}

\begin{exercise}
\label{exercise-difference-effective}
Soit $k$ être un corps fermé algébriquement. Soit $X$ être une courbe de projection lisse
sur $k$. Soit $r \geq s$ être entiers. Montrer qu'il y a un sous-ensemble fermé
naturel
$$
Z \subset X^r \times X^s
$$
de sorte qu'un point fermé $(x_1, \ldots, x_r, y_1, \ldots, y_s)$ de
$X^r \times X^s$ soit
dans $Z$ si et seulement si $x_1 + \ldots + x_r
- y_1 - \ldots - y_s$ est linéairement
équivalent à un diviseur efficace. Conseil: Choisissez un module inverse
auxiliaire $\mathcal{L}$ de très haut degré de sorte que $\mathcal{L}(-D)$ ait
une section non avantageuse pour tout diviseur efficace $D$ de
degré $r$. Ensuite, utilisez le résultat de l'exercice \ref{exercise-effective} deux fois.
\end{exercise}

\begin{exercise}
\label{exercise-genus-7}
Choisissez votre corps algébrique préféré $k$. Autant
que vous pouvez déterminer tous les $\mathfrak g^r_d$
possibles qui peuvent exister sur une
courbe du genre $7$.
\begin{enumerate}
\item déterminer dans quels cas le $\mathfrak g^r_d$ est libre de point de base, et
\item déterminer dans quels cas le $\mathfrak g^r_d$ donne un
encastrement fermé dans $\mathbf{P}^r$.
\end{enumerate}
Faites la même chose si vous supposez que votre courbe est "générale"
(faite votre propre notion de général - cela peut être plus facile que la question ci-dessus).
Faites la même chose si vous supposez que votre courbe est hyperelliptique.
Faites la même chose si vous supposez que votre courbe
est trigonale (et non hyperelliptique). Etc.
\end{exercise}





\section{Moduli}
\label{section-moduli}

\noindent
Dans cette section, nous examinons quelques approches naïves des
moduli d'objets géométriques algébriques.

\medskip\noindent
Laisse-moi. $k$ être un corps algébrique fermé. Supposons que $M$ est un problème de
moduli plus $k$. Nous ne définirons pas exactement ce que cela signifie ici, mais
dans chaque exercice il devrait être clair ce que nous voulons dire.
Pour comprendre ce qui suit il suffit de savoir ce
que les objets de $M$ sur $k$ sont, ce que les isomorphismes entre les objets
de $M$ sur $k$ sont, et ce que les familles de l'objet de $M$
sur une variété sont. Alors nous disons le {\it nombre de modules $M$} est $d
\geq 0$ si ce qui suit est vrai
\begin{enumerate}
\item il y a un nombre fini de familles $X_i \to V_i$, $i = 1, \ldots, n$
de sorte que chaque objet de $M$ sur $k$ est isomorphe à une fibre de l'un
d'eux et de telle sorte que $\max \dim(V_i) = d$, et
\item il n'y a aucun moyen de le faire avec un plus petit $d$.
\end{enumerate}
Ce n'est vraiment qu'une approximation très grossière des meilleures notions dans
la littérature.

\begin{exercise}
\label{exercise-number-moduli-lines-conics}
Soit $k$ être un corps algébrique fermé. Soit $d \geq 1$ et $n \geq
1$. Supposons que le module des hypersurfaces de degré $d$ dans $P^n$ est donné par
\begin{enumerate}
\item un objet est une hypersurface $X \subset \mathbf{P}^n_k$ de
degré $d$,
\item un isomorphisme entre deux objets $X \subset \mathbf{P}^n_k$ et
$Y \subset \mathbf{P}^n_k$ est un élément $g \in \text{PGL}_n(k)$ tel que
$g(X) = Y$, et
\item une famille d'hypersurfaces sur une variété $V$ est un sous-schème
fermé $X \subset \mathbf{P}^n_V$ de telle sorte que pour tous $v \in
V$ la fibre théorétique schématique $X_v$ de $X
\to V$ est une hypersurfaces dans $\mathbf{P}^n_v$.
\end{enumerate}
Calculer (si vous le pouvez -- ceux-ci deviennent progressivement plus difficiles)
\begin{enumerate}
\item le nombre de modules lorsque $n = 1$ et $d$ arbitraires,
\item le nombre de modules lorsque $n = 2$ et $d = 1$,
\item le nombre de modules lorsque $n = 2$ et $d = 2$,
\item le nombre de modules lorsque $n \geq 1$ et $d = 2$,
\item le nombre de modules lorsque $n = 2$ et $d = 3$,
\item le nombre de modules lorsque $n = 3$ et $d = 3$, et
\item le nombre de modules lorsque $n = 2$ et $d = 4$.
\end{enumerate}
\end{exercise}

\begin{exercise}
\label{exercise-number-moduli-hyperelliptic}
Soit $k$ être un corps algébrique fermé. Soit $g \geq 2$.
Disons que les moduli des courbes hyperelliptiques du genre $g$ sont donnés par
\begin{enumerate}
\item un objet est une courbe hyperelliptique projective lisse $X$ du
genre $g$,
\item un isomorphisme entre deux objets $X$ et $Y$ est un
isomorphisme $X \to Y$ de schémas sur $k$, et
\item une famille de courbes hyperelliptiques du genre $g$ sur une variété $V$ est
une bonne plate\footnote{Vous pouvez laisser tomber cette hypothèse sans changer la réponse
à la question.} morphisme $X \to Y$ de sorte que toutes les fibres
théorétiques du schéma de $X \to V$ sont des courbes hyperelliptiques projectives lisses du
genre $g$.
\end{enumerate}
Afficher que le nombre de modules est $2g - 1$.
\end{exercise}





\section{Exts mondiaux}
\label{section-global-exts}

\begin{exercise}
\label{exercise-ext-line-plane-p3}
Soit $k$ être un corps. Soit $X = \mathbf{P}^3_k$. Soit $L
\subset X$ et $P \subset X$ être une ligne et un plan,
considérés comme des sous-régimes fermés découpés par $1$, resp.,
$2$ équations linéaires.
$$
\Ext^i_X(\mathcal{O}_L, \mathcal{O}_P)
$$
Pour tous $i$.
Assurez-vous de faire à la fois le cas où $L$ est contenu dans $P$ et
le cas où $L$ n'est pas contenu dans $P$.
\end{exercise}

\begin{exercise}
\label{exercise-ext-point}
Soit $k$ être un corps. Soit $X = \mathbf{P}^n_k$. Soit $Z \subset
X$ être un point rationnel $k$ fermé, vu comme un sous-régime fermé.
Par exemple, le point avec des coordonnées homogènes $(1 :
0 : \ldots : 0)$. Calculer les dimensions de
$$
\Ext^i_X(\mathcal{O}_Z, \mathcal{O}_Z)
$$
pour tous $i$.
\end{exercise}

\begin{exercise}
\label{exercise-ext-pairings}
Soit $X$ être un espace sonné.
$$
\Ext^i_X(\mathcal{G}, \mathcal{H})
\times
\Ext^j_X(\mathcal{F}, \mathcal{G})
\longrightarrow
\Ext^{i + j}_X(\mathcal{F}, \mathcal{H})
$$
pour $\mathcal{O}_X$-modules $\mathcal{F}, \mathcal{G}, \mathcal{H}$. (Astuce : c'est
une chose super générale.)
\end{exercise}

\begin{exercise}
\label{exercise-move-out-locally-free}
Soit $X$ être un espace sonné. Soit $\mathcal{E}$
être un module local
fini sans $\mathcal{O}_X$ avec double $\mathcal{E}^\vee = \SheafHom_{\mathcal{O}_X}(\mathcal{E}, \mathcal{O}_X)$.
Prouver les instructions suivantes
\begin{enumerate}
\item $\SheafExt^i_{\mathcal{O}_X}(
\mathcal{F} \otimes_{\mathcal{O}_X} \mathcal{E}, \mathcal{G})
= \SheafExt^i_{\mathcal{O}_X}(
\mathcal{F},
\mathcal{E}^\vee \otimes_{\mathcal{O}_X} \mathcal{G}) =
\SheafExt^i_{\mathcal{O}_X}(
\mathcal{F}, \mathcal{G}) \otimes_{\mathcal{O}_X} \mathcal{E}^\vee$, et
\item
$\Ext^i_X( \mathcal{F} \otimes_{\mathcal{O}_X} \mathcal{E}, \mathcal{G})
=
\Ext^i_X(\mathcal{F}, \mathcal{E}^\vee \otimes_{\mathcal{O}_X} \mathcal{G})$.
\end{enumerate}
Ici, $\mathcal{F}$ et $\mathcal{G}$ sont $\mathcal{O}_X$-modules.

$$
\Ext^i_X(\mathcal{E}, \mathcal{G}) =
H^i(X, \mathcal{E}^\vee \otimes_{\mathcal{O}_X} \mathcal{G})
$$
\end{exercise}

\begin{exercise}
\label{exercise-trace-on-ext}
Soit $X$ être un espace sonné. Soit $\mathcal{E}$ être un module local
fini sans $\mathcal{O}_X$. Construire une carte de trace
$$
\Ext^i_X(\mathcal{E}, \mathcal{E}) \to H^i(X, \mathcal{O}_X)
$$
pour tous $i$. Généraliser vers une carte de trace
$$
\Ext^i_X(\mathcal{E}, \mathcal{E} \otimes_{\mathcal{O}_X} \mathcal{F})
\to H^i(X, \mathcal{F})
$$
pour tout module $\mathcal{O}_X$ $\mathcal{F}$.
\end{exercise}

\begin{exercise}
\label{exercise-duality}
Soit $k$ être un corps. Soit $X = \mathbf{P}^d_k$.
Définissez $\omega_{X/k} = \mathcal{O}_X(-d - 1)$. Prouver
que pour les modules finis sans local $\mathcal{E}$,
$\mathcal{F}$ le produit de tasse sur Ext combiné avec
la carte de trace sur Ext
$$
\Ext^i_X(\mathcal{E}, \mathcal{F} \otimes_{\mathcal{O}_X} \omega_{X/k})
\times
\Ext^{d - i}_X(\mathcal{F}, \mathcal{E})
\to
\Ext_X^d(\mathcal{F}, \mathcal{F} \otimes_{\mathcal{O}_X} \omega_{X/k})
\to
H^d(X, \omega_{X/k}) = k
$$
Vous pouvez soit reprouver la dualité dans ce cadre, soit réduire
à la cohomologie des gerbes et appliquer le théorème de la
dualité Serre comme le prouvent les conférences.
\end{exercise}







\section{Diviseurs}
\label{section-divisors}

\noindent
Nous recueillons toutes les définitions pertinentes ici en un seul endroit pour plus de commodité.

\begin{definition}
\label{definition-divisor}
Tout au long, laissez $S$ être n'importe quel schéma
et laissez $X$ être un schéma noethérien, intégral.
\begin{enumerate}
\item A {\it Le diviseur de Weil} sur $X$ est une combinaison
linéaire formelle $\Sigma n_i[Z_i]$ de diviseurs primaires $Z_i$ avec des coefficients entiers.
\item Un {\it diviseur primaire} est un sous-régime fermé $Z \subset X$,
qui fait partie intégrante du point générique $\xi \in Z$
tel que ${\mathcal O}_{X, \xi}$ a une dimension $1$. Nous utiliserons la notation
${\mathcal O}_{X, Z} = {\mathcal O}_{X, \xi}$ lorsque
$\xi \in Z \subset X$ sera comme ci-dessus. Notez que ${\mathcal O}_{X, Z} \subset
K(X)$ est un sous-ring du corps de fonction $X$.
\item Les {\it Diviseur faible associé à une fonction rationnelle
$f \in K(X)^\ast$} est la somme $\Sigma v_Z(f)[Z]$Ici. $v_Z(f)$ est
défini comme suit:
\begin{enumerate}
\item Si $f \in {\mathcal O}_{X, Z}^\ast$ puis $v_Z(f) = 0$.
\item Si $f \in {\mathcal O}_{X, Z}$ alors
$$
v_Z(f) = \text{length}_{{\mathcal O}_{X, Z}}({\mathcal O}_{X, Z}/(f)).
$$
\item Si $f = \frac{a}{b}$ avec $a, b \in {\mathcal O}_{X, Z}$
alors
$$
v_Z(f) = \text{length}_{{\mathcal O}_{X, Z}}({\mathcal O}_{X, Z}/(a)) -
\text{length}_{{\mathcal O}_{X, Z}}({\mathcal O}_{X, Z}/(b)).
$$
\end{enumerate}
\item Un {\it diviseur Cartier efficace} sur un schéma $S$ est
un sous-régime fermé $D \subset S$ de telle sorte que chaque point $d\in D$
a un quartier ouvert affine $\Spec(A) = U \subset S$ dans $S$ de
sorte que $D \cap U = \Spec(A/(f))$ avec $f \in A$ un non-zérodiviseur.
\item Les {\it Diviseur faible $[D]$ associé à un diviseur
Cartier efficace $D \subset X$} de notre système intégral noethérien
$X$ est définie comme la somme $\Sigma v_Z(D)[Z]$ où
$v_Z(D)$ est défini comme suit:
\begin{enumerate}
\item Si le point générique $\xi$ de $Z$ n'est pas dans $D$,
alors $v_Z(D) = 0$.
\item Si le point générique $\xi$ de $Z$ est dans $D$
alors
$$
v_Z(D) = \text{length}_{{\mathcal O}_{X, Z}}({\mathcal O}_{X, Z}/(f))
$$
où $f \in {\mathcal O}_{X, Z} = {\mathcal O}_{X, \xi}$ est le non-zérodiviseur qui
définit $D$ dans un quartier affine de $\xi$ (comme dans (4) ci-dessus).
\end{enumerate}
\item Laisser $S$ être un schéma. La gaine {\it des
anneaux quotients totaux ${\mathcal K}_S$} est la gaine de ${\mathcal O}_S$-algèbres qui est
la sheafification de la pré-file ${\mathcal K}'$ définie comme suit. Pour $U
\subset S$ ouvert, nous définissons ${\mathcal K}'(U) = S_U^{-1}{\mathcal O}_S(U)$ où $S_U \subset
{\mathcal O}_S(U)$ est le sous-ensemble multiplicatif composé de sections $f \in {\mathcal
O}_S(U)$ de sorte que le germe de $f$ dans ${\mathcal O}_{S,
u}$ est un diviseur non nul pour chaque $u\in U$. En particulier
les éléments de $S_U$ sont tous des diviseurs non nuls. Ainsi
${\mathcal O}_S$ est un sous-sheaf de ${\mathcal K}_S$, et nous obtenons une
courte séquence exacte
$$
0 \to {\mathcal O}_S^\ast \to {\mathcal K}_S^\ast \to
{\mathcal K}_S^\ast/{\mathcal O}_S^\ast \to 0.
$$
\item Un {\it diviseur Cartier} sur un schéma $S$
est une section globale de la gaine du quotient ${\mathcal K}_S^\ast/{\mathcal O}_S^\ast$.
\item Les {\it Diviseur Weil associé à un diviseur
Cartier} $\tau \in \Gamma(X, {\mathcal K}_X^\ast/{\mathcal O}_X^\ast)$ sur notre
système intégral noethérien $X$
est la somme $\Sigma v_Z(\tau)[Z]$ où $v_Z(\tau)$ est défini
par la recette suivante:
\begin{enumerate}
\item Si le germe de $\tau$ au point générique
$\xi$ de $Z$ est zéro -- en d'autres termes, l'image de
$\tau$ dans la tige $({\mathcal K}^\ast/{\mathcal O}^\ast)_\xi$ est ``zero'' -- alors $v_Z(\tau) = 0$.
\item Trouver un quartier ouvert d'affine $\Spec(A) = U \subset X$ de
sorte que $\tau|_U$ soit l'image d'une section $f \in {\mathcal K}(U)$ et
de plus $f = a/b$ avec $a, b \in A$.
$$
v_Z(f) = \text{length}_{{\mathcal O}_{X, Z}}({\mathcal O}_{X, Z}/(a)) -
\text{length}_{{\mathcal O}_{X, Z}}({\mathcal O}_{X, Z}/(b)).
$$
\end{enumerate}
\end{enumerate}
\end{definition}

\begin{remarks}
\label{remarks-divisors}
Voici quelques remarques insignifiantes.
\begin{enumerate}
\item Sur un schéma intégral de noethérien $X$,
la feuille ${\mathcal K}_X$ est constante avec la valeur du corps de fonction $K(X)$.
\item Pour avoir un sens à partir des définitions ci-dessus, il faut
montrer que
$$
\text{length}_{\mathcal O}({\mathcal O}/(ab)) =
\text{length}_{\mathcal O}({\mathcal O}/(a)) +
\text{length}_{\mathcal O}({\mathcal O}/(b))
$$
pour n'importe quelle paire $(a, b)$ d'éléments non nuls d'un domaine local noetherien
à une dimension ${\mathcal O}$. Ceci sera fait dans les conférences.
\end{enumerate}
\end{remarks}

\begin{exercise}
\label{exercise-effective-cartier-cartier}
(Sur n'importe quel
schéma.) Décrire comment attribuer un diviseur Cartier à un diviseur Cartier efficace.
\end{exercise}

\begin{exercise}
\label{exercise-rational-function-cartier}
(Sur un schéma intégral.)
Décrire comment attribuer un diviseur Cartier $D$ à une fonction rationnelle $f$
telle que le diviseur Weil associé à $D$ et à $f$ soit
d'accord. (C'est idiot.)
\end{exercise}

\begin{exercise}
\label{exercise-weil-not-cartier}
Donner un exemple d'un diviseur de Weil sur une variété qui n'est
pas le diviseur de Weil associé à un diviseur Cartier.
\end{exercise}

\begin{exercise}
\label{exercise-weil-Q-cartier}
Donner un exemple d'un diviseur de Weil $D$ sur une variété qui n'est
pas le diviseur de Weil associé à un diviseur Cartier, mais telle
que $nD$ est le diviseur de Weil associé à un diviseur Cartier
pour certains $n > 1$.
\end{exercise}

\begin{exercise}
\label{exercise-weil-not-Q-cartier}
Donner un exemple d'un diviseur de Weil $D$ sur une variété qui n'est pas
le diviseur de Weil associé à un diviseur Cartier et de telle sorte
que $nD$ n'est PAS le diviseur de Weil associé à un diviseur Cartier pour
n'importe quel $n > 1$.
(Considérez un cône, par exemple $X : xy - zw = 0$ dans
$\mathbf{A}^4_k$. Essayez de montrer que $D = [x = 0, z = 0]$ fonctionne.)
\end{exercise}

\begin{exercise}
\label{exercise-cartier-not-difference-effective-cartier}
Sur un schéma séparé $X$ de type fini sur
un corps: Donner un exemple d'un diviseur Cartier qui n'est pas
la différence de deux diviseurs
Cartier efficaces. Conseil: Trouver quelques $X$ qui n'ont pas de diviseurs Cartier
efficaces non vides par exemple le schéma
construit dans \cite[III Exercise 5.9]{H}. Il y a même
un exemple avec $X$ une variété --
à savoir la variété d'exercice \ref{exercise-nonprojective}.
\end{exercise}

\begin{exercise}
\label{exercise-nonprojective}
Exemple d'une bonne variété non-projective. Soit $k$
être un corps. Soit $L \subset \mathbf{P}^3_k$ être une ligne et laissez
$C \subset \mathbf{P}^3_k$ être un conique non-singulaire. Supposons que $C \cap L
= \emptyset$. Choisissez un isomorphisme
$\varphi : L \to C$. Soit $X$ être la variable $k$ obtenue
en collant $C$ à $L$ via $\varphi$. En d'autres termes, il
y a un morphisme birationnel approprié surjectif
$$
\pi : \mathbf{P}^3_k \longrightarrow X
$$
et un $U \subset X$ ouvert tel que $\pi : \pi^{-1}(U)
\to U$ est un isomorphisme, $\pi^{-1}(U) = \mathbf{P}^3_k \setminus (L \cup
C)$ et tel que $\pi|_L = \pi|_C \circ \varphi$. (Ces conditions ne définissent
pas encore uniquement $X$. Pour ce faire, vous devez spécifier la
gaine de la structure de $X$ le long de points de
$Z = X \setminus U$.) Afficher $X$ existe, est une
bonne variété, mais n'est pas projective.
(Conseil : Pour l'existence, utilisez le résultat de l'exercice \ref{exercise-prepare-glueing}. Pour une
utilisation non-projective, $\Pic(\mathbf{P}^3_k) = \mathbf{Z}$ pour montrer que $X$ ne peut
pas avoir de gaine inverse.)
\end{exercise}



\section{Differentiels}
\label{section-differentials}

\noindent
{\bf Définitions et résultats.} K\"Hler différentielles.
\begin{enumerate}
\item Laisser $R \to A$ être une carte annulaire. Le module
{\it des différentiels K\"ahler de $A$ sur $R$} est désigné $\Omega_{A/R}$.
Il est généré par les éléments $\text{d}a$, $a \in
A$ sous réserve des relations:
$$
\text{d}(a_1 + a_2) = \text{d}a_1 + \text{d}a_2,\quad
\text{d}(a_1a_2) = a_1\text{d}a_2 + a_2\text{d}a_1,\quad
\text{d}r = 0
$$
Les cartes universelles canoniques $R$-derivation $\text{d} : A \to
\Omega_{A/R}$ $a\mapsto \text{d}a$.
\item Considérez la courte séquence exacte
$$
0 \to I \to A \otimes_R A \to A \to 0
$$
qui définit l'idéal $I$. Il y a une dérivation canonique $\text{d} :
A \to I/I^2$ qui map $a$ à la classe de
$a \otimes 1 - 1 \otimes a$. Ceci est une autre présentation
du module de dérivation de $A$ sur $R$, en d'autres termes
$$
(I/I^2, \text{d}) \cong (\Omega_{A/R}, \text{d}).
$$
\item Pour les sous-ensembles multiplicatifs $S_R \subset R$ et $S_A
\subset A$ tels que $S_R$ nous avons des cartes dans $S_A$
$$
\Omega_{S_A^{-1}A / S_R^{-1}R} =
S_A^{-1}\Omega_{A/R}.
$$
\item Si $A$ est un modèle finiment présenté $R$-algèbre
alors $\Omega_{A/R}$ est un modèle finiment présenté $A$. Par conséquent,
dans ce cas, les idéaux {\it correspondant} de $\Omega_{A/R}$ sont définis.
\item Laisser $f : X \to S$ être un morphisme
des schémas. Il y a une gaine quasi-cohérente de ${\mathcal O}_X$-modules
$\Omega_{X/S}$ et une dérivation linéaire ${\mathcal O}_S$
$$
\text{d} : {\mathcal O}_X \longrightarrow \Omega_{X/S}
$$
de telle sorte que pour tout affine s'ouvre $\Spec(A) = U \subset
X$, $\Spec(R) = V \subset S$
avec $f(U) \subset V$ nous avons
$$
\Gamma(\Spec(A), \Omega_{X/S}) = \Omega_{A/R}
$$
compatible avec $\text{d}$.
\end{enumerate}

\begin{exercise}
\label{exercise-dual-numbers}
Soit $k[\epsilon]$ être l'anneau de nombres doubles sur
le corps $k$, c'est-à-dire $\epsilon^2 = 0$.
\begin{enumerate}
\item Considérez la carte de l'anneau
$$
R = k[\epsilon] \to A = k[x, \epsilon]/(\epsilon x)
$$
Montrer que les idéaux de montage de $\Omega_{A/R}$ sont (à partir de
l'idéal de montage zéro)
$$
(\epsilon), A, A, \ldots
$$
\item Considérez la carte $R = k[t]
\to A = k[x, y, t]/(x(y-t)(y-1), x(x-t))$. Montrer que les idéaux d'aménagement
de $\Omega_{A/R}$ dans $A$ sont (la caractéristique d'ensemble $k$ est zéro
pour la simplicité)
$$
x(2x-t)(2y-t-1)A, \ (x, y, t)\cap (x, y-1, t), \ A, \ A, \ldots
$$
Donc, les $0$- l'idéal de montage est découpé par un seul élément de
$A$, les $1$st L'idéal de montage définit deux points fermés de $\Spec(A)$, et
les autres sont tous triviaux.
\item Considérez la carte $R = k[t] \to A = k[x, y,
t]/(xy-t^n)$. Calculer les idéaux de montage de $\Omega_{A/R}$.
\end{enumerate}
\end{exercise}

\begin{remark}
\label{remark-fitting-omega-not-sings}
L'idéal $k$th Fitting de $\Omega_{X/S}$ est couramment utilisé pour définir le
schéma singulier du morphisme $X \to S$ lorsque $X$ a une dimension relative $k$
sur $S$. Mais comme le montre la partie a), vous devez faire attention lorsque
votre famille n'a pas de dimension fibreuse "constante", par exemple lorsqu'elle n'est pas plate.
Comme le montre la partie b), la planéité ne le garantit pas non plus
(et oui c'est une famille plate). Dans les "bons cas" -- comme dans c)
-- pour les familles de courbes, on s'attend à ce que l'idéal $0$-th
Fitting soit zéro et le $1$st Fitting idéal pour définir (scheme-théoriquement) le locus singulier.
\end{remark}

\begin{exercise}
\label{exercise-formally-smooth}
Supposons que $R$ soit un anneau et
$$
A = R[x_1, \ldots, x_n]/(f_1, \ldots, f_n).
$$
Notez que nous supposons que $A$ est
présenté par le même nombre d'équations que les
variables.
$$
( \partial f_i / \partial x_j )
$$
est $n \times n$, c'est-à-dire une matrice carrée. Supposons que
son déterminant est invertible en tant qu'élément dans $A$. Notez que
c'est exactement la condition qui dit que $\Omega_{A/R} = (0)$ dans ce
cas de $n$-générateurs et de relations $n$. Soit $\pi
: B' \to B$ être une surjection de $R$-algèbres dont
le noyau $J$ a un zéro carré (comme un idéal
dans $B'$). Soit $\varphi : A \to B$ être un homomorphisme
de $R$-algèbres. Montrer qu'il existe un homomorphisme unique de $R$-algèbres
$\varphi' : A \to B'$ tel que $\varphi = \pi \circ \varphi'$.
\end{exercise}

\begin{exercise}
\label{exercise-formally-smooth-one-equation}
Trouver une généralisation du résultat de l'exercice \ref{exercise-formally-smooth} dans le cas
où $A = R[x, y]/(f)$.
\end{exercise}

\begin{exercise}
\label{exercise-Jacobian-criterion}
Soit $k$ être un corps, laissez $f_1, \ldots, f_c \in k[x_1, \ldots, x_n]$
et laissez $A = k[x_1, \ldots, x_n]/(f_1, \ldots, f_c)$.
Supposons que $f_j(0, \ldots, 0) = 0$. Cela signifie que
$\mathfrak m = (x_1, \ldots, x_n)A$ est un idéal maximal.
Prouver que l'anneau local $A_\mathfrak m$ est régulier si
le rang de la matrice
$$
(\partial f_j/ \partial x_i)|_{(x_1, \ldots, x_n) = (0, \ldots, 0)}
$$
est $c$. Quelle est la dimension de $A_\mathfrak m$ dans ce
cas? Montrer que l'inverse est faux en donnant un exemple où
$A_\mathfrak m$ est régulier mais le rang est inférieur à $c$;
quelle est la dimension de $A_\mathfrak m$ dans votre exemple?
\end{exercise}




\section{Schemes, Examen final, automne 2007}
\label{section-final-exam-fall-2007}

\noindent
Il s'agissait des questions de l'examen final d'un cours sur les
programmes, au printemps 2007 à l'Université Columbia.

\begin{exercise}[Definitions]
\label{exercise-definitions}
Fournir des définitions des concepts suivants.
\begin{enumerate}
\item $X$ est un schéma {\it}
\item le morphisme des schémas $f : X \to Y$ est {\it fini}
\item les morphismes des schémas $f : X \to Y$ sont {\it de type fini}
\item le schéma $X$ est {\it noethérien}
\item le module ${\mathcal O}_X$ ${\mathcal L}$ du
schéma $X$ est {\it invertible}
\item le genre {\it} d'une courbe
de projection non singulaire sur un corps algébrique fermé
\end{enumerate}
\end{exercise}

\begin{exercise}
\label{exercise-kill-global-sections}
Soit $X = \Spec({\mathbf Z}[x, y])$, et laissez ${\mathcal F}$
être un
module quasi-cohérent ${\mathcal O}_X$. Supposons que ${\mathcal F}$ soit zéro lorsqu'il est limité
à
l'affine standard ouverte $D(x)$.
\begin{enumerate}
\item Montrer que chaque section globale $s$ de ${\mathcal F}$ est tuée par une certaine
puissance de $x$, c'est-à-dire $x^ns = 0$ pour certains $n\in {\mathbf N}$.
\item Pensez-vous qu'il en soit de même si nous ne supposons pas que ${\mathcal
F}$ est quasi-cohérent?
\end{enumerate}
\end{exercise}

\begin{exercise}
\label{exercise-empty-fibre-empty}
Supposez que $X \to \Spec(R)$ est un morphisme approprié et que
$R$ est un anneau d'évaluation discret avec le corps de résidus $k$.
Supposez que $X \times_{\Spec(R)} \Spec(k)$ est le schéma vide. Montrer que $X$
est le schéma vide.
\end{exercise}

\begin{exercise}
\label{exercise-curve-p1-p1}
Considérez la
variété projective \footnote{L'intégration projective est $((X_0,
X_1), (Y_0, Y_1))\mapsto (X_0Y_0, X_0Y_1, X_1Y_0, X_1Y_1)$ en d'autres
termes $(x, y)\mapsto (1, y, x, xy)$.}
$$
{\mathbf P}^1 \times {\mathbf P}^1 = {\mathbf P}^1_{{\mathbf C}}
\times_{\Spec({\mathbf C})} {\mathbf P}^1_{\mathbf C}
$$
sur le corps des nombres complexes ${\mathbf C}$. Il est couvert par quatre
pièces d'affine,
correspondant à des paires de pièces d'affine standard ${\mathbf P}^1_{\mathbf C}$. Par exemple,
supposons que nous
utilisons des coordonnées homogènes $X_0, X_1$ sur le premier facteur et $Y_0, Y_1$
Sur la seconde. $x = X_1/X_0$, et $y = Y_1/Y_0$. Ensuite, les 4
pièces d'affine ouvertes sont les spectres des anneaux
$$
{\mathbf C}[x, y], \quad
{\mathbf C}[x^{-1}, y], \quad
{\mathbf C}[x, y^{-1}], \quad
{\mathbf C}[x^{-1}, y^{-1}].
$$
Laisser $X \subset {\mathbf P}^1 \times {\mathbf P}^1$ être le sous-régime fermé
qui est
la fermeture du sous-ensemble fermé du premier morceau d'affine donné par l'équation
$$
y^3(x^4 + 1) = x^4 -1.
$$
\begin{enumerate}
\item Montrer que $X$ est contenu dans l'union du premier et
du dernier des 4 pièces d'affine ouvertes.
\item Afficher que $X$ est une courbe projective non singulaire.
\item Considérez le morphisme $pr_2 : X \to {\mathbf P}^1$ (projection sur le premier
facteur). Sur le premier morceau d'affine, c'est la carte $(x, y) \mapsto x$. Expliquer
brièvement pourquoi il a un degré $3$.
\item Calculer les points de ramification et les indices de ramification pour
la carte $pr_2 : X \to {\mathbf P}^1$.
\item Calculer le genre de $X$.
\end{enumerate}
\end{exercise}

\begin{exercise}
\label{exercise-finite-type-over-Z}
Soit $X \to \Spec({\mathbf Z})$ être un morphisme de type
fini. Supposons qu'il y ait un nombre infini de premiers
$p$ tels que $X \times_{\Spec({\mathbf Z})} \Spec({\mathbf F}_p)$ n'est pas vide.
\begin{enumerate}
\item Afficher que $X \times_{\Spec({\mathbf Z})}\Spec(\mathbf{Q})$ n'est
pas vide.
\item Pensez-vous qu'il en soit de même si nous remplaçons
la condition ``type fini'' par la condition ``type fini local''?
\end{enumerate}
\end{exercise}




\section{Schemes, Examen final, printemps 2009}
\label{section-final-exam-spring-2009}

\noindent
Il s'agissait des questions de l'examen final d'un cours sur les
programmes, au printemps 2009 à l'Université Columbia.

\begin{exercise}
\label{exercise-Noetherian-coherent}
Soit $X$ être un schéma noethérien.
Soit $\mathcal{F}$ être une gaine cohérente sur $X$.
Soit $x \in X$ être
un point. Supposons que $\text{Supp}(\mathcal{F}) = \{ x \}$.
\begin{enumerate}
\item Afficher que $x$ est un point fermé de $X$.
\item Afficher que $H^0(X, \mathcal{F})$ n'est pas zéro.
\item Afficher que $\mathcal{F}$ est généré par des sections globales.
\item Afficher que $H^p(X, \mathcal{F}) = 0$ pour $p > 0$.
\end{enumerate}
\end{exercise}

\begin{remark}
\label{remark-invertible-projective-space}
Soit $k$
être un corps. Soit $\mathbf{P}^2_k = \text{Proj}(k[X_0,
X_1, X_2])$. Toute gaine invertible sur $\mathbf{P}^2_k$
est isomorphe à $\mathcal{O}_{\mathbf{P}^2_k}(n)$ pour certains $n \in
\mathbf{Z}$.
$$
\Gamma(\mathbf{P}^2_k, \mathcal{O}_{\mathbf{P}^2_k}(n)) =
k[X_0, X_1, X_2]_n
$$
est la partie de degré
$n$ de l'anneau polynôme. Pour une gaine
quasi-cohérente $\mathcal{F}$
sur
$\mathbf{P}^2_k$ set $\mathcal{F}(n) =
\mathcal{F} \otimes_{\mathcal{O}_{\mathbf{P}^2_k}} \mathcal{O}_{\mathbf{P}^2_k}(n)$ comme
d'habitude.
\end{remark}

\begin{exercise}
\label{exercise-nonsplit-vectorbundle}
Soit $k$ être
un corps. Soit $\mathcal{E}$ être un faisceau vectoriel sur $\mathbf{P}^2_k$,
c'est-à-dire un module localement libre fini
$\mathcal{O}_{\mathbf{P}^2_k}$. Nous disons que $\mathcal{E}$ est {\it split} si
$\mathcal{E}$ est isomorphe à une somme directe inverse $\mathcal{O}_{\mathbf{P}^2_k}$-modules.
\begin{enumerate}
\item Afficher que $\mathcal{E}$ est fractionné si et seulement si $\mathcal{E}(n)$ est
fractionné.
\item Afficher que si $\mathcal{E}$ est fractionné,
alors $H^1({\mathbf{P}^2_k}, \mathcal{E}(n)) = 0$ pour
tous $n \in \mathbf{Z}$.
\item Laisser
$$
\varphi :
\mathcal{O}_{\mathbf{P}^2_k}
\longrightarrow
\mathcal{O}_{\mathbf{P}^2_k}(1) \oplus
\mathcal{O}_{\mathbf{P}^2_k}(1) \oplus
\mathcal{O}_{\mathbf{P}^2_k}(1)
$$
être donné par des
formes linéaires $L_0, L_1, L_2 \in \Gamma(\mathbf{P}^2_k, \mathcal{O}_{\mathbf{P}^2_k}(1))$. Assumer $L_i \not = 0$
pour certains $i$. Quelle est la
condition sur $L_0, L_1, L_2$ telle que le cokernel
de $\varphi$ est un faisceau de vecteurs?
Pourquoi?
\item Exemple d'un tel $\varphi$.
\item Afficher que $\Coker(\varphi)$ n'est pas divisé (si c'est
un paquet vectoriel).
\end{enumerate}
\end{exercise}

\begin{remark}
\label{remark-recall-dimension-theory}
Utilisez gratuitement les faits suivants sur la théorie des dimensions (et
ajoutez-en plus si vous en avez besoin).
\begin{enumerate}
\item La dimension d'un schéma est le supreme de la longueur des
chaînes de sous-ensembles fermés irréductibles.
\item La dimension d'un schéma de type fini sur un corps est le
maximum des dimensions de son affine s'ouvre.
\item La dimension d'un schéma noethérien est le maximum des dimensions
de ses composants irréductibles.
\item La dimension d'un schéma
d'affine coïncide avec la dimension de l'anneau correspondant.
\item Laisser $k$ être un corps et laisser $A$ être un type fini $k$-algèbre. Si
$A$ est un domaine, et $x \not = 0$, alors $\dim(A) = \dim(A/xA) + 1$.
\end{enumerate}
\end{remark}

\begin{exercise}
\label{exercise-irreducible-fibres-same-dimension-irreducible}
Soit $k$ être
un corps. Soit $X$ être un schéma projectif, réduit
sur $k$. Soit $f : X \to \mathbf{P}^1_k$ être un morphisme de
schémas sur $k$. Supposons qu'il existe un entier $d
\geq 0$ tel que pour chaque point $t \in \mathbf{P}^1_k$ la fibre
$X_t = f^{-1}(t)$ est irréductible de dimension $d$. (Rappelez qu'un espace irréductible n'est
pas vide.)
\begin{enumerate}
\item Afficher cette $\dim(X) = d + 1$.
\item Laisser $X_0 \subset X$ être une composante irréductible de $X$ dimension $d
+ 1$. Prouver que pour chaque $t \in \mathbf{P}^1_k$ la fibre $X_{0,
t}$ a une dimension $d$.
\item Que pouvez-vous conclure à propos de $X_t$ et $X_{0, t}$ à partir de ce qui précède?
\item Montrer que $X$ est irréductible.
\end{enumerate}
\end{exercise}

\begin{remark}
\label{remark-chi}
Compte tenu d'un schéma projectif $X$ sur un corps $k$ et
d'une feuille cohérente $\mathcal{F}$ sur $X$, nous avons défini
$$
\chi(X, \mathcal{F}) =
\sum\nolimits_{i \geq 0} (-1)^i\dim_k H^i(X, \mathcal{F}).
$$
\end{remark}

\begin{exercise}
\label{exercise-complete-intersection}
Soit $k$ être
un corps. Écrivez $\mathbf{P}^3_k = \text{Proj}(k[X_0, X_1, X_2, X_3])$. Soit
$C \subset \mathbf{P}^3_k$ être un {\it
type $(5, 6)$ courbe d'intersection complète}. Cela signifie
qu'il existe $F \in k[X_0, X_1, X_2, X_3]_5$ et $G
\in k[X_0, X_1, X_2, X_3]_6$ tel que
$$
C = \text{Proj}(k[X_0, X_1, X_2, X_3]/(F, G))
$$
est une variété de dimension $1$. (La variety implique une immersion
réduite et irréductible, mais n'hésitez pas à supposer
que $C$ est non-singulaire si vous le souhaitez.) Soit $i :
C \to \mathbf{P}^3_k$ être l'immersion fermée correspondante.
$$
\xymatrix{
0 \ar[r] &
\mathcal{O}_{\mathbf{P}^3_k}(-11)
\ar[r]^-{
\left(
\begin{matrix}
-G \\
F
\end{matrix}
\right)
} &
\mathcal{O}_{\mathbf{P}^3_k}(-5) \oplus \mathcal{O}_{\mathbf{P}^3_k}(-6)
\ar[r]^-{(F, G)} &
\mathcal{O}_{\mathbf{P}^3_k} \ar[r] &
i_*\mathcal{O}_C \ar[r] &
0
}
$$
est une séquence exacte de gerbes. S'il vous plaît utiliser ces faits pour:
\begin{enumerate}
\item calculer $\chi(C, i^*\mathcal{O}_{\mathbf{P}^3_k}(n))$ pour tout $n
\in \mathbf{Z}$, et
\item calculer la dimension de $H^1(C, \mathcal{O}_C)$.
\end{enumerate}
\end{exercise}

\begin{exercise}
\label{exercise-glueing}
Soit $k$ être un
corps. Considérez les anneaux
\begin{align*}
A & = k[x, y]/(xy) \\
B & = k[u, v]/(uv) \\
C & = k[t, t^{-1}] \times k[s, s^{-1}]
\end{align*}
et les cartes $k$-algèbres
$$
\begin{matrix}
A \longrightarrow C, &
x \mapsto (t, 0), &
y \mapsto (0, s) \\
B \longrightarrow C, &
u \mapsto (t^{-1}, 0), &
v \mapsto (0, s^{-1})
\end{matrix}
$$
Il est vrai que ces cartes induisent des isomorphismes $A_{x
+ y} \to C$ et $B_{u + v} \to C$. D'où les cartes
$A \to C$ et $B \to C$ identifient $\Spec(C)$ avec
des sous-ensembles ouverts de $\Spec(A)$ et $\Spec(B)$. Soit $X$ être le schéma
obtenu en collant $\Spec(A)$ et $\Spec(B)$ le long de $\Spec(C)$:
$$
X = \Spec(A) \amalg_{\Spec(C)} \Spec(B).
$$
Comme nous l'avons vu dans le cours, un tel schéma existe et il
y a des ouvertures d'affine $\Spec(A) \subset X$ et $\Spec(B) \subset
X$ dont le chevauchement est exactement $\Spec(C)$ identifié avec une ouverture de chacune
d'entre elles à l'aide des cartes ci-dessus.
\begin{enumerate}
\item Pourquoi $X$ est-il séparé?
\item Pourquoi $X$ de type fini sur $k$?
\item Calculer $H^1(X, \mathcal{O}_X)$, ou quelle est sa dimension?
\item Qu'est-ce qu'une façon plus géométrique de décrire $X$?
\end{enumerate}
\end{exercise}





\section{Schemes, examen final, automne 2010}
\label{section-final-exam-fall-2010}

\noindent
Ce sont les questions de l'examen final d'un cours sur les
programmes, à l'automne 2010 à l'Université Columbia.

\begin{exercise}[Definitions]
\label{exercise-definitions-fall-2010}
Fournir des définitions des concepts suivants.
\begin{enumerate}
\item un schéma séparé,
\item un morphisme quasi-compact des schémas,
\item un morphisme affine des schémas,
\item un sous-ensemble multiplicatif d'un anneau,
\item un régime noethérien,
\item une variété.
\end{enumerate}
\end{exercise}

\begin{exercise}
\label{exercise-prime-avoidance}
Évitement primaire.
\begin{enumerate}
\item Soit $A$ être un anneau. Soit $I \subset A$ être
un idéal et laissez $\mathfrak q_1$, $\mathfrak q_2$ être des idéaux premiers
tels que $I \not \subset \mathfrak q_i$. Montrer
que $I \not \subset \mathfrak q_1 \cup \mathfrak q_2$.
\item Qu'est-ce qu'une interprétation géométrique de (1)?
\item Laisser $X = \text{Proj}(S)$ pour un anneau classé $S$.
Soit $x_1, x_2 \in X$. Montrer qu'il existe un standard ouvert
$D_{+}(F)$ qui contient à la fois $x_1$ et $x_2$.
\end{enumerate}
\end{exercise}

\begin{exercise}
\label{exercise-compose-affine}
Pourquoi une composition de morphismes d'affines est-elle affine ?
\end{exercise}

\begin{exercise}[Examples]
\label{exercise-examples}
Donner des exemples de ce qui suit :
\begin{enumerate}
\item Un schéma de projection réductible sur un corps $k$.
\item Un schéma de 100 points.
\item Un morphisme non affine des schémas.
\end{enumerate}
\end{exercise}

\begin{exercise}
\label{exercise-chevalley-hilbert-nullstellensatz}
Chevalley est théorème et Hilbert Nullstellensatz.
\begin{enumerate}
\item Soit $\mathfrak p \subset \mathbf{Z}[x_1, \ldots, x_n]$ être
un idéal maximal. Que signifie le théorème de Chevalley
sur $\mathfrak p \cap \mathbf{Z}$?
\item À son tour, que signifie le Hilbert Nullstellensatz
sur $\kappa(\mathfrak p)$?
\end{enumerate}
\end{exercise}

\begin{exercise}
\label{exercise-P0}
Soit $A$ être un anneau. Soit $S = A[X]$ comme un
rangé $A$-algèbre où $X$ a un degré $1$. Montrer que $\text{Proj}(S) \cong \Spec(A)$ comme
des schémas sur $A$.
\end{exercise}

\begin{exercise}
\label{exercise-finite-is-projective}
Soit $A \to B$ être une homomorphisme d'anneaux finie.
Montrer que $\Spec(B)$ est un schéma de projection H sur $\Spec(A)$.
\end{exercise}

\begin{exercise}
\label{exercise-not-geometrically-irreducible}
Donner un exemple d'un schéma $X$ sur un corps $k$
tel que $X$ est irréductible et tel que pour une extension
finie $k'/k$ le changement de base $X_{k'} = X \times_{\Spec(k)}
\Spec(k')$ est connecté mais reductible.
\end{exercise}





\section{Schemes, Examen final, printemps 2011}
\label{section-final-exam-spring-2011}

\noindent
Ce sont les questions de l'examen final d'un cours sur les
programmes, au printemps 2011 à l'Université Columbia.

\begin{exercise}[Definitions]
\label{exercise-definitions-spring-2011}
Fournir des définitions des concepts en italique.
\begin{enumerate}
\item un schéma {\it séparé},
\item un morphisme {\it universellement fermé} des schémas,
\item {\it $A$ domine $B$} pour les anneaux locaux $A, B$ contenus dans
un corps commun,
\item la dimension {\it} d'un schéma $X$,
\item des {\it codimensions} d'un sous-régime fermé irréductible $Y$ d'un
régime $X$,
\end{enumerate}
\end{exercise}

\begin{exercise}[Results]
\label{exercise-results-spring-2011}
Indiquez quelque chose d'équivalent officiel au fait discuté dans
le cours.
\begin{enumerate}
\item Le critère de valeur de la pertinence pour un morphisme
$X \to Y$ des variétés par exemple.
\item La relation entre $\dim(X)$ et le corps de fonction $k(X)$ de
$X$ pour une variété $X$ sur un corps $k$.
\item Remplir le blanc : La catégorie des courbes projectives non-singulières
sur $k$ et des morphismes non-constants est anti-équivalente à $\ldots\ldots\ldots$.
\item Normalisation nulle.
\item critère jacobin.
\end{enumerate}
\end{exercise}

\begin{exercise}
\label{exercise-genus-plane-curve}
Soit $k$ être
un corps. Soit $F \in k[X_0, X_1, X_2]$ être une forme homogène
de degré $d$. Supposons que $C = V_{+}(F) \subset \mathbf{P}^2_k$ est une courbe
lisse sur $k$. Dénotez $i : C \to \mathbf{P}^2_k$ l'immersion fermée correspondante.
\begin{enumerate}
\item Afficher qu'il y a une courte séquence exacte
$$
0 \to \mathcal{O}_{\mathbf{P}^2_k}(-d)
\to \mathcal{O}_{\mathbf{P}^2_k} \to i_*\mathcal{O}_C \to 0
$$
des gerbes cohérentes sur $\mathbf{P}^2_k$: dites-moi ce que sont les cartes
et brièvement pourquoi c'est exact.
\item Conclure que $H^0(C, \mathcal{O}_C) = k$.
\item Calculer le genre de $C$.
\item Supposons maintenant que $P = (0 : 0 : 1)$ n'est pas sur $C$.
Prouver que $\pi : C \to \mathbf{P}^1_k$ donné par $(a_0 : a_1 : a_2) \mapsto (a_0 : a_1)$
a un degré $d$.
\item Assumer $k$ est algébriquement fermé, supposer que tous les indices de ramification
(les ``$e_i$') sont $1$ ou $2$, et supposer que la caractéristique de $k$
n'est pas égale à $2$. Combien de points
de ramification ont $\pi : C \to \mathbf{P}^1_k$?
\item En termes de $F$, que pensez-vous d'un ensemble d'équations de
l'ensemble de points de ramification de $\pi$?
\item Pouvez-vous deviner $K_C$?
\end{enumerate}
\end{exercise}

\begin{exercise}
\label{exercise-Pic-triangle}
Soit $k$ être un corps. Soit $X$ être un ``triangle'' au-dessus de $k$,
c'est-à-dire que vous obtenez $X$ en collant trois copies de $\mathbf{A}^1_k$ entre elles en
identifiant $0$ sur la première copie à $1$ sur la deuxième copie, $0$ sur la
deuxième copie à $1$ sur la troisième copie, et $0$ sur la troisième copie à
$1$ sur la première copie. Il s'avère que $X$ est isomorphe
à $\Spec(k[x, y]/(xy(x + y + 1)))$; n'hésitez pas à utiliser
ceci. Calculer le groupe Picard de $X$.
\end{exercise}

\begin{exercise}
\label{exercise-birational-morphism-curves-ample}
Soit $k$ être un corps. Soit $\pi : X \to Y$ être un
morphisme birationnel fini de courbes avec $X$ une courbe non-singulaire projective sur $k$.
Il découle du matériau dans le cours que $Y$ est une courbe correcte et que
$\pi$ est le morphisme de normalisation de $Y$. Nous avons
également vu dans le cours qu'il existe un dense ouvert $V \subset Y$ tel que
$U = \pi^{-1}(V)$ est un dense ouvert dans $X$ et $\pi : U \to
V$ est un isomorphisme.
\begin{enumerate}
\item Montrer qu'il existe un diviseur Cartier efficace $D \subset X$ tel
que $D \subset U$ et tel que $\mathcal{O}_X(D)$ est suffisant sur $X$.
\item Laisser $D$ être comme dans (1). Montrer que $E = \pi(D)$ est un diviseur Cartier
efficace sur $Y$.
\item Indiquer brièvement pourquoi
\begin{enumerate}
\item la carte $\mathcal{O}_Y \to \pi_*\mathcal{O}_X$ a un cokernel cohérent $Q$ qui
est supporté dans $Y \setminus V$, et
\item pour chaque $n$ il y a
une carte correspondante $\mathcal{O}_Y(nE) \to \pi_*\mathcal{O}_X(nD)$ dont
le cokernel est isomorphe à $Q$.
\end{enumerate}
\item Montrer
que $\dim_k H^0(X, \mathcal{O}_X(nD)) - \dim_k H^0(Y, \mathcal{O}_Y(nE))$ est délimité
(de quoi?) et concluez que la gaine inverse $\mathcal{O}_Y(nE)$
a beaucoup de sections pour les grandes $n$ (pourquoi?).
\end{enumerate}
\end{exercise}





\section{Schemes, examen final, automne 2011}
\label{section-final-exam-fall-2011}

\noindent
Il s'agissait des questions de l'examen final d'un
cours sur l'algèbre communautaire, à l'automne 2011 à l'Université Columbia.


\begin{exercise}[Definitions]
\label{exercise-definitions-fall-2011}
Fournir des définitions des concepts en italique.
\begin{enumerate}
\item un anneau {\it noethérien},
\item un schéma {\it noethérien},
\item un homomorphisme du cycle {\it fini},
\item un {\it fini} morphisme des schémas,
\item la dimension {\it} d'un anneau.
\end{enumerate}
\end{exercise}

\begin{exercise}[Results]
\label{exercise-results-fall-2011}
Indiquez quelque chose d'équivalent officiel au fait discuté dans
le cours.
\begin{enumerate}
\item Théorème principal de Zariski.
\item Normalisation nulle.
\item Théorème du reste chinois.
\item Monter pour les cartes à anneaux finies.
\end{enumerate}
\end{exercise}

\begin{exercise}
\label{exercise-dimension-of-ring}
Soit $(A, \mathfrak m, \kappa)$ être un anneau local noethérien dont le
corps de résidus n'a pas la caractéristique $2$. Supposons
que $\mathfrak m$ soit généré par trois éléments $x, y, z$ et
que $x^2 + y^2 + z^2 = 0$ dans $A$.
\begin{enumerate}
\item Quelles sont les valeurs possibles de $\dim(A)$?
\item Donner un exemple pour montrer que chaque valeur est possible.
\item Afficher que $A$ est un domaine si $\dim(A)
= 2$.
(Aperçu : voir $\bigoplus_{n \geq 0} \mathfrak m^n/\mathfrak m^{n + 1}$.)
\end{enumerate}
\end{exercise}

\begin{exercise}
\label{exercise-localization}
Soit $A$ être un
anneau. Soit $S \subset T \subset A$ être des sous-ensembles multiplicatifs.
Supposons que
$$
\{\mathfrak q \mid \mathfrak q \cap S = \emptyset\} =
\{\mathfrak q \mid \mathfrak q \cap T = \emptyset\}.
$$
Montrer que $S^{-1}A \to T^{-1}A$ est un isomorphisme.
\end{exercise}

\begin{exercise}
\label{exercise-locus-of-rank-1}
Laisser $k$ être un corps algébrique fermé.
$$
V_0 = \{ A \in \text{Mat}(3 \times 3, k) \mid \text{rank}(A) = 1\}
\subset \text{Mat}(3 \times 3, k) = k^9.
$$
\begin{enumerate}
\item Afficher que $V_0$ est l'ensemble de points
fermés d'un sous-ensemble (Zariski) fermé localement $V \subset \mathbf{A}^9_k$.
\item Est-ce que $V$ irréductible?
\item Qu'est-ce que $\dim(V)$?
\end{enumerate}
\end{exercise}

\begin{exercise}
\label{exercise-not-complete-intersection}
Prouver que l'idéal $(x^2, xy, y^2)$ dans $\mathbf{C}[x, y]$ ne peut
pas être généré par des éléments $2$.
\end{exercise}

\begin{exercise}
\label{exercise-finite-projection}
Soit $f \in \mathbf{C}[x, y]$ être un polynôme non
constant. Montrer que pour certains $\alpha, \beta \in
\mathbf{C}$ la carte $\mathbf{C}$-algèbre
$$
\mathbf{C}[t] \longrightarrow \mathbf{C}[x, y]/(f),\quad
t \longmapsto \alpha x + \beta y
$$
est fini.
\end{exercise}

\begin{exercise}
\label{exercise-union-of-two-affines}
Montrer que, compte tenu de plusieurs points $p_1, \ldots, p_n \in
\mathbf{C}^2$, le schéma $\mathbf{A}^2_\mathbf{C} \setminus \{p_1, \ldots, p_n\}$ est une
union de deux affine s'ouvre.
\end{exercise}

\begin{exercise}
\label{exercise-surjection-A1-P1}
Montrer qu'il existe un morphisme surjectif des
schémas $\mathbf{A}^1_\mathbf{C} \to \mathbf{P}^1_\mathbf{C}$. (Surjectif signifie
simplement surjectif sur des ensembles sous-jacents de points.)
\end{exercise}

\begin{exercise}
\label{exercise-almost-surjective}
Soit $k$ être un corps algébrique fermé. Soit $A \subset B$ être
une extension de domaines qui sont tous deux de type fini
$k$-algèbres. Prouver que l'image de $\Spec(B) \to \Spec(A)$
contient un sous-ensemble ouvert non vide de $\Spec(A)$ en utilisant les étapes suivantes:
\begin{enumerate}
\item Prouver-le si $A \to B$ est aussi fini.
\item Prouver-le si le corps de fraction de $B$ est une extension finie
du corps de fraction de $A$.
\item Réduisez l'énoncé au cas précédent.
\end{enumerate}
\end{exercise}





\section{Schemes, examen final, automne 2013}
\label{section-final-exam-fall-2013}

\noindent
Il s'agissait des questions de l'examen final d'un
cours sur l'algèbre communautaire, à l'automne 2013 à l'Université Columbia.

\begin{exercise}[Definitions]
\label{exercise-definitions-fall-2013}
Fournir des définitions des concepts en italique.
\begin{enumerate}
\item un {\it radical idéal} d'un anneau,
\item un homomorphisme du cycle {\it de type fini}
\item un différentiel {\it a la Weil},
\item un schéma {\it}.
\end{enumerate}
\end{exercise}

\begin{exercise}[Results]
\label{exercise-results-fall-2013}
Indiquez quelque chose d'équivalent officiel au fait discuté dans
le cours.
\begin{enumerate}
\item résultat sur les polynômes hilbert des modules classés.
\item dimension d'un anneau local noéthérien $(R, \mathfrak m)$ et
$\bigoplus_{n \geq 0} \mathfrak m^n/\mathfrak m^{n + 1}$.
\item Riemann-Roch.
\item Clifford est théorème.
\item Chevalley est théorème.
\end{enumerate}
\end{exercise}

\begin{exercise}
\label{exercise-surjective-after-localization}
Soit $A \to B$ être une carte annulaire. Soit $S \subset A$ être un sous-ensemble
multiplicatif. Supposons que $A \to B$ est de type fini et $S^{-1}A \to
S^{-1}B$ est surjectif. Montrer qu'il existe un $f \in S$ tel que $A_f \to
B_f$ est surjectif.
\end{exercise}

\begin{exercise}
\label{exercise-injective-local-ring-map-not-surjective}
Donner un exemple d'un homomorphisme local par injection $A \to B$
des cycles locaux, de sorte que $\Spec(B) \to \Spec(A)$ n'est pas surjectif.
\end{exercise}

\begin{situation}[Notation plane curve]
\label{situation-curve-in-the-plane}
Laisser $k$ être un corps algébrique fermé.
Laisser $F(X_0, X_1, X_2) \in k[X_0, X_1, X_2]$ être un polynôme
irréductible homogène de degré $d$.
$$
D = V(F) \subset \mathbf{P}^2
$$
soit la courbe du plan projectif donnée par la disparition de
$F$. Ensemble $x = X_1/X_0$ et $y
= X_2/X_0$ et $f(x, y) = X_0^{-d}F(X_0, X_1, X_2)
= F(1, x, y)$. Nous indiquons $K$ le corps fractionnaire du
domaine $k[x, y]/(f)$. Nous laissons $C$ soit la courbe
abstraite correspondant à $K$. Rappelez-vous (à partir des conférences) qu'il y a
une carte surjective $C \to D$ qui est bijectif sur le
locus non-singulaire de $D$ et un
isomorphisme si $D$ est non-singulaire. $f_x = \partial f/\partial x$ et $f_y
= \partial f/\partial y$. Enfin, nous indiquons $\omega = \text{d}x/f_y = -
\text{d}y/f_x$ l'élément de $\Omega_{K/k}$ discuté dans les conférences. $K_C$
le diviseur de zéros et de pôles de $\omega$C'est vrai.
\end{situation}

\begin{exercise}
\label{exercise-node-in-the-plane}
Dans la situation \ref{situation-curve-in-the-plane} supposer $d \geq 3$ et que
la courbe $D$ a exactement un point singulier, à savoir $P =
(1 : 0 : 0)$. Supposons plus loin que nous avons l'expansion
$$
f(x, y) = xy + h.o.t
$$
autour de $P = (0, 0)$. Puis $C$ a deux points $v$ et $w$ allongés
sur $P$ caractérisés par
$$
v(x) = 1, v(y) > 1
\quad\text{and}\quad
w(x) > 1, w(y) = 1
$$
\begin{enumerate}
\item Montrer que l'élément $\omega
= \text{d}x/f_y = - \text{d}y/f_x$ de $\Omega_{K/k}$ a un pôle
de premier ordre à la fois $v$ et $w$. (Le comportement
de $\omega$ aux points non-singulaires est comme discuté dans les conférences.)
\item Dans les conférences, nous avons montré que $\omega$ disparaît pour commander $d
- 3$ au diviseur $X_0 = 0$ ramené à $C$ sous la carte
$C \to D$. Combiné avec l'information de (1) quel est le degré du
diviseur de zéros et de pôles de $\omega$ sur $C$?
\item Quel est le genre de la courbe $C$?
\end{enumerate}
\end{exercise}

\begin{exercise}
\label{exercise-smooth-plane-curve-linear-system}
Dans la situation \ref{situation-curve-in-the-plane} supposons $d = 5$ et que la courbe
$C = D$ est non-singulaire. Dans les conférences, nous avons montré que le
genre de $C$ est $6$ et que le système linéaire $K_C$ est donné par
$$
L(K_C) = \{h\omega \mid h \in k[x, y],\ \deg(h) \leq 2\}
$$
où $\deg$ indique le degré total\footnote{Nous obtenons $\leq 2$ parce que $d - 3 =
5 - 3 = 2$.}. Soit $P_1, P_2, P_3, P_4, P_5 \in
D$ être deux fois des points distincts dans l'affine ouverte $X_0 \not = 0$. Nous indiquons $\sum
P_i = P_1 + P_2 + P_3 + P_4 + P_5$ le diviseur correspondant de $C$.
\begin{enumerate}
\item Décrire $L(K_C - \sum P_i)$ en termes de polynômes.
\item Quelles sont les possibilités pour $l(\sum P_i)$?
\end{enumerate}
\end{exercise}

\begin{exercise}
\label{exercise-rational-curve-high-degree}
Notez un $F$ comme dans Situation \ref{situation-curve-in-the-plane} avec $d = 100$ de
telle sorte que le genre de $C$ soit $0$.
\end{exercise}

\begin{exercise}
\label{exercise-high-degree-curve-quadratic-equation}
Soit $k$ être un corps algébrique fermé. Soit $K/k$
générer finement l'extension du corps du degré de transcendance $1$. Soit
$C$ être la courbe abstraite correspondant à $K$.
Soit $V \subset K$ être un $g^r_d$ et $\Phi : C
\to \mathbf{P}^r$ être le morphisme correspondant. Montrer que l'image de
$C$ est contenue dans un quadric\footnote{Un quadric est un degré $2$ hypersurface,
c'est-à-dire le zéro défini dans $\mathbf{P}^r$ d'un degré $2$ polynôme homogène.}
si $V$ est un système linéaire complet
et $d$ est suffisamment grand par rapport au genre $C$.
(Crédit extra : bon lié au degré requis.)
\end{exercise}

\begin{exercise}
\label{exercise-surjective-map-a2-p2}
Notation comme dans Situation \ref{situation-curve-in-the-plane}. Soit $U \subset \mathbf{P}^2_k$
être le sous-régime ouvert dont le complément
est $D$. Décrire la $k$-algèbre $A =
\mathcal{O}_{\mathbf{P}^2_k}(U)$. Donner une limite supérieure pour le nombre de générateurs
de $A$ comme étant une $k$-algèbre.
\end{exercise}



\section{Schemes, Examen final, printemps 2014}
\label{section-final-exam-spring-2014}

\noindent
Ce sont les questions de l'examen final d'un cours
sur les programmes, au printemps 2014 à l'Université Columbia.

\begin{exercise}[Definitions]
\label{exercise-definitions-spring-2014}
Soit $(X, \mathcal{O}_X)$ être un schéma. Fournissez des définitions
des concepts italisés.
\begin{enumerate}
\item l'anneau local {\it de $X$ à un point $x$},
\item une gaine {\it quasi-cohérente} de $\mathcal{O}_X$-modules,
\item a {\it cohérente} d'une teneur en poids de matières grasses égale ou supérieure à 50 % en
poids $\mathcal{O}_X$-modules (s'il vous plaît supposer $X$ est localement noethérien,
\item un {\it affine ouvert} de $X$,
\item a {\it morphisme fini des schémas} $X \to Y$.
\end{enumerate}
\end{exercise}

\begin{exercise}[Theorems]
\label{exercise-results-spring-2014}
Précisez précisément un fait non trivial discuté dans les conférences relatives à
chaque point.
\begin{enumerate}
\item sur l'invariance birationnelle du pluri-genre des variétés,
\item être un morphisme d'affine est une propriété locale,
\item la topologie d'une fibre théorétique schématique d'un morphisme, et
\item critère d'appréciation de la pertinence.
\end{enumerate}
\end{exercise}

\begin{exercise}
\label{exercise-miss-curve}
$X = \mathbf{A}^2_\mathbf{C}$ où $\mathbf{C}$ est le corps des nombres complexes. Une ligne
{\it} signifiera un sous-régime fermé de $X$ défini par une équation
linéaire $ax + by + c = 0$ pour certains $a, b, c \in
\mathbf{C}$ avec $(a, b) \not = (0, 0)$. Une courbe {\it}
signifiera un sous-régime fermé irréductible (donc non vide) $C \subset X$ de dimension $1$.
Un {\it quadrique} signifiera une courbe définie
par une équation quadratique $ax^2 + bxy + cy^2
+ dx + ey + f = 0$ pour certains $a, b,
c, d, e, f \in \mathbf{C}$ et $(a, b,
c) \not = (0, 0, 0)$.
\begin{enumerate}
\item Trouver une courbe $C$ telle que chaque ligne ait une intersection non vide avec $C$.
\item Trouver une courbe $C$ telle que chaque ligne et chaque quadrium ait une intersection
non vide avec $C$.
\item Afficher que pour chaque courbe $C$ il existe une autre courbe
telle que $C \cap C' = \emptyset$.
\end{enumerate}
\end{exercise}

\begin{exercise}
\label{exercise-normal-bundle-exceptional-curve}
Soit $k$ être un corps. Soit $b : X \to
\mathbf{A}^2_k$ être l'explosion du plan de l'affine dans l'origine. En
d'autres termes, si $\mathbf{A}^2_k = \Spec(k[x,
y])$, alors $X = \text{Proj}(\bigoplus_{n \geq 0} \mathfrak m^n)$
où $\mathfrak m = (x, y) \subset k[x,
y]$. Prouver les instructions suivantes
\begin{enumerate}
\item la fibre théorique schématique $E$ de $b$ sur
l'origine est isomorphe à $\mathbf{P}^1_k$,
\item $E$ est un diviseur Cartier efficace sur $X$,
\item la restriction de $\mathcal{O}_X(-E)$ à $E$ est
un faisceau de lignes de degré $1$.
\end{enumerate}
(Rappelez que $\mathcal{O}_X(-E)$ est la gaine idéale de $E$ dans $X$.)
\end{exercise}

\begin{exercise}
\label{exercise-surjective-map-affine-variety-projective-variety}
Soit $k$ être un corps. Soit $X$ être une variété projective sur
$k$. Montrer qu'il existe une variété d'affines $U$ sur $k$ et un morphisme
surjectif des variétés $U \to X$.
\end{exercise}

\begin{exercise}
\label{exercise-vandermonde}
Soit $k$ être un corps de caractéristiques $p > 0$ différent
de $2,3$. Considérez le sous-régime fermé $X$ de $\mathbf{P}^n_k$ défini par
$$
\sum\nolimits_{i = 0, \ldots, n} X_i = 0,\quad
\sum\nolimits_{i = 0, \ldots, n} X_i^2 = 0,\quad
\sum\nolimits_{i = 0, \ldots, n} X_i^3 = 0
$$
Pour quelles paires $(n, p)$ cette variété est-elle singulière?
\end{exercise}






\section{Algèbre commune, examen final, automne 2016}
\label{section-final-exam-fall-2016}

\noindent
Ce sont les questions de l'examen final d'un
cours sur l'algèbre commutative, à l'automne 2016 à l'Université Columbia.

\begin{exercise}[Definitions]
\label{exercise-definitions-fall-2016}
Soit $R$ être
un anneau. Fournir des définitions des concepts italisés.
\begin{enumerate}
\item l'anneau local {\it de $R$ à un premier $\mathfrak p$},
\item un module {\it fini} $R$,
\item un {\it présenté finiement} $R$-module,
\item $R$ est {\it régulier},
\item $R$ est {\it caténaire},
\item $R$ est {\it Cohen-Macaulay}.
\end{enumerate}
\end{exercise}

\begin{exercise}[Theorems]
\label{exercise-results-fall-2016}
Précisez précisément un fait non trivial discuté dans les conférences relatives à
chaque point.
\begin{enumerate}
\item anneaux réguliers,
\item les premiers modules associés de Cohen-Macaulay,
\item dimension d'un domaine de type fini sur un corps, et
\item Chevalley est théorème.
\end{enumerate}
\end{exercise}

\begin{exercise}
\label{exercise-finite-injective}
Soit $A \to B$ être une homomorphisme d'anneaux telle que
\begin{enumerate}
\item $A$ est local avec un idéal maximal $\mathfrak m$,
\item $A \to B$ est une homomorphisme d'anneaux finie\footnote{Rappelons que cela signifie
que $B$ est finie sous forme de $A$-module.}
\item $A \to B$ est injectable (on pense à $A$ comme un sous-anneau de $B$).
\end{enumerate}
Montrer qu'il y a un premier idéal $\mathfrak q \subset B$ avec
$\mathfrak m = A \cap \mathfrak q$.
\end{exercise}

\begin{exercise}
\label{exercise-find-components}
Soit $k$ être un corps. Soit $R =
k[x, y, z, w]$. Considérez l'idéal $I =
(xy, xz, xw)$. Quelles sont les composantes
irréductibles de $V(I) \subset \Spec(R)$ et quelles sont leurs dimensions?
\end{exercise}

\begin{exercise}
\label{exercise-no-nonconstant-morphism}
Soit $k$ être un corps. Soit $A = k[x, x^{-1}]$ et $B =
k[y]$. Affichez que toute carte $k$-algèbre $\varphi : A \to B$ map $x$ à une constante.
\end{exercise}

\begin{exercise}
\label{exercise-regular-over-Z}
Considérez l'anneau $R = \mathbf{Z}[x, y]/(xy - 7)$. Prouver
que $R$ est régulier.
\end{exercise}

\noindent
Avec un anneau local noethérien $(R, \mathfrak m, \kappa)$
pour $n \geq
0$, nous laissons $\varphi_R(n) = \dim_\kappa(\mathfrak m^n/\mathfrak m^{n + 1})$.

\begin{exercise}
\label{exercise-hilbert-function-allowed}
Existe-t-il un anneau local noethérien $R$ avec $\varphi_R(n) =
n + 1$ pour tous $n \geq 0$?
\end{exercise}

\begin{exercise}
\label{exercise-hilbert-function-prohibited}
Soit $R$ être un anneau local
noethérien. Supposez que $\varphi_R(0) = 1$, $\varphi_R(1) =
3$, $\varphi_R(2) = 5$. Montrer que $\varphi_R(3) \leq 7$.
\end{exercise}





\section{Schemes, Examen final, printemps 2017}
\label{section-final-exam-spring-2017}

\noindent
Ce sont les questions de l'examen final d'un cours sur les
régimes, au printemps 2017 à l'Université Columbia.

\begin{exercise}[Definitions]
\label{exercise-definitions-spring-2017}
Soit $f : X \to Y$ être un
morphisme des schémas. Fournir de brèves définitions des concepts italisés.
\begin{enumerate}
\item la fibre théorique du schéma {\it } de $f$ à $y \in Y$,
\item $f$ est un {\it Morphisme fini},
\item un {\it quasi-cohérent} $\mathcal{O}_X$-module,
\item $X$ est {\it variété},
\item $f$ est un {\it morphisme lisse},
\item $f$ est un {\it morphisme approprié}.
\end{enumerate}
\end{exercise}

\begin{exercise}[Theorems]
\label{exercise-results-spring-2017}
Précisément mais brièvement, énoncez un fait non trivial discuté dans les conférences
relatives à chaque point.
\begin{enumerate}
\item Pousser les gerbes quasi-cohérentes,
\item cohomologie des gerbes cohérentes sur les variétés projectives,
\item Serre dualité pour un schéma projectif sur un corps, et
\item Riemann-Hurwitz.
\end{enumerate}
\end{exercise}

\begin{exercise}
\label{exercise-compute-degree}
Soit $k$ être un corps algébrique fermé.
Soit $\ell > 100$ être un nombre primaire différent
de la caractéristique de $k$. Soit $X$ être le modèle projectif non-singulier
de la courbe affine donnée par l'équation
$$
y^\ell = x(x - 1)^3
$$
dans $\mathbf{A}^2_k$. Répondre aux questions suivantes:
\begin{enumerate}
\item Quel est le genre de $X$?
\item Donner une limite supérieure
pour la gonalité\footnote{La gonalité est le plus petit degré de morphisme
non constant de $X$ à $\mathbf{P}^1_k$.} de $X$.
\end{enumerate}
\end{exercise}

\begin{exercise}
\label{exercise-count-components}
Soit $k$ être un corps algébrique fermé. Soit $X$ être
un schéma projectif réduit sur $k$ dont tous les composants irréductibles
ont la même dimension $1$. Soit $\omega_{X/k}$ être le
module de dualisation relatif. Montrer que
si $\dim_k H^1(X, \omega_{X/k}) > 1$, alors $X$ est déconnecté.
\end{exercise}

\begin{exercise}
\label{exercise-sections-L-L-inverse}
Donner un exemple d'un schéma $X$ et d'un module
$\mathcal{O}_X$ non trivial invertible $\mathcal{L}$ de telle sorte
que $H^0(X, \mathcal{L})$ et $H^0(X, \mathcal{L}^{\otimes -1})$ ne soient
pas nuls.
\end{exercise}

\begin{exercise}
\label{exercise-in-product}
Laisser $k$ être un corps algébrique fermé. Laisser $g \geq 3$. Laisser
$X$ et $X'$ être des courbes projectives lisses sur $k$ du genre
$g$ et $g + 1$. Laisser $Y \subset X \times X'$
être une courbe telle que les projections $Y \to X$ et $Y \to
X'$ ne soient pas constantes. Prouver que le modèle projectif non-singulaire de $Y$
a le genre $\geq 2g + 1$.
\end{exercise}

\begin{exercise}
\label{exercise-finitely-many}
Soit $k$ être un corps fini. Soit $g > 1$. Sketch une preuve de
ce qui suit: il n'y a qu'un nombre fini de
classes d'isomorphisme de courbes projectives lisses sur $k$ du
genre $g$. (Vous aurez le mérite d'essayer de répondre à cela.)
\end{exercise}







\section{Algèbre commune, examen final, automne 2017}
\label{section-final-exam-fall-2017}

\noindent
Ce sont les questions de l'examen final d'un cours sur l'algèbre commutative,
à l'automne 2017 à l'Université Columbia.

\begin{exercise}[Definitions]
\label{exercise-definitions-fall-2017}
Fournir de brèves définitions des concepts italisés.
\begin{enumerate}
\item des {\it adjoint gauche} d'un functeur $F : \mathcal{A} \to \mathcal{B}$,
\item des {\it degré de transcendance} d'une extension $L/K$ des champs,
\item une fonction régulière {\it} sur une variété d'affines classiques
$X \subset k^n$,
\item une gaine {\it} sur un espace topologique,
\item a {\it bague locale}, et
\item un morphisme des schémas $f : X \to Y$ étant {\it affine}.
\end{enumerate}
\end{exercise}

\begin{exercise}[Theorems]
\label{exercise-results-fall-2017}
Précisément mais brièvement énoncez un fait non trivial discuté dans les conférences relatives à
chaque élément (s'il y en a plus d'un, il suffit de choisir
l'un d'eux).
\begin{enumerate}
\item Yoneda lemma,
\item Mayer-Vietoris,
\item dimension et cohomologie,
\item Hilbert polynôme, et
\item dualité pour l'espace projectif.
\end{enumerate}
\end{exercise}

\begin{exercise}
\label{exercise-compute-dimension}
Laisse-moi. $k$ être un corps algébrique fermé. Considérer le sous-ensemble fermé $X$
des $k^5$ avec topologie et coordonnées de Zariski $x_1, x_2, x_3, x_4, x_5$
données par les équations
$$
x_1^2 - x_4 = 0,\quad
x_2^5 - x_5 = 0,\quad
x_3^2 + x_3 + x_4 + x_5 = 0
$$
Quelle est la dimension de $X$ et pourquoi?
\end{exercise}

\begin{exercise}
\label{exercise-can-there-be}
Soit $k$ être un corps. Soit $X = \mathbf{P}^1_k$
être l'espace projectif de la dimension $1$ sur $k$. Soit
$\mathcal{E}$ être un module fini localement libre
$\mathcal{O}_X$. Pour $d \in \mathbf{Z}$ indiquer $\mathcal{E}(d) = \mathcal{E} \otimes_{\mathcal{O}_X}
\mathcal{O}_X(d)$ la $d$e torsion Serre de
$\mathcal{E}$ et $h^i(X, \mathcal{E}(d)) = \dim_k H^i(X, \mathcal{E}(d))$.
\begin{enumerate}
\item Pourquoi n'y a-t-il pas $\mathcal{E}$
avec $h^0(X, \mathcal{E}) = 5$ et $h^0(X, \mathcal{E}(1)) = 4$?
\item Pourquoi n'y a-t-il pas $\mathcal{E}$
avec $h^1(X, \mathcal{E}(1)) = 5$ et $h^1(X, \mathcal{E}) = 4$?
\item Pour lequel $a \in \mathbf{Z}$ peut exister un paquet vectoriel
$\mathcal{E}$ sur $X$ avec
$$
\begin{matrix}
h^0(X, \mathcal{E})\phantom{(1)} = 1 &
h^1(X, \mathcal{E})\phantom{(1)} = 1 \\
h^0(X, \mathcal{E}(1)) = 2 &
h^1(X, \mathcal{E}(1)) = 0 \\
h^0(X, \mathcal{E}(2)) = 4 &
h^1(X, \mathcal{E}(2)) = a
\end{matrix}
$$
\end{enumerate}
Des réponses partielles sont bienvenues et encouragées.
\end{exercise}

\begin{exercise}
\label{exercise-banana}
Laisser $X$ être un espace topologique qui est l'union $X =
Y \cup Z$ de deux sous-ensembles fermés
$Y$ et $Z$ dont l'intersection est indiquée $W = Y \cap Z$. Dénote $i :
Y \to X$, $j : Z \to X$ et $k : W \to X$ les cartes
d'inclusion.
\begin{enumerate}
\item Montrer qu'il y a une courte séquence exacte de gerbes
$$
0 \to \underline{\mathbf{Z}}_X \to
i_*(\underline{\mathbf{Z}}_Y) \oplus
j_*(\underline{\mathbf{Z}}_Z) \to
k_*(\underline{\mathbf{Z}}_W) \to 0
$$
où $\underline{\mathbf{Z}}_X$ indique la masse constante avec la
valeur $\mathbf{Z}$ sur $X$, etc.
\item Que pouvez-vous conclure sur la relation entre les
groupes de cohomologie de $X$, $Y$, $Z$, $W$ et $\mathbf{Z}$-coefficients?
\end{enumerate}
\end{exercise}

\begin{exercise}
\label{exercise-cohomology-infinite-punctured}
Soit $k$ être un corps. Soit $A = k[x_1, x_2, x_3, \ldots]$ être
l'anneau polynôme dans infiniment de variables.
Dénotez $\mathfrak m$ l'idéal maximal de $A$ généré par toutes
les variables. Soit $X = \Spec(A)$ et $U = X \setminus \{\mathfrak m\}$.
\begin{enumerate}
\item Afficher $H^1(U, \mathcal{O}_U) = 0$. Conseil: {\v
C}ech cohomology computing.
\item Quelle est votre hypothèse pour $H^i(U, \mathcal{O}_U)$ pour $i \geq 1$?
\end{enumerate}
\end{exercise}

\begin{exercise}
\label{exercise-principal}
Soit $A$ être un anneau local. Soit $a \in A$ être
un non-zérodiviseur. Soit $I, J \subset A$ être des idéaux tels que
$IJ = (a)$. Montrer que l'idéal $I$ est principal, c'est-à-dire généré
par un élément (qui se révélera être un non-zérodiviseur).
\end{exercise}



\section{Schemes, Examen final, printemps 2018}
\label{section-final-exam-spring-2018}

\noindent
Ce sont les questions de l'examen final d'un cours sur les
régimes, au printemps 2018 à l'Université Columbia.

\begin{exercise}[Definitions]
\label{exercise-definitions-spring-2018}
Fournir de brèves définitions des concepts italisés. Laisser $k$
être un corps algébrique fermé. Laisser $X$
être une courbe projective sur $k$.
\begin{enumerate}
\item une algèbre lisse } { \it sur $k$,
\item le degré {\it} d'un module inverse $\mathcal{O}_X$ sur $X$,
\item le genre {\it} de $X$,
\item des {\it Groupe de classes de division de Weil} des $X$,
\item $X$ est {\it hyperelliptique}, et
\item des {\it numéro d'intersection} de
deux courbes sur une surface projective lisse sur $k$C'est vrai.
\end{enumerate}
\end{exercise}

\begin{exercise}[Theorems]
\label{exercise-results-spring-2018}
Précisément mais brièvement énoncez un fait non trivial discuté dans les conférences relatives à
chaque élément (s'il y en a plus d'un, il suffit de choisir
l'un d'eux).
\begin{enumerate}
\item Riemann-Hurwitz théorème,
\item Le théorème de Clifford,
\item factorisation des cartes entre surfaces projectives lisses,
\item Théorème de l'indice Hodge, et
\item Hypothèse de Riemann pour les courbes sur les champs finis.
\end{enumerate}
\end{exercise}

\begin{exercise}
\label{exercise-hyperelliptic}
Laisser $k$ être un corps algébrique fermé. Laisser $X \subset \mathbf{P}^3_k$ être une
courbe lisse de degré $d$ et genre $\geq 2$. Assumer $X$ n'est
pas contenu dans un plan et qu'il y
a une ligne $\ell$ dans $\mathbf{P}^3_k$ rencontrant $X$ dans $d - 2$
points. Montrer que $X$ est hyperelliptique.
\end{exercise}

\begin{exercise}
\label{exercise-singular}
Soit $k$ être un corps algébrique fermé. Soit $X$ être une
courbe projective avec des points singuliers distincts par paire $p_1, \ldots,
p_n$. Expliquer pourquoi le genre de normalisation de $X$ est au
plus $-n + \dim_k H^1(X, \mathcal{O}_X)$.
\end{exercise}

\begin{exercise}
\label{exercise-how-many-blowups}
Soit $k$ être un corps. Soit $X = \Spec(k[x, y])$ être affine $2$ espace.
Soit
$$
I = (x^3, x^2y, xy^2, y^3) \subset k[x, y].
$$
Que $Y \subset X$ soit le sous-régime fermé correspondant à $I$.
Que $b : X' \to X$ soit l'explosion de l'idéal $(x, y)$,
c'est-à-dire l'explosion de l'espace affine à l'origine.
\begin{enumerate}
\item Montrer que l'image inverse théorique du schéma $b^{-1}Y \subset X'$
est un diviseur Cartier efficace.
\item Étant donné un exemple d'idéal $J \subset k[x,
y]$ avec $I \subset J \subset (x, y)$ tel que si
$Z \subset X$ est le sous-régime fermé correspondant à $J$, alors l'image
inverse théorique du schéma $b^{-1}Z$ n'est pas un diviseur Cartier efficace.
\end{enumerate}
\end{exercise}

\begin{exercise}
\label{exercise-rational-surface}
Soit $k$ être un corps algébrique fermé. Considérez les types
de surfaces suivants
\begin{enumerate}
\item $S = C_1 \times C_2$ où $C_1$ et $C_2$ sont des courbes projectives lisses,
\item $S = C_1 \times C_2$ où $C_1$ et $C_2$ sont des courbes projectives lisses et
le genre $C_1$ est $> 0$,
\item $S \subset \mathbf{P}^3_k$ est une hypersurface de degré $4$, et
\item $S \subset \mathbf{P}^3_k$ est une hypersurface lisse de degré $4$.
\end{enumerate}
Pour chaque type, indiquez brièvement pourquoi ou pourquoi la classe de surfaces
de ce type contient des surfaces rationnelles.
\end{exercise}

\begin{exercise}
\label{exercise-self-square-line}
Laisser $k$ être un corps algébrique fermé. Laisser $S \subset \mathbf{P}^3_k$ être une
hypersurface lisse de degré $d$. Supposons que $S$ contient une ligne
$\ell$. Quel est le carré de $\ell$ considéré comme un diviseur sur $S$?
\end{exercise}






\section{Algèbre commune, examen final, automne 2019}
\label{section-final-exam-fall-2019}

\noindent
Ce sont les questions de l'examen final d'un cours sur l'algèbre commutative,
à l'automne 2019 à l'Université Columbia.

\begin{exercise}[Definitions]
\label{exercise-definitions-fall-2019}
Fournir de brèves définitions des concepts italisés.
\begin{enumerate}
\item a {\it sous-ensemble constructible} d'un espace topologique noétherien,
\item la localisation {\it} d'un module $R$ $M$ à un premier $\mathfrak p$,
\item des {\it longueur} d'un module sur une bague locale
noethérien $(A, \mathfrak m, \kappa)$,
\item un module projectif {\it} sur un anneau $R$, et
\item un module {\it Cohen-Macaulay} sur
un anneau local noethérien $(A, \mathfrak m, \kappa)$.
\end{enumerate}
\end{exercise}

\begin{exercise}[Theorems]
\label{exercise-results-fall-2019}
Précisément mais brièvement énoncez un fait non trivial discuté dans les conférences relatives à
chaque élément (s'il y en a plus d'un, il suffit de choisir
l'un d'eux).
\begin{enumerate}
\item images de jeux constructibles,
\item Hilbert Nullstellensatz,
\item dimension des algèbres de type fini sur les champs,
\item Normalisation de la noéther, et
\item anneaux locaux réguliers.
\end{enumerate}
\end{exercise}

\noindent
Pour un anneau $R$ et un idéal $I \subset R$ rappeler
que $V(I)$ indique l'ensemble de $\mathfrak p \in \Spec(R)$ avec $I \subset \mathfrak p$.

\begin{exercise}[Making primes]
\label{exercise-infinitely-many-primes}
Construire infiniment plusieurs idéaux principaux distincts $\mathfrak p
\subset \mathbf{C}[x, y]$ tels que $V(\mathfrak
p)$ contient $(x, y)$ et $(x - 1, y - 1)$.
\end{exercise}

\begin{exercise}[No prime]
\label{exercise-no-prime}
Laisser $R = \mathbf{C}[x, y, z]/(xy)$. Argue brièvement il n'y a pas un premier idéal
$\mathfrak p \subset R$ tel que $V(\mathfrak p)$ contient
$(x, y - 1, z - 5)$ et $(x - 1, y, z - 7)$.
\end{exercise}

\begin{exercise}[Frobenius]
\label{exercise-frobenius}
Soit $p$ être un nombre principal (vous pouvez supposer $p = 2$ pour simplifier les formules). Soit $R$
être un anneau tel que $p = 0$ dans $R$.
\begin{enumerate}
\item Afficher que la carte $F : R \to R$, $x \mapsto x^p$ est un homomorphisme cyclique.
\item Afficher que $\Spec(F) : \Spec(R) \to \Spec(R)$ est la carte d'identité.
\end{enumerate}
\end{exercise}

\noindent
Rappelons qu'une spécialisation {\it} $x \leadsto y$ de points d'un espace topologique signifie simplement
$y$ est à la fermeture de $x$. Nous disons $x \leadsto y$ est
une spécialisation {\it immédiate} s'il n'y a pas de $z$ différente de $x$
et $y$ telle que $x \leadsto z$ et $z \leadsto y$.

\begin{exercise}[Dimension]
\label{exercise-dimension}
Supposons que nous ayons un espace topologique sobre $X$ contenant $5$ points distincts
$x, y, z, u, v$ ayant les spécialisations suivantes
$$
\xymatrix{
x \ar[d] \ar[r] & u & v \ar[l] \ar[dl] \\
y \ar[r] & z
}
$$
Quelle est la dimension minimale qu'un $X$ peut avoir? Si $X$ est le
spectre d'une algèbre de type fini sur un corps et $x
\leadsto u$ est une spécialisation immédiate, que pouvez-vous dire de
la spécialisation $v \leadsto z$?
\end{exercise}

\begin{exercise}[Tor computation]
\label{exercise-tor-computation}
Laisser $R = \mathbf{C}[x, y, z]$. Laisser $M = R/(x, z)$ et $N =
R/(y, z)$. Pour laquelle $i \in \mathbf{Z}$ est $\text{Tor}_i^R(M, N)$ non zéro?
\end{exercise}

\begin{exercise}
\label{exercise-depth-goes-up}
Laisser $A \to B$ être un homomorphisme local plat des anneaux Noetheriens locaux. Montrer que
si $A$ a la profondeur $k$, alors $B$ a la profondeur au moins $k$.
\end{exercise}









\section{ Géométrie algébrique, examen final, printemps 2020}
\label{section-final-exam-spring-2020}

\noindent
Il s'agissait des questions de l'examen final d'un cours sur la géométrie
algébrique, au printemps 2020 à l'Université Columbia.

\begin{exercise}[Definitions]
\label{exercise-definitions-spring-2020}
Fournir de brèves définitions des concepts italisés.
\begin{enumerate}
\item un schéma {\it},
\item un morphisme {\it des schémas},
\item un module {\it quasi-cohérent} sur un schéma,
\item une variété {\it} sur un corps $k$,
\item une courbe {\it} sur un corps $k$,
\item a {\it morphisme fini} des schémas,
\item des {\it cohomologie} d'une gerbe de groupes abeliens $\mathcal{F}$ sur
un espace topologique $X$,
\item un {\it dualisant la gaine} sur un schéma $X$ de dimension $d$
propre sur un corps $k$, et
\item a {\it carte rationnelle} d'une variété $X$ à une variété $Y$C'est vrai.
\end{enumerate}
\end{exercise}

\begin{exercise}[Theorems]
\label{exercise-results-spring-2020}
Précisément mais brièvement énoncez un fait non trivial discuté dans les conférences relatives à
chaque élément (s'il y en a plus d'un, il suffit de choisir
l'un d'eux).
\begin{enumerate}
\item cohomologie des gerbes abéliennes sur un espace topologique noétherien $X$
de dimension $d$,
\item gerbe de différentiels $\Omega^1_{X/k}$ d'une variété lisse sur un
corps $k$,
\item dualisant la gerbe $\omega_X$ d'une variété projective lisse $X$ sur
le corps $k$,
\item une courbe lisse
du genre $0$ sur un corps à fermeture algébrique $k$, et
\item le genre d'une courbe plane de degré $d$.
\end{enumerate}
\end{exercise}

\begin{exercise}
\label{exercise-coh-P1-infty-doubled}
Soit $k$ être un corps. Soit $X$ être un schéma sur $k$. Assumer
$X = X_1 \cup X_2$ est une couverture ouverte avec $X_1$, $X_2$
à la fois isomorphe à $\mathbf{P}^1_k$ et $X_1 \cap X_2$ isomorphe à
$\mathbf{A}^1_k$. (Ce schéma existe, par exemple, vous pouvez prendre $\mathbf{P}^1_k$ avec
$\infty$ doublé.) Montrer que $\dim_k H^1(X, \mathcal{O}_X)$ est infini.
\end{exercise}

\begin{exercise}
\label{exercise-no-morphism}
Laisser $k$ être un corps algébrique fermé. Laisser $Y$ être une courbe projective lisse
du genre $10$. Trouver une bonne limite inférieure pour le
genre d'une courbe projective lisse
$X$ de sorte qu'il existe un morphisme non constant $f : X \to Y$
qui n'est pas un isomorphisme.
\end{exercise}

\begin{exercise}
\label{exercise-element-order-n-in-pic}
Laisser $k$ être un corps algébrique de caractéristiques $0$.

$$
X : T_0^d + T_1^d - T_2^d = 0 \subset \mathbf{P}^2_k
$$
soit la courbe de degré de Fermat $d \geq 3$. Considérons les points fermés $p = [1 : 0
: 1]$ et $q = [0 : 1 : 1]$ le $X$. Ensemble $D = [p] - [q]$C'est vrai.
\begin{enumerate}
\item Afficher que $D$ n'est pas trivial dans le groupe de classe de diviseur de Weil.
\item Montrer que $d D$ est trivial dans le groupe de classe des diviseurs de
Weil. (Enfin, essayez de montrer que $d[p]$ et $d[q]$ sont tous les deux l'intersection de
$X$ avec une ligne dans le plan.)
\end{enumerate}
\end{exercise}

\begin{exercise}
\label{exercise-number-equations}
Soit $k$ être un corps algébrique fermé. Considérez l'intégration $2$-uple
$$
\varphi : \mathbf{P}^2 \longrightarrow \mathbf{P}^5
$$
En termes de matériel/notation dans les conférences c'est le morphisme
$$
\varphi = \varphi_{\mathcal{O}_{\mathbf{P}^2}(2)} :
\mathbf{P}^2
\longrightarrow
\mathbf{P}(\Gamma(\mathbf{P}^2, \mathcal{O}_{\mathbf{P}^2}(2)))
$$
En termes de coordonnées homogènes, elle est donnée par
$$
[a_0 : a_1 : a_2] \longmapsto
[a_0^2 : a_0a_1 : a_0a_2 : a_1^2 : a_1a_2 : a_2^2]
$$
Il s'agit d'une immersion fermée (il suffit
d'utiliser ceci). Soit $I \subset k[T_0, \ldots, T_5]$ être l'idéal
homogène de $\varphi(\mathbf{P}^2)$, c'est-à-dire que les éléments de
la partie homogène $I_d$
sont les polynômes homogènes $F(T_0, \ldots, T_5)$ de degré $d$
qui se limitent à zéro sur le sous-schème fermé $\varphi(\mathbf{P}^2)$. Calculer
$\dim_k I_d$ en fonction de $d$.
\end{exercise}

\begin{exercise}
\label{exercise-dualizing-sheaf-rank-1}
Soit $k$ être un corps algébrique fermé. Soit $X$ être un schéma
approprié de dimension $d$ sur $k$ avec module de
dualisation $\omega_X$. Vous recevez les informations suivantes:
\begin{enumerate}
\item $\text{Ext}^i_X(\mathcal{F}, \omega_X) \times H^{d - i}(X,
\mathcal{F}) \to H^d(X, \omega_X) \xrightarrow{t} k$ est non dégénéré
pour tous $i$ et pour tous les modules cohérents $\mathcal{O}_X$
$\mathcal{F}$, et
\item $\omega_X$ est fini localement sans un certain rang $r$.
\end{enumerate}
Affichez que $r = 1$. (Aperçu : voir ce qui se passe si vous
prenez $\mathcal{F}$ un module approprié supporté à un point fermé.)
\end{exercise}







\section{Algèbre commutative, examen final, automne 2021}
\label{section-final-exam-fall-2021}

\noindent
Il s'agissait des questions de l'examen final d'un cours sur l'algèbre commutative,
à l'automne 2021 à l'Université Columbia.

\begin{exercise}[Definitions]
\label{exercise-definitions-fall-2021}
Fournir de brèves définitions des concepts italisés.
\begin{enumerate}
\item a {\it sous-ensemble multiplicatif} d'une bague $A$,
\item un anneau artinien {\it } $A$,
\item le spectre {\it d'un anneau} $A$ comme espace topologique,
\item une carte à anneau plat {\it } $A \to B$,
\item la hauteur {\it} d'un $\mathfrak p$ idéal dans $A$, et
\item les functeurs {\it $\text{Tor}^A_i(-, -)$} sur un anneau $A$.
\end{enumerate}
\end{exercise}

\begin{exercise}[Theorems]
\label{exercise-results-fall-2021}
Précisément mais brièvement énoncez un fait non trivial discuté dans les conférences relatives à
chaque élément (s'il y en a plus d'un, il suffit de choisir
l'un d'eux).
\begin{enumerate}
\item Anneaux artiniens,
\item planéité et idéaux premiers,
\item longueurs de $A/\mathfrak m^n$ pour $(A, \mathfrak m)$ noethérien local,
\item la formule de dimension pour les anneaux caténaires universels Noétheriens,
\item complétement d'un anneau local noétherien, et
\item La dualité Matlis pour les anneaux locaux artiniens.
\end{enumerate}
\end{exercise}

\begin{exercise}[Units]
\label{exercise-simple-units}
Quelle est la structure du groupe d'unités de $\mathbf{Z}[x, 1/x]$ en
tant que groupe abelien? Aucune explication nécessaire.
\end{exercise}

\begin{exercise}[Ideals]
\label{exercise-ideals}
Soit $A = \mathbf{F}_2[x, y]/(x^2, xy, y^2)$ et indiquez
$\overline{x}$ et $\overline{y}$ les images de $x$ et $y$ dans $A$. Listez
les idéaux de $A$. Aucune explication n'est nécessaire.
\end{exercise}

\begin{exercise}[Tor and Ext]
\label{exercise-compute-tor-ext}
Soit $(A, \mathfrak m, \kappa)$ être un anneau local
noethérien. Réglez $\varphi(n) = \dim_\kappa \mathfrak m^n/\mathfrak m^{n + 1}$.
\begin{enumerate}
\item Afficher que $\text{Tor}_1^A(A/\mathfrak m^n, \kappa)$ a
une dimension $\varphi(n)$ comme espace-vecteur $\kappa$.
\item Afficher que $\text{Ext}^1_A(A/\mathfrak m^n, \kappa)$ a
une dimension $\varphi(n)$ comme espace-vecteur $\kappa$.
\end{enumerate}
\end{exercise}

\begin{exercise}[Two vectors]
\label{exercise-two-vectors}
Soit $A = \mathbf{Z}[a_1, a_2, a_3, b_1, b_2, b3]$. Définissez
$a = (a_1, a_2, a_3)$ et $b = (b_1, b_2, b_3)$ dans $A^{\oplus
3}$. Considérez l'ensemble
$$
Z = \{\mathfrak p \in \Spec(A) \mid a, b
\text{ map to linearly dependent vectors of }
\kappa(\mathfrak p)^{\oplus 3}\}
$$
\begin{enumerate}
\item Prouver que $Z$ est un sous-ensemble fermé de $\Spec(A)$.
\item Quelle est la dimension $\dim(Z)$ de $Z$?
\item Qu'adviendrait-il de $\dim(Z)$ si nous remplaçions $\mathbf{Z}$ par un corps?
\end{enumerate}
\end{exercise}

\begin{exercise}[Injectives]
\label{exercise-injective-artinian-local}
Laisser $(A, \mathfrak m, \kappa)$ être un anneau local Artinien.
Assumer $A$ est injecteur sous forme de module $A$.
Montrer que $\Hom_A(\kappa, A)$ a une dimension $1$
a un espace $\kappa$-vecteur.
\end{exercise}



\section{ Géométrie algébrique, examen final, printemps 2022}
\label{section-final-exam-spring-2022}

\noindent
Il s'agissait des questions de l'examen final d'un cours sur la géométrie algébrique,
au printemps de 2022 à l'Université Columbia.

\begin{exercise}[Definitions]
\label{exercise-definitions-spring-2022}
Fournir de brèves définitions des concepts italisés.
\begin{enumerate}
\item un schéma {\it},
\item un module {\it quasi cohérent} sur un schéma $X$,
\item un morphisme {\it plat} des schémas $X \to Y$,
\item un {\it fini} morphisme des schémas $X \to Y$,
\item un schéma de groupe {\it} $G$ sur un schéma de base $S$,
\item une famille de variétés {\it } sur un schéma de base $S$,
\item le degré {\it} d'un point fermé $x$ sur une variété $X$
sur le corps $k$,
\item comme d'habitude {\it Hauteur logarithmique} d'un point
$p = (a_0 : \ldots : a_n)$ dans $\mathbf{P}^n(\mathbf{Q})$, et
\item un corps {\it $C_i$}.
\end{enumerate}
\end{exercise}

\begin{exercise}[Theorems]
\label{exercise-results-spring-2022}
Précisément mais brièvement énoncez un fait non trivial discuté dans les conférences relatives à
chaque élément (s'il y en a plus d'un, il suffit de choisir
l'un d'eux).
\begin{enumerate}
\item morphismes d'un schéma $X$ au schéma d'affine $\Spec(A)$,
\item cohomologie d'un module quasi-cohérent $\mathcal{F}$ sur un
schéma d'affine $X$,
\item le groupe Picard de $\mathbf{P}^1_k$ où $k$ est un corps,
\item les dimensions des fibres d'un morphisme plat approprié $X
\to S$ pour $S$ Noéthérien,
\item $\mathbf{G}_m$-modules équivalents sur un schéma $S$, et
\item Théorème de Bezout sur les intersections (se limiter à un
cas spécial si vous le souhaitez).
\end{enumerate}
\end{exercise}

\begin{exercise}[Cubic hypersurfaces]
\label{exercise-cubics}
Laisse-moi. $F \in \mathbf{C}[T_0, \ldots, T_n]$ être homogène de degré $3$. Compte tenu
$3$ vecteurs $x, y, z \in \mathbf{C}^{n + 1}$ prendre en considération la condition
$$
(*)\quad
F(\lambda x + \mu y + \nu z) = 0 \text{ in } \mathbf{C}[\lambda, \mu, \nu]
$$
\begin{enumerate}
\item Quelle est la dimension de l'espace de tout choix de $x, y, z$?
\item Combien d'équations sur les coordonnées de $x$, $y$ et $z$ est
la condition (*)?
\item Quelle est la dimension attendue de l'espace de tous les triples $x,
y, z$ telle que (*) est vraie?
\item Quelle est la dimension de l'espace de tous les triples de telle sorte
que $x, y, z$ soit linéairement dépendante?
\item Conclure que sur une hypersurface de degré $3$ dans $\mathbf{P}^n$ nous nous
attendons à trouver un sous-espace linéaire de dimension $2$ fourni $n \geq
a$ où il est à vous de trouver $a$.
\end{enumerate}
\end{exercise}

\begin{exercise}[Heights]
\label{exercise-heights}
Soit $K$ être un corps. Soit $h_n : \mathbf{P}^n(K) \to \mathbf{R}$, $n \geq
0$ être une collection de fonctions satisfaisant les axiomes $2$ dont nous
avons parlé dans les conférences. Soit $X$ être une
variété projective sur $K$. Soit $\mathcal{L}$ être un module $\mathcal{O}_X$ invertible et
rappelez-vous que nous avons construit dans les conférences une fonction de
hauteur associée $h_\mathcal{L} : X(K) \to \mathbf{R}$. Soit $\alpha :
X \to X$ être un automorphisme de $X$ sur $K$.
\begin{enumerate}
\item Prouver que $P \mapsto h_\mathcal{L}(\alpha(P))$ diffère
de la fonction $h_{\alpha^*\mathcal{L}}$ par une quantité
limitée. (Enfin, rappelez-vous que s'il y a un morphisme
$\varphi : X \to
\mathbf{P}^n$ avec $\mathcal{L} = \varphi^*\mathcal{O}_{\mathbf{P}^n}(1)$, puis par la construction $h_\mathcal{L}(P) =
h_n(\varphi(P))$ et jouez avec cela. En général, écrivez
$\mathcal{L}$ comme une différence de deux de ces deux.)
\item Supposons que $h_\mathcal{L}(P) - h_\mathcal{L}(\alpha(P))$ est non
lié sur $X(K)$. Montrer
que $h_\mathcal{N}$ avec
$\mathcal{N} = \mathcal{L} \otimes \alpha^*\mathcal{L}^{\otimes -1}$ est non
lié sur $X(K)$.
\item Assumer $X$ est une courbe elliptique et
que $\mathcal{L}$ est un module symétrique ample et invertible sur
$X$ de telle sorte que $h_\mathcal{L}$ est non lié
sur $X(K)$. Montrer qu'il existe un module invertible $\mathcal{N}$ de degré $0$
tel que $h_\mathcal{N}$ est non lié. (Enfin,
rappelez-vous que $X$ est une variété abelienne de dimensions $1$. Ainsi
$h_\mathcal{L}$ est quadratique jusqu'à une constante par les résultats dans
les conférences. Choisissez un point approprié $P_0 \in X(K)$. Soit
$\alpha : X \to X$ être traduit par $P_0$.
Considérez $P \mapsto h_\mathcal{L}(P) - h_\mathcal{L}(P + P_0)$. Appliquez les
résultats que vous avez prouvés ci-dessus.)
\end{enumerate}
\end{exercise}

\begin{exercise}[Monomorphisms]
\label{exercise-monomorphisms}
Soit $f : X \to Y$ être un monomorphisme dans la catégorie des
schémas : pour toute paire de morphismes $a, b : T \to
X$ des schémas si $f \circ a = f \circ b$, alors $a
= b$. Montrer que $f$ est injecteur sur des points.
\end{exercise}

\begin{exercise}[Fixed points]
\label{exercise-fix-points}
Laisser $k$ être un corps algébrique fermé.
\begin{enumerate}
\item Si $G = \mathbf{G}_{m, k}$ montre que si $G$ agit sur
une variété projective $X$ sur $k$, alors l'action a un point fixe, c'est-à-dire qu'il
existe un point $x \in X(k)$ tel que $a(g, x) = x$ pour tous
$g \in G(k)$.
\item Idem avec $G = (\mathbf{G}_{m, k})^n$ égal au produit des
copies $n \geq 1$ du groupe multiplicatif.
\item Donner un exemple d'action d'un schéma de groupe connecté $G$
sur une variété projective lisse $X$ qui n'a pas de point fixe.
\end{enumerate}
\end{exercise}



\section{ Géométrie algébrique, examen final, printemps 2025}
\label{section-final-exam-spring-2025}

\noindent
Ce sont les questions de l'examen final d'un cours sur la géométrie algébrique,
au printemps de 2025 à l'Université Columbia.

\begin{exercise}[Definitions]
\label{exercise-definitions-spring-2025}
Fournir de brèves définitions des concepts italisés.
\begin{enumerate}
\item un schéma {\it} $X$,
\item un module {\it quasi cohérent} sur un schéma $X$,
\item le groupe {\it Picard} d'un schéma $X$,
\item donnent un morphisme $f : X \to
Y$ des schémas et un module $\mathcal{O}_X$ $\mathcal{F}$, le {\it pushfore} $f_*\mathcal{F}$,
\item a {\it immersion fermée} des schémas, et
\item une variété {\it} sur un corps donné $k$.
\end{enumerate}
\end{exercise}

\begin{exercise}[Theorems]
\label{exercise-results-spring-2025}
Précisément et succinctement, indiquez un fait non trivial discuté dans les conférences relatives
à chaque élément (s'il y en a plus d'un, il suffit d'en choisir un).
\begin{enumerate}
\item espace topologique d'un schéma $X$,
\item un critère de représentativité pour les functeurs de la catégorie des programmes,
\item le functeur Picard pour une courbe projective lisse $X$ sur un
corps à fermeture algébrique $k$,
\item torsion dans le groupe Picard d'une courbe projective lisse $X$ sur un
corps à fermeture algébrique $k$,
\item un théorème sur le groupe Picard d'un produit $X \times Y \times
Z$ de variétés projectives lisses $X$, $Y$, $Z$ sur un corps à
fermeture algébrique $k$.
\end{enumerate}
\end{exercise}

\begin{exercise}
\label{exercise-affine-line-infinite-nr-points}
Soit $k$ être un corps. Expliquer pourquoi $\Spec(k[x])$ a infiniment beaucoup de points.
\end{exercise}

\begin{exercise}
\label{exercise-hypersurface-variety}
Soit $k$ être un corps et laissez $x_1, \ldots, x_n$ être
des variables. Soit $f \in k[x_1, \ldots, x_n]$ être un polynôme
non constant. Soit $X = \Spec(k[x_1, \ldots, x_n]/(f))$. Quelle propriété
devrait avoir $f$ pour que $X$ soit une variété sur $k$?
\end{exercise}

\begin{exercise}
\label{exercise-find-singular-point}
Trouver un point singulier (vous n'avez pas besoin de les trouver
tous) sur le spectre de $\mathbf{C}[x, y]/(x^5 - 5xy^4 + 20y -
16)$ vu comme une variété sur $\mathbf{C}$.
\end{exercise}

\begin{exercise}
\label{exercise-kernel-restriction-pic-curve}
Soit $X$ être une courbe de projection lisse sur un corps fermé algébriquement
$k$. Soit $x \in X(k)$ être un point fermé sur $X$.
Soit $U = X \setminus \{x\}$. Que pouvez-vous dire sur le noyau
de la carte de restriction $\Pic(X) \to \Pic(U)$?
\end{exercise}

\begin{exercise}
\label{exercise-sections-degree-2g-minus-4}
Soit $X$ être une courbe projective lisse sur un corps algébrique
fermé $k$ du genre $g \geq 100$.
Quelles valeurs peuvent $h^0(\mathcal{L}) = \dim_k H^0(X, \mathcal{L})$ prendre
pour $\mathcal{L}$ être un $\mathcal{O}_X$-module invertible de degré $2g - 4$ sur
$X$? Répondez aussi complètement que possible.
\end{exercise}









\begin{multicols}{2}[\section{Autres chapitres}]
\noindent
Préliminaires
\begin{enumerate}
\item \hyperref[introduction-section-phantom]{Introduction}
\item \hyperref[conventions-section-phantom]{Conventions}
\item \hyperref[sets-section-phantom]{Théorie des ensembles}
\item \hyperref[categories-section-phantom]{Catégories}
\item \hyperref[topology-section-phantom]{Topologie}
\item \hyperref[sheaves-section-phantom]{Faisceaux sur les espaces}
\item \hyperref[sites-section-phantom]{Sites et faisceaux}
\item \hyperref[stacks-section-phantom]{Champs}
\item \hyperref[fields-section-phantom]{Corps}
\item \hyperref[algebra-section-phantom]{Algèbre commutative}
\item \hyperref[brauer-section-phantom]{Groupes de Brauer}
\item \hyperref[homology-section-phantom]{Algèbre homologique}
\item \hyperref[derived-section-phantom]{Catégories dérivées}
\item \hyperref[simplicial-section-phantom]{Méthodes simpliciales}
\item \hyperref[more-algebra-section-phantom]{Compléments d'algèbre}
\item \hyperref[smoothing-section-phantom]{Lissage des morphismes d'anneaux}
\item \hyperref[modules-section-phantom]{Faisceaux de modules}
\item \hyperref[sites-modules-section-phantom]{Modules sur les sites}
\item \hyperref[injectives-section-phantom]{Objets injectifs}
\item \hyperref[cohomology-section-phantom]{Cohomologie des faisceaux}
\item \hyperref[sites-cohomology-section-phantom]{Cohomologie sur les sites}
\item \hyperref[dga-section-phantom]{Algèbre différentielle graduée}
\item \hyperref[dpa-section-phantom]{Algèbre à puissances divisées}
\item \hyperref[sdga-section-phantom]{Faisceaux différentiels gradués}
\item \hyperref[hypercovering-section-phantom]{Hyperrecouvrements}
\end{enumerate}
Schémas
\begin{enumerate}
\setcounter{enumi}{25}
\item \hyperref[schemes-section-phantom]{Schémas}
\item \hyperref[constructions-section-phantom]{Constructions de schémas}
\item \hyperref[properties-section-phantom]{Propriétés des schémas}
\item \hyperref[morphisms-section-phantom]{Morphismes de schémas}
\item \hyperref[coherent-section-phantom]{Cohomologie des schémas}
\item \hyperref[divisors-section-phantom]{Diviseurs}
\item \hyperref[limits-section-phantom]{Limites de schémas}
\item \hyperref[varieties-section-phantom]{Variétés}
\item \hyperref[topologies-section-phantom]{Topologies sur les schémas}
\item \hyperref[descent-section-phantom]{Descente}
\item \hyperref[perfect-section-phantom]{Catégories dérivées de schémas}
\item \hyperref[more-morphisms-section-phantom]{Compléments sur les morphismes}
\item \hyperref[flat-section-phantom]{Compléments sur la platitude}
\item \hyperref[groupoids-section-phantom]{Schémas en groupoïdes}
\item \hyperref[more-groupoids-section-phantom]{Compléments sur les schémas en groupoïdes}
\item \hyperref[etale-section-phantom]{Morphismes étales de schémas}
\end{enumerate}
Thèmes de théorie des schémas
\begin{enumerate}
\setcounter{enumi}{41}
\item \hyperref[chow-section-phantom]{Homologie de Chow}
\item \hyperref[intersection-section-phantom]{Théorie de l'intersection}
\item \hyperref[pic-section-phantom]{Schémas de Picard des courbes}
\item \hyperref[weil-section-phantom]{Théories cohomologiques de Weil}
\item \hyperref[adequate-section-phantom]{Modules adéquats}
\item \hyperref[dualizing-section-phantom]{Complexes dualisants}
\item \hyperref[duality-section-phantom]{Dualité pour les schémas}
\item \hyperref[discriminant-section-phantom]{Discriminants et différentes}
\item \hyperref[derham-section-phantom]{Cohomologie de de Rham}
\item \hyperref[local-cohomology-section-phantom]{Cohomologie locale}
\item \hyperref[algebraization-section-phantom]{Géométrie algébrique et formelle}
\item \hyperref[curves-section-phantom]{Courbes algébriques}
\item \hyperref[resolve-section-phantom]{Résolution des surfaces}
\item \hyperref[models-section-phantom]{Réduction semi-stable}
\item \hyperref[functors-section-phantom]{Foncteurs et morphismes}
\item \hyperref[equiv-section-phantom]{Catégories dérivées de variétés}
\item \hyperref[pione-section-phantom]{Groupes fondamentaux de schémas}
\item \hyperref[etale-cohomology-section-phantom]{Cohomologie étale}
\item \hyperref[crystalline-section-phantom]{Cohomologie cristalline}
\item \hyperref[proetale-section-phantom]{Cohomologie pro-étale}
\item \hyperref[relative-cycles-section-phantom]{Cycles relatifs}
\item \hyperref[more-etale-section-phantom]{Compléments de cohomologie étale}
\item \hyperref[trace-section-phantom]{La formule des traces}
\end{enumerate}
Espaces algébriques
\begin{enumerate}
\setcounter{enumi}{64}
\item \hyperref[spaces-section-phantom]{Espaces algébriques}
\item \hyperref[spaces-properties-section-phantom]{Propriétés des espaces algébriques}
\item \hyperref[spaces-morphisms-section-phantom]{Morphismes d'espaces algébriques}
\item \hyperref[decent-spaces-section-phantom]{Espaces algébriques décents}
\item \hyperref[spaces-cohomology-section-phantom]{Cohomologie des espaces algébriques}
\item \hyperref[spaces-limits-section-phantom]{Limites d'espaces algébriques}
\item \hyperref[spaces-divisors-section-phantom]{Diviseurs sur les espaces algébriques}
\item \hyperref[spaces-over-fields-section-phantom]{Espaces algébriques sur des corps}
\item \hyperref[spaces-topologies-section-phantom]{Topologies sur les espaces algébriques}
\item \hyperref[spaces-descent-section-phantom]{Descente et espaces algébriques}
\item \hyperref[spaces-perfect-section-phantom]{Catégories dérivées d'espaces}
\item \hyperref[spaces-more-morphisms-section-phantom]{Compléments sur les morphismes d'espaces}
\item \hyperref[spaces-flat-section-phantom]{Platitude sur les espaces algébriques}
\item \hyperref[spaces-groupoids-section-phantom]{Groupoïdes dans les espaces algébriques}
\item \hyperref[spaces-more-groupoids-section-phantom]{Compléments sur les groupoïdes dans les espaces}
\item \hyperref[bootstrap-section-phantom]{Amorçage}
\item \hyperref[spaces-pushouts-section-phantom]{Sommes amalgamées d'espaces algébriques}
\end{enumerate}
Thèmes de géométrie
\begin{enumerate}
\setcounter{enumi}{81}
\item \hyperref[spaces-chow-section-phantom]{Groupes de Chow des espaces}
\item \hyperref[groupoids-quotients-section-phantom]{Quotients de groupoïdes}
\item \hyperref[spaces-more-cohomology-section-phantom]{Compléments sur la cohomologie des espaces}
\item \hyperref[spaces-simplicial-section-phantom]{Espaces simpliciaux}
\item \hyperref[spaces-duality-section-phantom]{Dualité pour les espaces}
\item \hyperref[formal-spaces-section-phantom]{Espaces algébriques formels}
\item \hyperref[restricted-section-phantom]{Algébrisation des espaces formels}
\item \hyperref[spaces-resolve-section-phantom]{Résolution des surfaces revisitée}
\end{enumerate}
Théorie des déformations
\begin{enumerate}
\setcounter{enumi}{89}
\item \hyperref[formal-defos-section-phantom]{Théorie formelle des déformations}
\item \hyperref[defos-section-phantom]{Théorie des déformations}
\item \hyperref[cotangent-section-phantom]{Le complexe cotangent}
\item \hyperref[examples-defos-section-phantom]{Problèmes de déformation}
\end{enumerate}
Champs algébriques
\begin{enumerate}
\setcounter{enumi}{93}
\item \hyperref[algebraic-section-phantom]{Champs algébriques}
\item \hyperref[examples-stacks-section-phantom]{Exemples de champs}
\item \hyperref[stacks-sheaves-section-phantom]{Faisceaux sur les champs algébriques}
\item \hyperref[criteria-section-phantom]{Critères de représentabilité}
\item \hyperref[artin-section-phantom]{Axiomes d'Artin}
\item \hyperref[quot-section-phantom]{Espaces de Quot et de Hilbert}
\item \hyperref[stacks-properties-section-phantom]{Propriétés des champs algébriques}
\item \hyperref[stacks-morphisms-section-phantom]{Morphismes de champs algébriques}
\item \hyperref[stacks-limits-section-phantom]{Limites de champs algébriques}
\item \hyperref[stacks-cohomology-section-phantom]{Cohomologie des champs algébriques}
\item \hyperref[stacks-perfect-section-phantom]{Catégories dérivées de champs}
\item \hyperref[stacks-introduction-section-phantom]{Introduction aux champs algébriques}
\item \hyperref[stacks-more-morphisms-section-phantom]{Compléments sur les morphismes de champs}
\item \hyperref[stacks-geometry-section-phantom]{La géométrie des champs}
\end{enumerate}
Thèmes de théorie des espaces de modules
\begin{enumerate}
\setcounter{enumi}{107}
\item \hyperref[moduli-section-phantom]{Champs de modules}
\item \hyperref[moduli-curves-section-phantom]{Espaces de modules de courbes}
\end{enumerate}
Divers
\begin{enumerate}
\setcounter{enumi}{109}
\item \hyperref[examples-section-phantom]{Exemples}
\item \hyperref[exercises-section-phantom]{Exercices}
\item \hyperref[guide-section-phantom]{Guide bibliographique}
\item \hyperref[desirables-section-phantom]{Desiderata}
\item \hyperref[coding-section-phantom]{Style de programmation}
\item \hyperref[obsolete-section-phantom]{Obsolète}
\item \hyperref[fdl-section-phantom]{Licence de documentation libre GNU}
\item \hyperref[index-section-phantom]{Index généré automatiquement}
\end{enumerate}
\end{multicols}


\bibliography{my}
\bibliographystyle{amsalpha}

\end{document}

